% =====================================================================
%  R. Nielsen, OM AANDSDANNELSE med særligt Hensyn paa vort Opdragelsesvæsen.
%  Sex Forelæsninger.  Kjøbenhavn: Gyldendal, 1877.  109 pp.
%
%  TRANSCRIPTION -- GENERATED; DO NOT EDIT BY HAND.  Made in three steps, all rerunnable:
%    python3 ebook-method/om-aandsdannelse/draft.py --emit --force
%      this template + the scan's text layer (the Royal Library's OCR, in KB's scan;
%      ~/bibliotek/Nielsen, Rasmus/aandsdannelse.pdf, see SCANS.tsv), then
%      patches.PATCHES (corrections read at the image that a word decision cannot express);
%    python3 ebook-method/collate/check.py om-aandsdannelse --italics --apply
%      the decisions in ebook-method/om-aandsdannelse/decisions.tsv, taken against an
%      independent tesseract reading of the images (rules.py writes the ones settled by
%      rule, each marked with its rule; the rest were read at the image), and the
%      italics decisions (--italics must be in the same call as --apply);
%    python3 ebook-method/om-aandsdannelse/draft.py --post
%      patches.POST: corrections written against the collated text.
%  Run them chained (&&): each step exits without writing if a patch fails to match.
%  WORD-ACCURACY (2026-10-06).  No e-book and no earlier transcription exist.  The draft
%  (KB's OCR layer, lines rebuilt from the word boxes) was collated word by word, with
%  punctuation, quote kinds and word division, against tesseract's independent reading of
%  the same images: 722 disagreements.  206 were settled by rule: 82 by check.py (glyph
%  faults, noise, moved words), 124 by rules.py (tesseract's „ read as a comma and its
%  guillemets for the print's „ “ -- the print sets „...“ only; accents tesseract adds; the
%  layer's own glyph faults where tesseract's word is attested in the corpus and the
%  layer's is not).  The rest were read at the image: 518 crops read BLIND (the reader saw
%  only the printed lines; .parts/reads/), compared by script (reads.py) or decided by
%  hand from the readers' transcriptions, and 32 more read by the main session.  OPEN 0
%  (words and italics); no decision left unapplied that changes the text.  Then, at the
%  image: letterspacing (114 candidate lines, 38 spaced; spaced runs garbled by both OCRs
%  set by patch), the title page, Indhold and the six lecture openings; the --splits and
%  --moved lists; the words attested nowhere else in the corpus near an attested word (14
%  read: 8 misprints marked PRINTED AS IS, 5 OCR faults mended).  Spot-read in full, pp. 47
%  and 88, with two planted control errors (2 of 2 caught): 2 letterspaced words had been
%  missed (now set).  109 of 109 page markers exact (check.py --pages).
%  TRANSLATION READING (2026-10-06).  Translating the book (translation.tex) was a full
%  second reading; its 55 odd spots were read at the image (spotsheets.py, sheets t01-t06,
%  ambiguous dots at 600 dpi in z01; .parts/reads/spots-t.tsv): 29 faults mended (28 by
%  patches.POST, 1 by a decision): specks read as . , ' ‘ or a breve accent; dots of risen spacing read as
%  stops or commas; two split words; a dash and a comma not in the print; four false
%  paragraph breaks, pp. 37, 54 (two), 67; p. 67 „tale at ~ om,“ for „tale om, at“; two
%  printed quotation marks missed, pp. 64, 77; a bold „først“, p. 66; 3 letterspaced runs
%  set (pp. 7, 31, 60); 8 more misprints marked PRINTED AS IS (pp. 26, 29, 59, 68 three,
%  84, 86) and one unclosed quotation noted (p. 35); 12 confirmed as printed.
%  NOT ESTABLISHED: letterspacing of short words (under five letters) outside the lines
%  read -- the detector misses them, so some \emph spans may be missing; misprints that
%  both OCRs read alike as real words (no full image pass).  The book has no footnotes and
%  no italics; one word is bold (p. 66, \textbf).  This copy is heavily marked in pen and pencil; such marks are not
%  transcribed.
%
%  CONVENTIONS.  "% --- p. N ---" + \opage{N} at the first word of printed page N
%  (a word broken over the turn is split with %).  Footnotes as printed, numbered
%  per page.  Italics \textit; letterspacing \emph.  Folios, signature marks and the
%  library's stamps are not transcribed.  Print governs; misprints are kept and
%  marked PRINTED AS IS.
% =====================================================================
\documentclass[12pt,a4paper,openany]{book}
\usepackage[utf8]{inputenc}
\usepackage[T1]{fontenc}
\usepackage{libertinus}
\usepackage{libertinust1math}
\usepackage[danish]{babel}
\usepackage[protrusion=true,expansion=true]{microtype}
\usepackage{setspace}
\usepackage[margin=1.3in]{geometry}
\usepackage{fancyhdr}
\usepackage{hyperref}
\hypersetup{pdftitle={Om Aandsdannelse}, pdfauthor={Rasmus Nielsen}, hidelinks}
\pagestyle{fancy}
\fancyhf{}
\fancyhead[LE,RO]{\thepage}
\fancyhead[C]{\textit{R. Nielsen: Om Aandsdannelse (1877)}}
\setstretch{1.3}
\usepackage{marginnote}
\renewcommand{\thefootnote}{\arabic{footnote})}
\newcommand{\tocline}[3]{\noindent\makebox[2.2em][r]{#1}\hspace{0.6em}#2\dotfill\ #3\par}
\newcommand{\opage}[1]{\marginnote{\footnotesize\textit{[#1]}}}
% a lecture opens a new page: heading, a short rule; the print's large initial is not reproduced
\newcommand{\lecture}[1]{\clearpage\begin{center}{\large #1}\\[0.4em]\rule{3em}{0.4pt}\end{center}\medskip\noindent}

\begin{document}

% --- front matter (PDF 6, title; PDF 8, Indhold), read at the image 2026-10-06 (sheet m01).
% p. [I], title.  The publisher's device (a round monogram seal) is not reproduced; KB's stamp is not transcribed.
\begin{titlepage}
\begin{center}
\vspace*{3em}
{\Huge Om Aandsdannelse}\\[1em]
{\large med særligt Hensyn paa vort Opdragelsesvæsen.}\\[2.5em]
\emph{Sex Forelæsninger}\\[1em]
af\\[1em]
{\large \emph{R. Nielsen.}}\\[5em]
% [publisher's device]
\emph{Kjøbenhavn.}\\[0.5em]
Gyldendalske Boghandels Forlag (F. Hegel \& Søn).\\[0.5em]
{\small Trykt hos J. Jørgensen \& Co.}\\[0.5em]
{\small 1877.}
\end{center}
\end{titlepage}
% p. [III], Indhold.  „Forelæsning“ and the ordinals are letterspaced in every entry.
\begin{center}{\large\textbf{Indhold.}}\\[0.3em]\rule{3em}{0.4pt}\end{center}
\hfill{\small Side.}\par
{\small
\tocline{}{\emph{Første Forelæsning}: Forskjel mellem Kultur- og Aandsdannelse.}{1}
\tocline{}{\emph{Anden Forelæsning}: Autoritetstro og personlig Selvforstaaelse}{17}
\tocline{}{\emph{Tredie Forelæsning}: Aand og Bogstav: det kritiske Vendepunkt}{38}% PRINTED AS IS: the lecture begins on p. 34 (a reader's pencil corrects it)
\tocline{}{\emph{Fjerde Forelæsning}: S. Kierkegaards Optræden}{52}
\tocline{}{\emph{Femte Forelæsning}: S. Kierkegaards Selvmodsigelse}{71}
\tocline{}{\emph{Sjette Forelæsning}: Religion og Humanitet}{91}
}


% ===== TEXT (pp. 1--109) =====
% --- p. 1 ---
\setcounter{footnote}{0}%
\lecture{Første Forelæsning.}\opage{1}
Der gjøres i vore Dage store og betydelige Anstrengelser
til Bedste for Oplysning og Dannelse. Ved Siden
af den lærde Dannelse og de mangfoldige Arter af Fagdannelse
lægges der Vægt paa almindelig Dannelse:
politisk, æsthetisk, selskabelig Dannelse. Enhver i Samfundet
sætter en Ære i, at gjælde for et dannet Menneske.
Det gaaer med Dannelse som med Paaklædning. Dannelse
er en klædelig Pynt; hvo ønsker ikke at være
velklædt og tage sig ud til sin Fordeel?

Men trods den paaskjønnelsesværdige Iver for Dannelsens
Alsidighed turde det dog, naar vi see nøiere
efter, vise sig, at der er Noget glemt, at den lovpriste
Alsidighed er i høi Grad eensidig. Hvad Nutidens Oplyste
forstaae ved Dannelse er saa godt som udelukkende
en Dannelse i Forhold til \emph{det Mangfoldige}, er hvad
vi med eet Ord ville kalde \emph{Kulturdannelse}; at der
ogsaa maa gives en Dannelse i Forhold til det Enfoldige,
og at det er den egenlige \emph{Aandsdannelse},
bliver, saa vidt vi skjønne, alt for letfærdig overseet.
Det forudsættes uden videre, at Kulturdannelse er Et
og Alt, saa at, naar man blot har Kultur i rigeligt
Maal, følger Aandsdannelsen med af sig selv; det er
en betænkelig Vildfarelse.

% --- p. 2 ---
\setcounter{footnote}{0}%
\opage{2}Kulturdannelsen vender Blikket udad; hvad som helst
der i Mangfoldighedens Verden glider det speidende Øie
forbi, forvandler sig til Phænomen. Tanken, det Usandselige,
er Phænomen lige saa vel som det Sandselige, Samvittigheden
Phænomen lige saa vel som Elektriciteten, Geniet
et Phænomen saa vel som Lynet, Kristus et Phænomen som
Buddha og Muhamed. Kalde vi nu det Menneske, hvis hele
Bevidsthedsindhold tømmer sig ud i Kulturdannelsens
Former, et Kulturmenneske, da er det let at forstaae, at
et saadant Menneske, hvor travlt han end har med Tingene
og med sig selv, dog aldrig kan blive saa indadvendt,
at han i Alvor faaer Øie enten paa Tingenes eller sit
eget Væsen. Han negter ingenlunde, at der er Væsen
i Tingene, men, siger han, dette Væsen i og for sig er
ikke til at erkjende. Han negter ikke, at han jo paa
en Maade har en Sjæl, men hvad Sjælen er for et
Væsen og om den virkelig er et Væsen, hvordan skulde
han kunne vide det. Vil han betragte sin Sjæl for
derved at kunne opdage sit eget Selv, da bliver dette
Selv ham øieblikkelig et Phænomen. Og hvad ligger
saa til Grund for dette Phænomen? Noget Mystisk,
hvorover man ikke skal gruble. Ligegyldig for hvad
han i sig selv er, danner han sig da til at synes Noget,
at blive synet og tage sig ud som et anseeligt Phænomen
i Phænomenernes Række.

Udad vendt mod det Mangfoldige, er Kulturbevidstheden
alsidig; den er som en Argus besaaet med Øine,
hundrede klare Øine; men disse hundrede klare Øine
ere alle den Trolddom underlagte, at hvad de fæste
sig paa, øieblikkelig skjuler sit Væsen og forvandler sig
til blot og bart Phænomen. Men i et Phænomen, der
holder sit Væsen skjult, kan man ikke fordybe sig, og
% --- p. 3 ---
\setcounter{footnote}{0}%
\opage{3}det, hvori man ikke kan fordybe sig, har kun flygtig,
forbigaaende Interesse. Man bliver kjed af bekjendte
Phænomener, som Barnet af sit Legetøi, og higer efter
Nyt. Intet enkelt Phænomen har Gyldighed i og for
sig; kun Phænomenernes Afvexling, Farvespil og Sammenhæng
i store Grupper har Interesse. Er man færdig
med dem i Hjemlandet, saa maa man afsted og see sig
om i Udlandet. Kulturmennesket er Kosmopolit, en
Odysseus, vidt bereist; han taler mange Sprog, har
været i mange Lande, har levet baade med Kykloper
og Nympher. Hvor er det oplivende at høre ham tale
om de Bjerge, han har besteget, de Musæer, han har
besøgt, de Udgravninger af uvurdeerlige Oldtidsskatte, han
selv har overværet. Han kan kritisere Forfattere, thi
han kjender Literaturerne; han kan beskrive Bygninger
og Monumenter, thi han kjender Stiilarterne; Malerier
og Statuer! ja, han har Kjenderblik, han har Smag.

Men hvad veed Kulturmennesket i Forhold til det
Enfoldige? Hvilket Spørgsmaal! Er Enfoldighed et
Phænomen lige som alt Andet i Verden, skulde den
dannede Mand saa ikke vide, at der foruden Dumhed
gives høiere Arter af Enfoldighed, i Barnets Uskyld, i
Kvindens jomfruelige Sjæl, i Kunstnerens geniale Naivetet,
i Almuetroen, i Religionen? Ja vel, siger han, er den
hellige Enfoldighed ædel og skjøn; men Kritiken er
Dannelsens Blomst, Kulturens Fuldendelse, Kritiken „den
tiende Muse“, forholder sig til Enfoldigheden som Falken
til Duen; et giftigt Stik af spids Kritik kan dræbe
Troen, og med Troen falder Religionen.

Aandsdannelsen vender Blikket indad men dog
ikke saaledes, at den udelukker det Mangfoldige. Man
har i Optiken Lindser af forskjellig Form,  nogle, der
% --- p. 4 ---
\setcounter{footnote}{0}%
\opage{4}sprede, andre, der samle Straalerne. Kulturdannelsen
foretrækker straalespredende Lindser med et indbildt
Brændpunkt; Aandsdannelsen kræver Straalesamlere og
stiller dem saaledes, at Brændpunktet falder i Selvet.
„Kjend dig selv!“ det er theoretisk: fæst Opmærksomheden
paa dit indre, usandselige Væsen, thi kun i og
ved dette Væsen formaaer du at erkjende Phænomenerne.
Det er dog dig selv, der seer og hører, fornemmer og
føler; det er dig selv, der glædes og sørger, tænker og
grunder; det er dig selv, der lever. Vil du vide, hvad
Livet er, hvad Sjælen er, hvad Bevidstheden er, da søg
fremfor Alt Besvarelsen hos dig selv. Himlens Stjerner,
Jordens Dyr og Hjernens Atomer kunne ikke yde dig
ringeste Bidrag paa første Haand, thi hvad du veed
om disse Ting, veed du kun i Kraft af din egen Bevidsthed,
hvad du dømmer om disse Ting, dømmer du
i Kraft af din egen Forstand, hvad Værd du tillægger
disse Ting, tillægger du dem i Kraft af din egen Følelse.
Bild dig ikke ind, at du ved at gaae ud fra en kritisk
Betragtning af Dyreslægterne i deres Mangfoldighed
eller af Menneskeslægten i Almindelighed kan opdage,
hvad Livet er i sit Væsen; den kritiske Betragtning
maa have en Maalestok, men hvad du end drøfter og
hvad du end maaler, saa bærer du i sidste Instans
Maalestokken i dig selv. „Kjend dig selv!“ det er
praktisk: hold fast paa dig selv, ikke i selvisk Indbildning,
men i fri, bevidst Vexelvirkning med den omgivende
Verden. Der er en Tingenes Orden, under
hvilken du maa bøie dig, og hvoraf du maa lade dig
bestemme. Den foreskriver dig Betingelser, men du
kan vælge, den tilbyder dig Nydelser, den paalægger
dig Lidelser, den truer dig med Døden; men hvorledes
% --- p. 5 ---
\setcounter{footnote}{0}%
\opage{5}du forstaaer alt dette, beroer paa din Sindssamling
og din Selvforstaaelse; hvorledes du tager dig i alt
dette, afhænger af din Villiesretning, altsaa væsenlig
af, hvorledes du selv er. --- Fra Adspredelsen i det
Udvortes, fra Phænomenernes Mangfoldighed vil Aandsdannelsen
føre Mennesket tilbage til det enfoldige
Ene, det væsenlige Indre, hvori Selvet lever, og saa
igjennem dette Indre til Allivets Kilde, Naturens Ophav,
til Nødvendighedens Gaade, til Frihedens Udspring,
til Aanden som Aand, til Gud.

Uden Aandsdannelse er ingen sammenhængende,
heelstøbt Verdensanskuelse mulig; uden en sammenhængende
Verdensanskuelse er ingen sand Karakterens
Selvstændighed mulig. En indre fri og selvstændig
Karakteer kan lige saa lidt erhverves ved at
seile med alle Strømme og staae paa Udkig efter alle
Mærkværdigheder som ved at lukke Øiet til for det
Virkelige og lig en Navlebeskuer stirre paa det mystiske
Punkt i Sjælen. Hvad vi i væsenlig Forstand kalde
Verdensanskuelse, Livsanskuelse, beroer paa en Vexelbestemmelse
af det, der kommer uden fra, med det,
der oprindelig er i os selv og kommer indenfra, en
Vexelbestemmelse, der er saa gjennemgaaende, at Sjælen
under alle Indtryk fra Tingenes og Phænomenernes
Verden bevarer Indtrykket af sig selv. Selvet maa vise
sig som medbestemmende, grundbestemmende Factor,
ikke blot som Samler og Bærer af det Hele. Væsenet
i denne alsidige, i Selvbestemmelse rodfæstede Vexelbestemmelse
er Aand.

Aandsdannelse opnaaes ikke derved, at man paa
udvortes Autoritet antager faste Dogmer af en positiv
Religionslære. Der er Mange, som have fuldt op af
% --- p. 6 ---
\setcounter{footnote}{0}%
\opage{6}positive Dogmer om Aanden, og dog have de ingen
Aandsdannelse; hvorimod en „Fritænker“, der negter
alle Dogmer, kan have Aandsdannelse. Hvad det Negative
og Positive i Aandsdannelsen angaaer, kommer
det væsenlig an paa, hvorledes Fornegtelsen og Antagelsen
bliver begrundet. Vil Nogen f. Ex. negte Guds
Tilværelse og Sjælens Udødelighed, fordi han bilder sig
ind, at der i Naturlæren og Kritiken er ført tilstrækkelige,
uimodsigelige Beviser for Antagelsens Urimelighed,
da er han blind og dum i Forhold til Aand. Befinder
han sig derimod i Tilstand af, at kunne forstaae sig
selv og gjennemføre sit Liv conseqvent i Fornegtelsen;
finder han, at der i hans Væsen er Noget, der gjør
Oprør saa vel imod Gud som imod det evige Liv;
negter han, fordi han er sig klart bevidst, at han, ifølge
sin Sindsbeskaffenhed, vil negte og maa negte: da har
hans Fornegtelse Oprindelighedens Præg; han har Aand
i negativ Form, han har Aandsdannelse. En Aandsfornegter
med Aand kan være et Aandens Redskab, en
dygtig Tugtemester for en aandssvækket, aandelig uklar,
aandelig doven, aandelig forfjollet Slægt.

Den positive Aandsdannelse kan gaae i forskjellig
Retning, alt eftersom man stiller sig til de to ideelle
Magter: \emph{Nødvendigheden og Friheden}. Sættes
Nødvendigheden som det Oprindelige, Friheden som det
Afledede, vil Dannelsen gaae i rationel Retning, og Overbeviisningen
grundes paa Viden. Sættes Friheden som
det Oprindelige, Nødvendigheden som det Afledede, vil
Dannelsen gaae i religiøs Retning, og Overbeviisningen
bygges paa Tro. I første Tilfælde falder Eftertrykket
paa Erkjendelsen af den evige Fornuft, som udtaler sig
i Naturen; Mennesket fatter og forstaaer sig selv som
% --- p. 7 ---
\setcounter{footnote}{0}%
\opage{7}et Fornuftvæsen i en Verden af fornuftige, nødvendige
Love. I sidste Tilfælde falder Eftertrykket paa Erkjendelsen
af Gud og det Livets Ord, der aabenbarer hans
Villie. Mennesket fatter og forstaaer sig selv som et
„Guds Billede“, Villiesorgan for den guddommelige Aand.

I hvilken af disse to Hovedretninger Aandsdannelsen
end gaaer, maa Verdenserkjendelsen dog lade sig forlige
med Gudserkjendelsen, og Livsanskuelsen i begge
Henseender være sammenhængende. Fra dette Synspunkt
vil det bedst kunne sees, hvorledes det forholder
sig med det nittende Aarhundredes alsidige Dannelse,
ved nemlig at kaste et Blik paa vort Opdragelsesvæsen.

Man er i Fornuftlæren ikke kommet videre end til
en kritisk Opløsning af den positive Religions Dogmer,
og i Religionslæren ikke videre end til en simpel Gjentagelse
og modificeret Indskærpelse af de gamle Dogmer.
Skulde Opdragelsen bygges paa en Livsanskuelse, hvori
der virkelig var Enhed og Sammenhæng, da maatte
man, i det Mindste indtil videre, have to særskilte Opdragelsesanstalter,
to Slags Skoler, en \emph{Kulturskole}
paa Fornuftens og en \emph{Kirkeskole} paa den positive
Religions Grundlag, saa at der virkelig blev to Livsanskuelser
at vælge imellem.

I Kirkeskolen maatte da alt det strengt holdes
borte, som middelbart eller umiddelbart kunde svække
Troen paa den religiøse Overleverings historiske Ufeilbarhed
og vække Tvivl om Dogmernes absolute Sandhed.
Hele Opdragelsen maatte være indrettet derefter. Kirkeskolen
maatte da ikke indskrænke sig til en Underviisning
i Religion alene, den maatte paa sin Viis give en
virkelig Natur- og Verdenserkjendelse, normerende for
Kultur og Videnskab. Eleverne i en saadan Skole kunne
% --- p. 8 ---
\setcounter{footnote}{0}%
\opage{8}da nok lære Sprog, baade ældre og nyere; men til Oprindelsen
af de mange Sprog maae de ikke kjende
anden Forklaring end den babyloniske Taarnbygning.
At der maa læres Historie, baade Bibelhistorie og almindelig
Verdenshistorie, forstaaer sig; men det Afgjørende
er, at Historiens Gang og de store Begivenheder
stadig holdes under Belysning af Paradiset og
Syndefaldet og Djævelens fortsatte Kamp mod det styrende
Forsyn. Der er da heller Intet i Veien for, at
der ogsaa kan undervises i Mathematik, Physik og Naturhistorie;
men det maa gjælde som urokkelig Forudsætning,
at Verden er skabt i sex Dage, at Himlen er en
fast Hvælving, og at Jorden staaer stille, medens Solen
gaaer op og ned. Hypotheser om Dyreslægter, der skulde
være uddøde før Menneskets Fremkomst, om en naturlig
Raahedstilstand uden et forudgaaende Paradis o. desl.,
maae behandles som „Verdens Børnelærdomme“, Phantasier
af en udøbt Fornuft; thi det er da klart, at intet
Dyr kunde døe, før Døden var kommen ind i Verden,
men Døden kom jo først, efter at Mennesket havde
syndet. Paa denne Maade kommer der Eenhed i Underviisningen.
Mennesket erkjender, hvorledes det er skabt,
hvorledes det er faldet, og hvorledes det igjen kan opreises:
paa dette Grundlag kan en Opdragelse gjennemføres,
og Aandsdannelsen blive derefter.

I Kulturskolen maa Grunden ryddes og Localet
renses for alskens mugne Overleveringer. Kritiken med
„den frie Tanke“ bør sørge for, at der er behørig luftet
ud, saa at man kan aande frit i Fornuftens Atmosphære.
Istedenfor Theologiens Dogmer opstilles Naturlærens
store Grundsætninger: Materiens Evighed, Naturlovenes
Uforanderlighed. Astronomiens Lære om Kloderne i
% --- p. 9 ---
\setcounter{footnote}{0}%
\opage{9}Rummet, Geologiens Fremstilling af Jordperioderne,
Zoologiens Opfattelse af Menneskeslægtens Afstamning,
den almene Forudsætning om Naturen som den absolut
nødvendige, selvtilstrækkelige Grund til Alt, ogsaa til
Aanden, disse Fornufterkjendelser afgive Elementerne
i det Underlag, hvorpaa Livsanskuelsen skal hvile, og
Opdragelsen lægges an. En Sædelære er nødvendig;
men de sædelige Bud maae opfattes som almindelige
Fornuftbud, hvis Opfyldelse kræves af Samfundet, men
hvis Anvendelse indenfor Tvangslovenes Grændse faaer
Frihed og Friskhed, Farve og Liv fra den selvbevidste
Individualitets eiendommelige Naturgrund. Der kan
undervises i Religion lige saa vel som i Kulturhistorie
og Sprog, men den hele Udvikling falder ind under
Fornuftnødvendighedens Synspunkt; det er en Selvfølge,
at de positive Religioner med deres Mirakler og overnaturlige
Phænomener uden Undtagelse opfattes som
Mythologier. Føies hertil en Paaviisning af det Ideale
i Kunst og Poesi; en Anerkjendelse af det menneskelige
Hjertes ædlere Følelser, af Samvittighedens Ret og
den frie Personlighed: saa har man jo en sammenhængende
Verdensanskuelse, paa hvilken et Opdragelsessystem
kan gjennemføres conseqvent. Aandsdannelsen
bliver derefter.

Man lægge vel Mærke til, at de to Opdragelsessystemer
ere saa himmelvidt forskjellige i Aand og
Retning, at enhver Grundsætning, der laanes fra det
ene og optages i det andet, nødvendig maa forvirre
Sindet og svække Karakteren. Har man nu hidtil holdt
de to Skoler adskilte fra hinanden? Nei. Er det sandsynligt,
at man i Fremtiden vil kunne holde dem adskilte?
Nei. Hvad har man da gjort og hvad vil man
% --- p. 10 ---
\setcounter{footnote}{0}%
\opage{10}fremdeles gjøre? Blande dem sammen i al Uklarhed.
Lige fra Almueskolen op til Universitetet fortsættes jo
Blandingen i vor kritiske Nutid paa høist ukritisk
Maade. Man tænker som saa: et dannet Menneske
maa kjende den fornuftige Verdensorden, og et dannet
Menneske bør kjende Religionens Hovedlærdomme. Fornufterkjendelsen
veileder os i vore Foretagender; Religionen
styrker vor moralske Følelse og trøster os i
vore Lidelser. Samfundet kan ikke gjøre Mere for den
unge Slægt end aabne den Adgang til alsidig Dannelse;
Forholdet mellem Fornuften og den positive Religion
er en Gaade, som Enhver maa løse for sig selv, saa
godt han kan.

Men er det nu ogsaa forsvarligt at fylde de unge
Sjæle med et Sammensurium af modstridende Forestillinger
og grunde deres Opdragelse paa en Sum af
mørke Gaader, som man selv hverken kan eller gider
løse? Lad Religionen have sine Hemmeligheder og
Naturen sine, saa er det dog et uafviseligt Krav, at de,
der skulle være Ungdommens Lærere, maae kunne gjøre
Rede for, hvorledes det forholder sig med vor Erkjendelse
af det Overnaturlige i Religionen og vor fornuftige
Verdenserkjendelse. Kunne de ikke det, da maa
den humane Opdragelse i Nutiden, hvad Karakteerdannelsen
angaaer, komme til at staae langt under den
katholske i Middelalderen og den græske og romerske
i Oldtiden.

Man tænke sig hos os en Faders Tiltale til sin
Søn, naar denne efter endt Underviisning forlader Skolen.
„Min Søn! Du kjender nu de to Ledestjerner, som
skulle lyse paa din Vei, Religionen og Fornuften. Tab
dem aldrig af Sigte! De lyse sjeldent lige klart paa
% --- p. 11 ---
\setcounter{footnote}{0}%
\opage{11}een Gang, men naar den ene slukkes, saa tændes den
anden, saa Du vil aldrig være uden Veiledning. Den
ene sidder høit paa den overnaturlige Himmel og giver
Dig Vink om en Almagtsvillie, som, naar den vil, kan
ryste Jorden og lade Stjernerne falde; den anden lyser
i Naturens uendelige Rum og klarer Dig en Verdensorden
saa uforanderlig, at end ikke det mindste Fnug
kan bevæges deri uden en naturlig Aarsag. Den ene
gaaer op i Øst, den anden i Vest. Ved at følge den
første, oplyses Du om, at Jorden er under Forbandelse,
og at Veien til Maalet er Lidelsernes Vei til Salighed
hisset; ved at følge den anden erkjender Du, at Naturen
er god, som den er, og denne Verden en Tumleplads
for en rig Udvikling af dine Evner og Kræfter. Religionen
har lært Dig, at Du for Gud er en strafskyldig
Synder, evig fortabt, hvis Du ikke ved en Andens Fortjeneste
og Forbøn kan blive benaadet. Fornuften har
ladet Dig indsee, at Du selv maa bøde for dine Overtrædelser,
at ingen Stedfortræder eller vilkaarlig Naade
kan hjælpe Dig ud af de Livsmodsigelser, hvori Du har
hildet Dig selv. Saa veed jeg da, at Du er alsidig underviist.
Styr nu din Kurs efter de ledende Stjerner! Med
dit dobbelte Kompas kan Du trygt pløie Livets Sø og
vente at naae en rolig Havn.“

En saadan Tiltale vil lyde som en Parodi paa vort
hele Opdragelsesvæsen; og dog er den saa fri for al
Overdrivelse, at de antydede Modsigelser blive langt
mere skjærende, naar man gaaer i det Enkelte. Men
Et er Blandingen af det Uforligelige, et Andet er Forbindelsen
af det væsenlig Sammenhørende. Indrømmer
man, at en religiøs Gudserkjendelse er Mennesket lige
saa nødvendig, som en rationel Verdenserkjendelse, og
% --- p. 12 ---
\setcounter{footnote}{0}%
\opage{12}at Aandsdannelsen er betinget af begges Forening, da
har Kulturverdenen med Hensyn paa sit Opdragelsessystem
en stor Opgave at løse.

Ved Underviisningen i de fagvidenskabelige Discipliner
er her i denne Henseende Intet at rette. Valg
af Stof og Methode afhænger af Øiemed og særlige Vilkaar. Alt kommer an paa den Anvendelse, som
Videnskabens Resultater finde i Aandslæren. Vi vælge
til et Exempel \emph{Skabelsen}. Er denne Verden skabt
eller er dens Natur med absolut Nødvendighed, hvad
den er, fra Evighed af? At dømme efter det Hastværk,
hvormed den moderne Kritik har smykket den positive
Religion til Slagtoffer paa Naturalismens Alter, skulde
man rigtignok antage, at Skabelse var et forlorent
Dogme, en blind Forestilling uden fornuftig Mening.
Men er det nu virkelig saaledes? Naturlæren gaaer ud
fra den Grundsætning, at der er Fornuft i Naturen.
Det er mere end en Hypothese, thi hvis det Modsatte
antages, saa er Videnskaben umulig. Vil Nogen maaskee
paastaae, at Fornuft og Natur er Eet og det
Samme, saa at Udtrykket: der er Fornuft i Naturen,
ikke siger Andet, end at der er Fornuft i Fornuften,
Natur i Naturen, da røber en saadan Paastand just
ikke noget overvættes Maal af kritisk Skarpsindighed.
Skal der være fornuftig Mening i Udtrykket, da maa
der nødvendig underforstaaes, at Eenheden af Fornuft
og Natur indeslutter en væsenlig Forskjel. Ville vi nu
et Øieblik standse og besinde os paa, hvad der da
egenlig skal forstaaes ved Fornuften, mon saa ikke den
logiske Nødvendighed skulde føre os til Aanden?

„\emph{Naturlovene},“ siger H. C. Ørsted, „udgjøre i
Forbindelse med det Virksomme, det \emph{Almindelige} i
% --- p. 13 ---
\setcounter{footnote}{0}%
\opage{13}Mangfoldigheden og det Evige i Omskiftningerne.
Naturlovene kunne derfor ogsaa kaldes \emph{Naturtanker}.
Men ligesom der i al Tænken er en Virksomhed, og
Tanken tænkt uden denne kun er en Tænkningsform,
saaledes er der i enhver tilværende Ting en Virksomhed,
ved hvilken den hævder sin Tilværelse og yttrer den
mod andre Ting. Uden Virksomhed har den ingen
Virkelighed. I Tænkningen kalde vi dette Virksomme
Villie, i de udvortes Gjenstande Kraft. Men ligesom
Villen og Tænken er Eet, kun betragtet fra to Sider,
saaledes ogsaa Kraft og Naturlov; de ere Yttringer af
en uendelig Enhed af Kraft og Tanke, af den guddommelige
Villen og den guddommelige \emph{Tænken}\dots{} Det, som
vi erkjende i Naturen, er da ikke noget Fremmed, men
det Samme som det, der lever i os selv.“ --- Skulle vi
altsaa af „det, der lever i os selv“, kunne indsee, hvad
guddommelig Tænken og Villen har at betyde, da maa
det dog vel blive os klart, at, om vi end umulig kunne
forestille os Guddommen som et tænkende og villende
Individ paa menneskelig Viis, saa maae vi dog nødvendig
opfatte den som et tænkende, villende Selvvæsen
d. e. som et fra Naturen forskjelligt Fornuftvæsen. En
evig Tænken uden en evig Tænker, en evig Villen uden
en evig Villende, Tanker, som Ingen har tænkt, Villiesbestemmelser,
som Ingen har villet, slige Antagelser
høre hjemme i Indbildningernes Verden. Naturen kan
lige saa lidt tænke de Tanker, der udtale sig i dens
Love, som den kan ville de „Handlinger“, ved hvilke
disse Love opfyldes. Enhver Ting virker efter sin
Natur; men hvad er en Tings Natur vel Andet end
dens indre ved guddommelig Tænken og Villen bestemte
Anlæg. Da ingen Naturting kan give sig selv Anlæg,
% --- p. 14 ---
\setcounter{footnote}{0}%
\opage{14}saa maa det fornuftige Naturanlæg oprindelig være bestemt
af det Alt bestemmende Selvvæsen, d. e. af den
Alt vidende, Alt formaaende Aand. Gjennem Erkjendelsen
af Tingene og deres Love føres vi til Erkjendelsen
af en fornuftig Verdensorden, igjennem Erkjendelsen
af en given Verdensorden til Erkjendelsen af den
ordnende Fornuft; men den ordnende Fornuft er en
Indbildning, dersom den ikke har Aanden til Grund og
Bærer.

Med samme Ret som vi sige: der er Fornuft i
Naturen, kunne vi ogsaa sige: der er Fornuft i Aanden.
Aand og Fornuft ere Eet, men saaledes, at denne Enhed
indeslutter en Forskjel. Naar vi tale om Fornuften saa
tænke vi allernærmest paa det uendelige Indhold af
Viisdomstanker, Villiesbestemmelser, Begreber og Love,
der i Medfør af Tingenes Anlæg er virkeliggjort i Naturen.
Naar vi tale om Aanden, tænke vi paa det uendelige
Selvvæsen, der tænkende, villende, frembringende
udvikler Fornuftindholdet af sit indre Dyb. Fra dette
Synspunkt kan Forholdet mellem Nødvendighed og
Frihed, Natur og Skabelse opklares. Nødvendigheden
sees i Naturen, for saavidt den aabenbarer sig i Tingenes
fornuftige Sammenhæng, Følgen maa nødvendig
have sin Grund, Virkningen sin Aarsag. At antage en
vilkaarlig Villie, som uden naturlige Aarsager kunde
faae i Sinde at frembringe eller tilintetgjøre naturlige
Ting eller gjøre de naturlige Ting unaturlige, er ganske
vist et Tegn paa Drømmeri og Overtro.

„Man har vel ogsaa“, siger Ørsted, „villet forsvare
Overtroens Drømme mod Naturvidenskaben ved at vise
hen dertil, at en umaadelig større Deel af Naturen er
os mere ubekjendt end velbekjendt. Dette Forsvar kan,
% --- p. 15 ---
\setcounter{footnote}{0}%
\opage{15}selv hos de Uindviede, ikke vise Andet end en Mulighed,
som dog maa være desto mindre betydningsfuld, jo
flere af den gamle Overtroes Indbildninger ere blevne
saaledes belyste, at deres Intethed er bleven afgjort;
men for dem, som ikke havde ladet sig nøie med at
tilegne sig Naturvidenskaben med Hukommelsen og Forstanden,
men havde levet sig ind i dens Heelhed, er
det klart, at vore Indsigter, trods alle Ufuldkommenheder,
dog tilintetgjøre enhver Tanke om, at der i
Naturen skulde foregaae noget Fornuftstridigt.“ Denne
Udtalelse maae vi ubetinget tiltræde.

Indenfor det givne Hele maae Ting og Love nødvendig
svare til hverandre: andre Kræfter, andre Love; andre
Love, andre Kræfter. Men er det givne Heles absolute
Nødvendighed dermed beviist? Ingenlunde! At antage,
at den givne Erfaringsverden maa være den eneste
mulige Verden, fordi ethvert Brud paa dens Love er
fornuftstridigt, er lige saa ufornuftigt som at antage, at
den retvinklede Triangel er den eneste mulige Triangel,
fordi Benegtelsen af den pythagoræiske Læresætning er
fornuftstridig. Den givne Natur er ganske vist et Fornufthele,
men det Heles Eiendommelighed er bestemt
ved dets eiendommelige Anlæg, oprindelig bestemt ved
en almægtig Selvbestemmelse af Aanden som Aand,
altsaa ved en \emph{Skabelse}.

Er der nu ud fra Fornuftlærens Standpunkt aabnet
Udsigt til, at den rationelle Verdenserkjendelse lader
sig forene med en religiøs Gudserkjendelse, saa maa
det vel ogsaa ud fra Religionslærens Standpunkt lade
sig paavise, at den religiøse Gudserkjendelse er forenelig
med en rationel Verdenserkjendelse. Men paa hvilke
Betingelser? Her er Noget, der maa forandres, om%
% --- p. 16 ---
\setcounter{footnote}{0}%
\opage{16}dannes, fra Grunden af omarbeides. Hvad er det?
Dog vel ikke den aabenbarede Religion, dens overnaturlige
Udspring, dens Mirakler og hellige Kjendsgjerninger?
Langt fra! Rationalismen, den nye, saa
vel som den gamle, har i denne Henseende prostitueret
sig mere end tilstrækkeligt. Det er ikke Aabenbaringens
Verdensanskuelse som den staaer for os med alle
Oldtidens Mærker, hvis Karakteertræk skulle udviskes.
Nei, det er noget ganske Andet, der for Opdragelsens
og Aandsdannelsens Skyld maa kræves. Det er vort
indre, væsenlige Forhold til dette Aabenbarede, vor
Forstaaelse, vor personlige Tilegnelse deraf, kort sagt:
vor hele religiøse Bevidsthed, vor Tro, der fra Grunden
af maa omarbeides. Hvad en saadan Omarbeidelse
betyder, hvor den bærer hen, og hvor vidt den lader
sig gjennemføre, vil i det Følgende blive Gjenstand
for nærmere Overveielse.

% --- p. 17 ---
\setcounter{footnote}{0}%
\lecture{Anden Forelæsning.}\opage{17}

Religionen er Grundlag for Aandslivet, Religionen maa
opfattes i Tro. Det ligger da nær at gjøre et Menneskes
Tro til Maalestok for hans Aand: jo stærkere Tro, desto
kraftigere Aand, jo kraftigere Aand, desto stærkere Tro.
Dette er dog ikke altid saaledes. Vel kan der ikke
gives nogen sand og frugtbar Aandsudvikling uden Tro;
men der kan i Virkeligheden gives en haardnakket,
stærk og urokkelig Tro, som dog er saare fattig paa
Aand, ja ude af Stand til at begrunde nogen virkelig
Aandsdannelse.

Troen har nemlig to Hovedsider, ligesom den positive
Religion har to Hovedsider, en udvortes og en
indvortes, en sandselig og en oversandselig, en historisk
i sin Overlevering og en overhistorisk i sit Væsen. I
Begyndelsen ere begge Elementer sammensmeltede, det
Indvortes kommer tilsyne i et Udvortes, og det Udvortes
omsætter sig umiddelbart i et Indvortes; men efterhaanden
skilles de ad, og der indtræder betænkelige
Kriser i Aandslivet.

Det er et Karakteermærke ved alle positive Religioner
uden Undtagelse, at de ved deres første Fremtræden
anmelde sig som guddommelige Aabenbaringer.
Jo mindre Aand der er i en saadan Aabenbaring, desto
% --- p. 18 ---
\setcounter{footnote}{0}%
\opage{18}mere Eftertryk lægges der paa Antagelsen af det Udvortes.
Den rettroende Muhamedaner antager, at Koranen
har existeret fra Evighed i den øverste Himmel
paa den saa kaldte bevarede Tavle, og at den derfra
stykkeviis i en Tid af 23 Aar aabenbaredes Muhamed
ved Engelen Gabriel. De Modsigelser, som findes i den
hellige Bog, løses af Fortolkerne saaledes, at det senere
aabenbarede Vers ophæver det tidligere. Her er en fast
urokkelig Tro paa et Udvortes. Ogsaa Jødedommen
lægger afgjørende Vægt paa det udvortes Skete og det
udvortes Foreskrevne; men til et talende Beviis paa, at
der oprindelig er mere Aand i Jehovatroen end i Muhamedanismen,
traadte Propheterne op mod Præsteceremonierne
og de udvortes Bud, og gjorde i Aandens Navn
det Indvortes gjældende. I Christendommen, Aandens
absolute Religion, maa det Udvortes tjene det Indvortes.
Kristus siger: „mine Ord ere Aand og ere Liv.“ --- „Jeg
er Sandheden.“ --- „I skulle erkjende Sandheden, og
Sandheden skal gjøre Eder frie.“ --- Opgiver man det
positive, historiske Fodfæste, da bliver Inderligheden
mystisk, og den religiøse Bevidsthed svæver i Luften;
men sætter man omvendt Troens Forvisning i sandselig
Oplevelse af et overnaturligt Udvortes, da staaer den
muhamedanske Tro i Liv og Kraft med glødende Farver
paa øverste Trin.

Der fortælles i Korstogenes Historie vidunderlige
Ting om en vidunderlig Skikkelse i Østerlandet, „den
\emph{Gamle paa Bjerget}.“ Hassan-Ben-Sabbah, Overhovedet
for „Ismaeliternes Broderskab“ havde forstaaet
at bibringe sine Undergivne en saa ubetinget Tro
paa
Alkoranens Ord, en saa punktlig Lydighed mod dens
Bud og derved en saa blind Underkastelse under hans
% --- p. 19 ---
\setcounter{footnote}{0}%
\opage{19}egen hellige Autoritet, at han forbausede Verden. Hans
Troende foragtede verdslig Ære, Lidelse og Pinsler, ja
Livet selv, for at lyde hans mindste Vink. Det var
Opdragelse. En Grev Henrik af Champagne drog igjennem
Ismaeliternes Land, besøgte den Gamle paa hans
Borg og fik en hæderlig Modtagelse. Paa ethvert af de
høie Taarne stode to Mænd paa Vagt stive som Støtter.
Hassan gav et Par af disse Vagthavende et Vink, og
øieblikkelig styrtede de sig fra Taarnet og laae sønderknuste
for den forbausede Korsfarers Fødder. Den Gamle
yttrede koldblodig: „Hvis I skulde ønske det, vilde I
see dem alle paa mit Vink saaledes styrte sig til Jorden.“

Med denne ubetingede Tro, denne blinde Lydighed
havde det, efter Beretningerne, et eget Sammenhæng.
Det var saa langt fra, at de vagthavende Mænd paa
Taarnene imødesaae et saadant Øieblik med Rædsel eller
betragtede den Gamle som en grusom Tyran, man kun
adlød af Skræk, at de tvertimod ærede ham i deres
Sjæls Inderste og ansaae en Afløsning, som den ovenfor
beskrevne, for den lykkeligste af alle. De nærmere
Omstændigheder fremstilles saaledes: Ismaeliternes Sekt
var oprindelig deelt i to Klasser, Mestre og Lærlinge,
til disse føiede Hassan en tredie Klasse, der skulde
være uvidende om Sektens egenlige Hemmeligheder.
Medlemmerne af denne Klasse kaldtes Offrene; de
vare klædte i Hvidt, med røde Huer og Skjærf, og
maatte bestandig holde sig omkring Stormesteren, for
at forsvare og hævne ham. Midt i de Lande, den
Gamle raadede over, laa et Slot, et Muhameds Paradis
paa Jorden. Yndige Haver, løvfulde Træer, Blomster
beduggede af Vandspringenes Støvregn bød den Indtrædende
velkommen, skyggefulde Gange skjærmede  ham
% --- p. 20 ---
\setcounter{footnote}{0}%
\opage{20}mod Solens Brand. Boligen forinden straalede i Guld
og Marmor; Arabesker, Silke og kostbare Tepper i de
høie, kjølige Sale, unge Piger hvilende paa bløde Divaner
i vellystig Yppighed! her var Alt, hvad en troende
Sjæl kunde ønske og attraae. Naar nu en Yngling skulde
indvies til Optagelse i Offrenes (Fedawiernes) Klasse,
skete det paa følgende Maade. Han blev i den Gamles
Nærværelse bedøvet ved en Drik, tilberedt af Haschisch
(Bladene af indisk Hamp), og i denne bevidstløse Tilstand
ført til det fortryllende Sted. Naar han da slog
Øinene op, maatte han nødvendig troe, at han var i
Paradis og nød i fuldt Maal alle Salighedens Glæder.
Han blev dernæst atter bedøvet og bragt tilbage. Den
Gamle stod igjen ved hans Side og forsikkrede den Opvaagnede,
at han ikke legemlig havde været af Stedet;
forholdt det sig desuagtet som han sagde, da maa hans
Sjæl ved et Mirakel være blevet henrykket til Paradis.
Saa havde den unge Mand Troen i Hænderne; tungsindig
længtes han efter den Stund, da han skulde komme
derind for evig Tid, og var villig til at underkaste sig
de Betingelser, den hellige Autoritet foreskrev ham,
i ubetinget Lydighed.

Selv indenfor Aandens Religion er der noget Muhamedansk
i Kirketroen, særlig den katholske. Hvad
vi her forstaae ved det Muhamedanske ligger ikke fra
Først af i Læren, i Dogmerne, i Forskrifterne, i Ceremonierne,
ikke i det, som er Troens Gjenstand, men
oprindelig i Troen selv, i den Maade, hvorpaa Troen
opfattes og bestemmes. Vil man maaskee indvende, at
den Maade, hvorpaa Troens Væsen opfattes og bestemmes,
selv er et Dogme, da bevæger man sig kun i en
Cirkel. Man kan umulig foreskrive Tro paa Troens
% --- p. 21 ---
\setcounter{footnote}{0}%
\opage{21}Dogme uden at forudsætte et personligt Indre, en Samvittighed,
som kan troe sig forpligtet til at antage Dogmet.
Hvor megen Overtalelse eller Tvang man end
anvender, saa kan Antagelsen dog ikke blive til Tro,
med mindre der i Sjælen selv er Noget, som siger Ja
og Amen til Troens Dogme. Troen er altsaa i første Instans
en oprindelig Frihedsyttring, en Selvbestemmelsesact, en
Sandhedserkjendelse af personlige Grunde. Hvori bestaaer
nu det Muhamedanske? Ligefrem deri, at det personlige
Indtryk af Troens Gjenstand, som skulde danne
Udgangspunktet for en Tilegnelse i Aand og Sandhed,
med Eet forvandles til et Motiv for ubetinget Underkastelse
under en udvortes Autoritet. „Du maa dog føle
med Dig selv og erkjende, at Du er en falden, hjælpeløs
Skabning, der trænger til Opreisning og Frelse“; naar
denne Henvendelse til Menneskets Samvittighed har
gjort sin Virkning, saa følger Opfordringen: „grib Frelsen,
der tilbydes, som den Skibbrudne griber Redningsplanken,
lad Dig optage i den ene saliggjørende Kirkes
Skjød!“ Nu er det ude med Selvforstaaelsen og den
frie Tilegnelse. Tro paa Kirken, blind Underkastelse
under dens Autoritet, ubetinget Lydighed mod dens Bud
er Frelsens Betingelse. Vistnok forsømmer Kirken ikke
at oplyse sine Troende om, hvad det er, de skulle troe.
Den lærer dem hvad Gud er, hvad Hellighed er, hvad
Synd er; det kan Altsammen forstaaes tydelig til Punkt
og Prik. Men fra al denne Forstaaelse er der Et, som
væsenlig er udelukket, og det er Selvforstaaelsen. Det
er Kirkens Gud, Kirkens Hellighed, hvorpaa der troes.
Den Troende veed slet ikke, hvad Gud er i sig selv;
han veed kun, hvad Kirken lærer, at Gud er. Han veed
ikke engang hvad Synd er og hvad det er at synde;
% --- p. 22 ---
\setcounter{footnote}{0}%
\opage{22}Kirken maa bogstavelig lære ham det. Hvad Kirken
stempler som Synd, det er Synd, om end den naturlige
Følelse siger det Modsatte; naar Kirken har udslettet
en Synd, saa er den udslettet, hvormeget den oprindelige
Sands for Ret og Uret end kan have at indvende.
Hemmeligheden er, at den Troende ingen Samvittighed
har for sig selv, at han har givet Kirken sin Samvittighed.
Autoritetstro er Udvorteshed i det Indvortes;
Formlen for denne Udvorteshed er Dogmet.

Man har slaaet Allarm mod den katholske Kirkes
mange historiske Overgreb i det Udvortes; og det er
sandt, de have været oprørende og kunne maaske endnu,
naar „den moderne Bevidsthed“ ret har hjulpet til og
gjort sit Bedste, blive oprørende nok; men det er Magthaveriets
almindelige Overgreb, der mødes og standses
med udvortes Vaaben. Det radikale, aandsfortærende
Overgreb er Forvandlingen af det Indvortes selv til et
Udvortes. Dette Muhamedanske, at Autoritetstroen tager
det indvortes Menneske til Fange, lokker det ud af sig
selv, gjør det fremmed for sig selv, saa at det, bedøvet
som ved en Trylledrik, phantastisk og dog forstandig,
kan finde Saligheden udenfor sig, og derved opnaae en Art
sandselig Forvisning om dens Realitet, er en bedaarende
Illusion, en uoverstigelig Skranke for Aandsdannelsen.

Man har gjort Ophævelser over de mange vilkaarlige
Tilsætninger i den katholske Overlevering, dens
opdigtede Mirakler, dens Legender og Overtro; ikke
uden Grund. Miraklerne ere dens overnaturlige Bygningsemner,
dens Ornamenter; Legenderne snoe sig som
Vedbende op ad dens Mure: den er luxuriøst udstyret
med Overtroens Opfindelser. Og dog ere Overtroens
Gjenstande endnu ikke Overtroens Væsen, men kun
% --- p. 23 ---
\setcounter{footnote}{0}%
\opage{23}Udslag deraf. Overtroen sidder i Helligdommen, i Troen
selv som Autoritetstro.

Med alt dette er der dog Aand i det katholske Væsen.
Paa Autoritetens og den blinde Lydigheds Grundlag have
Jesuiterne viist, hvad der kan opnaaes ved mystisk comprimeret
Inderlighed og et consequent Opdragelsessystem.
En Karakteer som Franz Xaver skal man forgjæves lede
om i noget andet Kirkesamfund. Et saadant Troesmod,
en saadan Udholdenhed i Forsagelse og Selvfornegtelse,
en saadan Bekymring for at vinde Sjæle maa i deres
Øine, som eensidig see paa det Udvortes, stille denne
Jesuitmissionær paa lige Trin med Apostelen Paulus.
Men efter sit Indvortes, hvad er han saa? En Incarnation
af den Autoritetstro, som fra Grunden af magtstjæler
det personlige Selv. Det er ikke hans eget menneskelige
Væsen, der under Kirkens Indflydelse er kommet
til Udvikling, det er tvertimod det kirkelige Væsen, som
udenfra er trængt ind i ham. Kirken har taget hans
Hjerte, hans Villie ud af hans Bryst og sat sit Hjerte,
sin Aand, sin Villie i Stedet. Han troer i Kirkens Navn
Kirkens Tro, hans Verdensforsagelse er i ham Kirkens
Forsagelse, hans Omvendelsesiver Kirkens Iver, selv er
han det blinde Redskab, „villieløs som et Cadaver“,
lydig „som Staven i Oldingens Haand“. Hvad han gjør
er ganske vist noget langt Høiere og Bedre end hvad
hine Mænd gjorde, der styrtede sig ned fra Taarnet,
fordi det, hans Autoritet har foreskrevet ham, er høiere
og bedre; men Motivet er det samme: ubetinget Lydighed,
blind Underkastelse under en udvortes Autoritet.

Ved sin Autoritetstro og sit Opdragelsesvæsen har
den katholske Kirke sat et evigt Fjendskab mellem sig
og den menneskelige Fornuft. Vel maa det synes, som
% --- p. 24 ---
\setcounter{footnote}{0}%
\opage{24}om den systematiske Udvikling og Begrundelse af det
kirkelige Lærebegreb indeholdt en væsenlig Anerkjendelse
af Fornuftens Ret; men Anerkjendelsen er kun et Skin.
Vistnok har der været og er endnu store Videnskabsmænd,
berømte Lærde, ikke blot Historikere og Philologer,
men ogsaa Mathematikere, Astronomer, Physikere,
ja selv Philosopher i Kirkens Skjød. Men see vi nærmere
til, saa ere de Videnskaber, som bevæge sig paa
et verdsligt Omraade, kun til en vis Grad tolererede.
Kirken kan, efter Omstændighederne, lukke et Øie til
for videnskabelige Theorier og Hypotheser, saa længe
de ikke aabenlyst sætte sig op mod dens Autoritet;
men det skeer altid med den stiltiende Forudsætning,
at dens egen ufeilbare Viden er uendelig ophøiet over 
„denne Verdens Børnelærdomme“. Hvad Tænkningen i
Forhold til Dogmet angaaer, er Fornuftens Afhængighed
ligefrem slavisk. Fornuften er forpligtet til at indrømme,
at Kristus har to Villier, fordi Kirken siger, at han
har to; kvis Kirken sagde, at han havde halvtredie Villie, % PRINTED AS IS: kvis (for hvis; read at the image)
saa var Fornuften ikke alene nødt til at antage det
for Sandhed, men endog at gjøre sit Bedste for at faae
det beviist. At det Brød, som, før det indvies af Præsten,
er naturligt Brød, efter Indvielsen kan beholde sin
naturlige Smag, sit naturlige Udseende og sine naturlige
Egenskaber, og dog ikke længere være naturligt
Brød, men være forvandlet til det Legeme, hvormed
Kristus blev korsfæstet, er et saa vældigt Slag i Fornuftens
Ansigt, at det maa sortne for dens Øine og
overbevise den om, at den katholske Autoritetstro er
lige saa uforenelig med en rationel Verdenserkjendelse
som Alkoranens Aabenbaringer.

% --- p. 25 ---
\setcounter{footnote}{0}%
\opage{25}Med Reformationen begyndte Troesdogmet at undergaae
en væsenlig Omdannelse. Luther brød med den
katholske Autoritet, hævdede Samvittighedens Ret og
opfattede Troen som personlig Tilegnelse af Kristi Fortjeneste,
altsaa som Grundfactor i det religiøse Livs
Inderlighed. For et Øieblik saae det ud, som skulde
det Ydre med Eet nedsættes til Stof og Materiale for
det sig fremarbeidende Indre, og Skriftens hellige Ord,
gjennemglødet af Aand, tænde Selvforstaaelsens Lys i
Sjælen. Men det varede ikke længe, før Troens Gjenstand
drog Opmærksomheden bort fra Bevægelserne i
det troende Selv. Det stærke, religiøse Liv blev lidt
efter lidt til en Lære om Livet. Den protestantiske
Kirke fik sin egen Theologie og sin egen Dogmatik i
Lighed med den katholske: Udvortesheden feirede en
ny Triumph, da Kirken blev Statskirke og Orthodoxien
naaede saa vidt, at den fik sine Dogmer thinglæste af
Politimyndigheden. Under kongelig Varetægt var Autoritetstroen
nu bosat i Skriften, væsenlig lige saa udvortes
som forhen i Pavedømmet.

Den Paastand, at Noget er sandt, fordi det staaer 
i Bibelen, er lige saa muhamedansk som, at Noget er
sandt, fordi det staaer i Koranen. Den orthodoxe Bibeltro
kan lige saa lidt faae Øie paa den Cirkel, hvori den
bevæger sig, som den orthodoxe Pavetro. Den, der
troer paa Kirkens Ufeilbarhed, bilder sig ind, at han,
frigjort for al menneskelig Feilbarhed, har det Ufeilbare
over for sig, og mærker ikke, at det er ham selv, der
lægger Ufeilbarheden til af sit Eget. Han troer paa
Kirken, fordi den er ufeilbar, og see, saa er det kun hans
egen Tro paa dens Ufeilbarhed, der er ufeilbar. Ligesaa
med den, der udvortes bygger sin Tro paa Skriften. Han
% --- p. 26 ---
\setcounter{footnote}{0}%
\opage{26}vidste jo slet ikke, hvad han skulde troe, dersom ikke den
ufeilbare Skrift lærte ham det; men han veed heller ikke,
at det er hans egen --- ufeilbare! --- Tro, der lærer ham, at
Skriften er ufeilbar. Saa beraaber man sig fra begge
Sider paa Aandens Vidnesbyrd; men med Autoritetstroen
faaer man kun Dogmet om Aanden, ikke Aanden
selv, kun Dogmet om Vidnesbyrdet, ikke Vidnesbyrdet
selv. I Aandens personlige Vidnesbyrd maa. det Guddommelige % PRINTED AS IS: a full stop after „maa“ (read at the image, T011, z01)
lige saa fuldt bekræftes af det Menneskelige,
som det Menneskelige af det Guddommelige.

Det er en Kjendsgjerning, at den lutherske Orthodoxie,
fastholdt i dens positive Dogmer, staaer i samme
fjendtlige Forhold til Nutidens Kultur og Videnskab som
den katholske. Det er ikke alene i det Physiske og det
Historiske, men selv i det Psychologiske, i Bevidsthedslivet,
at Hindringerne for en gjensidig Forstaaelse vise
sig. Man tænke paa Læren om Menneskenaturens absolute
Fordærvelse ved Arvesynden, om Guds almindelige
Naade, som dog ikke er almindelig, men kun kommer de
Udvalgte tilgode; paa Læren om de to Naturer i Kristus, om
hans Sæde ved Faderens Høire, hans Legemes Allestedsnærværelse
og dog begrændsede Nærværelse i Nadveren
o. s. v., og man vil, uden at overanstrænge sin Forstand,
snart opdage, at om end Mange, Mange have troet bogstavelig
paa alt dette, saa er der dog Ingen, der i Aand
og Sandhed har troet det, af den simple Grund, at
Ingen har kunnet troe det. At troe paa Noget er at
have det i Dogmets Form som Troens Gjenstand; men
at troe Noget er ved inderlig Tilegnelse at forvandle
det til Troens Indhold. Hvad Troesgjenstande angaaer,
skal der neppe kunne opvises noget saa Urimeligt og
Absurd, at det jo fra en eller anden Side kan blive
% --- p. 27 ---
\setcounter{footnote}{0}%
\opage{27}Gjenstand for en Autoritetstro: det Absurde er jo netop
det, som alene kan antages paa Autoritet. Anderledes
med Troens Indhold. Hvad der inderlig skal optages
i den religiøse Bevidsthed til aandelig Næring for Siælen, % PRINTED AS IS: Siælen (for Sjælen; read at the image)
maa som ethvert andet Næringsmiddel lade sig opløse
og fordøie; det maa assimilere sig med Bevidsthedens
Væsen: Forstaaelsen deraf maa gaae i Eet med en personlig
Selvforstaaelse. De positive Dogmer, de være
for Øvrigt knudrede og kantede eller runde og slebne,
ere i deres faste Tilstandsform som Stene, klare Krystaller,
maaske Ædelstene! men ikke Livets Brød. Man
kan hænge dem udenpaa Bevidstheden og pynte sin
Sjæl med dem; men fordøie dem kan man ikke.

Forholdet mellem Fornuftskolen og Kirkeskolen viser
sig i en ny Belysning, naar Troen opfattes som personlig
Tilegnelse af de positive Dogmers almeenreligiøse Indhold.
Det er ingenlunde Opgaven at afskaffe Kirketroen
eller i nogensomhelst udvortes Forstand at angribe Autoriteten
som Autoritet. Tvertimod! I udvortes Forhold
kan Autoriteten ikke alene være ærværdig, men endog
nødvendig. End ikke den rationelle Underviisning kan
undvære Autoriteten. Discipelen maa paa Lærerens
Autoritet antage bestemte Læresætninger, og i Tillid til,
at det, der endnu ligger over hans Forstand, vil blive
ham forstaaligt, arbeide sig frem til Indsigt. Autoritetens
Vaabenmærke bør staae over Indgangen til Undersøgelsernes
Skole. En Videnskab, der endnu ikke har vundet
Anseelse i den offentlige Mening og derved en vis Autoritet,
vil være udsat for de mest overfladiske Indvendinger,
uden at Nogen bryder sig om at undersøge
Sagen. Gaaer det saaledes med det Rationelle, som
dog maa kunne hævde sig i Kraft af almeengyldige
% --- p. 28 ---
\setcounter{footnote}{0}%
\opage{28}Beviser, hvorledes maa det da ikke gaae med det
Religiøse, som kun kan erkjendes af personlige Grunde.
Har Bibelordet ingen Autoritet, Kirken og dens Dogmer
ingen Ærværdighed, ingen Anseelse, da er Religionen
priisgiven, og Enhver kan haane den efter Indfald og
Lune. Sandheden selv, den evige, der er, før Abraham
blev, maatte jo iføre sig Kjød og Blod og træde frem
for Menneskene i menneskelig Skikkelse som Autoritet.

Men selv den absolut sande Autoritet er kun et
Udvortes; Sandheden, det evige Væsen, er et Indvortes,
Sandheden kræver Erkjendelse, og for Erkjendelsens
Skyld nedsætter den det Udvortes til Middel og giver
sin Autoritet hen til Opløsning. „I skulle erkjende
Sandheden, og Sandheden skal gjøre Eder frie.“ At
Autoriteten opløses, betyder da ikke, at den afskaffes,
nei, det betyder, at den i Sandhed stadfæstes. Kun
med den falske Autoritet er det forbi, naar den opløses.
En falsk Autoritet, Theologien f. Ex., maa svækkes desto
mere, jo mere dens Videnskabelighed prøves, thi ethvert
Forsøg paa at erkjende dens Grunde opløser sig i det
Absurde. Maaske vil man indvende, at, naar Theologien
mister sin Autoritet, saa maa det gaae ud over Kirken
og Dogmerne. Det vilde være sandt, dersom Kirkens
Autoritet var bygget paa Theologien; det rationelle
Tøbrud, der opløser de sidste Stumper af den theologiske
Autoritet, vilde da ganske vist drage Opløsningen
af den kirkelige efter sig. Men Kirkens Autoritet er,
naar vi tænke os om, ikke bygget paa Theologien, men
paa Religionen; det er en Autoritet ikke for Videnskaben,
men for Troen, en Troesautoritet. Hvorvidt denne
Troesautoritet er sand eller falsk, beroer paa, hvorvidt
dens Dogmer ved religiøs, personlig Tilegnelse opløse
% --- p. 29 ---
\setcounter{footnote}{0}%
\opage{29}sig i Sandheden eller i Usandheden. Her er en Opgave
for Aandsdannelsen, et Grundproblem for religiøs-rationel
Opdragelse. Nogle Exempler.

\emph{Natur og Skabelse}. „I Begyndelsen skabte Gud
Himlen og Jorden“; disse Bibelens Ord med den derpaa
følgende Skabelseshistorie have kirkelig Autoritet.
„Naturen er uskabt; det er gaaet naturligt til med alle
Dannelser i Himlen og paa Jorden“; denne Paastand
har, ifølge den naturalistiske Synsmaade, videnskabelig
Autoritet. Kan nu en ubetinget Tro paa nogen af disse
to Autoriteter træde i Sandhedserkjendelsens Sted?
Umuligt! Lader man den kirkelige Autoritet udelukke
den videnskabelige, saa kan man jo let slippe fri for
Bryderier med det Fornuftige. Man lægger sit Hoved
til Hvile paa Autoritetens Dovnepude, Fornuften
slumrer ind, og den blidelig indslumrede Fornuft forstyrres
ikke let at af den rationelle Videns lydløse Fjed. % PRINTED AS IS: „at af“ (an extra „at“; read at the image, T013)
Man drømmer om Skabelsesmiraklet og Adam i Paradis;
men slige Drømmerier føre ikke til personlig Sandhedserkjendelse
i Troen. Lader man den videnskabelige
Autoritet udelukke den kirkelige, da har man unegtelig
en Spore til rationel Undersøgelse af alt det Givne i
dets Sammenhæng; men hvorfra Sammenhængen og det
Sammenhængende er, forbliver en Gaade; thi den personlige
Selverkjendelse mangler. Har Religionen ingen
Autoritet, saa bliver man ved Spørgsmaalet om Tingenes
Oprindelse alt for let færdig med Friheden. Natur er
Nødvendighed og Nødvendigheden Natur, Naturnødvendigheden
opsluger Friheden, først den guddommelige og
dernæst den menneskelige. Har Videnskaben ingen
Autoritet, saa bliver man alt for let færdig med Naturen.
Skaberen er almægtig og kan gjøre Alt, hvad han vil
% --- p. 30 ---
\setcounter{footnote}{0}%
\opage{30}med den skabte Natur, Menneskets Frihed er ikke
væsenlig afhængig af Naturens Love, men af Skaberens
frie Villie: det Hele løser sig op i guddommelig Vilkaarlighed.
Kun begge Autoriteters forenede Tryk kan sætte
alle Undersøgelsens Hjul i Gang og drive Bevidsthedslivet
frem fra Gudserkjendelse til Verdenserkjendelse,
fra det Religiøse til det Rationelle, fra Tro til Viden
og fra Viden til Tro.

\emph{Synd og Forløsning}. Religionen vil med kirkelig
Autoritet indestaae for, at Synden er Folkenes Fordærvelse,
at det syndige Menneske ikke ved egen Kraft
kan reise sig af sit Fald, men trænger til en guddommelig
Hjælper og Forløser. Fornuften vil med Videnskabens
Autoritet indestaae for, at den almindelige Syndighed
er en simpel Omskrivning af Menneskenaturens nødvendige
Ufuldkommenhed, at Frigjørelsens Betingelser
maae søges i denne Naturs Perfectibilitet, at Menneskeslægten
ved en fornuftig Brug af sine naturlige Kræfter
kan og maa i Tidernes Løb lidt efter lidt blive sin egen
Frelser. Kaster man nu Vrag paa Fornuftens Autoritet,
hvad saa? Saa voxer den dogmatiske Forestilling om
Synd og Fordærvelse op til en saadan Høide, at den
formørker hele Tilværelsen og ender med at gjøre
den religiøse Selverkjendelse umulig. Intet Troens Lys
er i Stand til at opklare Fordærvelsens Afgrund. Opfatte
vi Troen, ikke som en Tro paa Dogmet, en blind
Autoritetstro, men som en personlig, paa Selvforstaaelse
begrundet Tilegnelse af det aabenbarede Ord, da er
der visselig Ingen, der kan eller har kunnet troe, at
Menneskenaturen ved Adams Synd er bleven saa grundfordærvet,
at Slægtens Individer uden Undtagelse ere
hjemfaldne til evig Fordømmelse. Forkaster man Kirkens
% --- p. 31 ---
\setcounter{footnote}{0}%
\opage{31}og Dogmets Autoritet, saa bliver ikke alene Arvesynden
til en forskruet Indbildning, men selv de virkelige Synder
finde i de naturnødvendige Drifter saa mange Undskyldninger,
at Synd og Skyld blive til tomme Talemaader,
Syndsbevidstheden et Tegn paa blødagtig Vellyst. Der
kan hos Kulturens Ordførere endog forekomme saa overfladiske,
ja saa brutale Haansyttringer om Synd og Skyld,
at den moralske Tilregnelighed udslettes. Hvorfor?
Fordi man ikke finder det Uleiligheden værd at gaae
i sig selv og undersøge Syndsbevidstheden nærmere. ---
Begge Autoriteter, Kirken og Videnskaben, ere uundværlige,
naar Mennesket under indre Samvirken af Tro
og Fornuft skal opdage Syndens og Forløsningens Hemmelighed,
og faae sit personlige Væsen opklaret.

\emph{Det evige Liv}. Vi have Skriftens og Kirkens
Autoritet for, at Tilværelsen af et evigt Liv er historisk
godtgjort ved Kristi Opstandelse. Vi have Naturlærens
videnskabelige Autoritet for, at Opstandelseshistorien
staaer i Strid med Naturlovene. Skal nu Kirkeautoriteten
udelukkende gjælde, da er Forvisningen om det evige
Liv bygget paa Opstandelseshistorien, hvis Autoritet er
bygget paa Bibelen, hvis Autoritet er bygget paa
Kirkens Autoritet, som igjen er bygget paa Historiens
Autoritet. Det er en muhamedansk Forvisning
i Kraft af et Udvortes. Skal Videns Autoritet uden
videre gjælde, saa hjemfalder Historien med Dogmet til
Kritiken, og Resultatet bliver, at det ikke er Livet,
men Ideen om Livet, der er det Evige. I begge Tilfælde
søger man Visheden om Selvet udenfor Selvet, fra kirkelig
Side i en Historie, fra videnskabelig Side i en Naturbetragtning:
den personlige Forstaaelse af det personlige
Liv er udelukket. Den videnskabelige Alvor, hvormed
% --- p. 32 ---
\setcounter{footnote}{0}%
\opage{32}de organiske Betingelser for Sjælens Existens undersøges,
maa svare til den Troens Alvor, hvormed den religiøse
Inderlighed omfatter det personlige Livs Fuldendelse,
og omvendt. Fornuftundersøgelsen skal ikke indskrænkes
eller paa endelig Maade afpasses efter det Religiøse,
saa lidt som Troesundersøgelsen skal begrændses og
tillæmpes efter det Fornuftige. Opløsningen maa fra
begge Udgangspunkter være fuldstændig. Kun idet
begge Autoriteter opløses, og Opløsningen viser, at det
Religiøse netop er Muligheden for det Fornuftige, og
det Fornuftige Muligheden for det Religiøse, at begge
Muligheders Eenhed er det personlige Livs oprindelige,
for Evigheden bestemte Selveenhed, kan Troen forenes
med Viden, Religionen med Videnskaben.

Som eneste Udvei til at forlige Kulturbevidstheden
med den religiøse Trang har man begyndt at kalde den
for mere end en Menneskealder siden hensovede Rationalisme
tillive igjen ved at indblæse den Tidens Aand.
Men Forsøget vil neppe lykkes. Rationalisten har Ret,
naar han fordrer, at Orthodoxiens Dogmer maae opløses,
men han har Uret i at opløse dem i Viden, ikke i Tro;
at opløse dem i Viden er at knække dem som hule
Nødder; den religiøse Kjærne er borte. Rationalisten
vil have sine Troessætninger begrundede i Videnskaben
og mærker ikke, at han indsmugler noget Personligt,
som Videnskaben ikke kan begrunde. Han forudsætter
en hellig, retfærdig og kjærlig Gud, og mærker ikke,
at han herved er kommet langt ud over Videns Grændser.
Begreber som: Guds Retfærdighed, Guds Kjærlighed,
Guds Naade, Barmhjertighed o. desl. ere ikke Vidensbegreber,
men Personlighedsbegreber, Troesbegreber.
Kjærlighedens retfærdige Gud er en Personlighedens Gud,
% --- p. 33 ---
\setcounter{footnote}{0}%
\opage{33}Hjertets Gud, Samvittighedens Gud. Rationalisten antager,
at Læren om Sjælens Udødelighed er fornuftig,
i sin Grund videnskabelig, og mærker ikke, at han
sammenblander det, som Videnskaben ikke kan modbevise,
med det, som den kan bevise.

Mangel paa Indsigt i Forholdet mellem det Personlige
og det Upersonlige, det Religiøse og det Fornuftige,
er Mangel paa Aandsdannelse, Mangel paa Opdragelse.

% --- p. 34 ---
\setcounter{footnote}{0}%
\lecture{Tredie Forelæsning.}\opage{34}

Det er os fra det daglige Liv bekjendt, hvor let der
kan opkomme Tvist om Forholdet mellem Ord og Mening.
Man bør, som det hedder, ikke hænge sig i Ord, men
see paa Meningen; og dog har Meningen Tilhold i
Ordet. Man skal, som det udtrykkes i høiere Stiil, betænke,
at Aanden levendegjør, men Bogstaven ihjelslaaer;
men ved juridiske Kjendelser og i lovgivende Forsamlinger
skal man da ogsaa betænke, at Aanden hæfter
ved det bogstavelige Udtryk og vilde uden dette være
hjemløs og uvederhæftig. Striden om Ord og Mening,
Aand og Bogstav faaer i den protestantiske Kirke afgjørende
Indflydelse paa Fortolkningen af den hellige
Skrift. Det 16de og 17de Aarhundredes Dogmatikere
holdt Aanden fast i Bogstaven og paastode, at Gud
selv ordret havde dicteret de hellige Forfattere, hvad
de skulde skrive. I den Forstand var Skriftens Ord
Guds Ord, at Propheter, Evangelister og Apostle, kun
havde været „Guds Skrivere“, „Guds Penne“. Der kan
altsaa ikke være end den mindste Feil eller Uovereensstemmelse
i det skrevne Aabenbaringsdocument, hverken
hvad de historiske Kjendsgjerninger eller hvad Udtrykket
for det Religiøse angaaer: den hellige Skrift er i Et og
% --- p. 35 ---
\setcounter{footnote}{0}%
\opage{35}Alt ordret absolut ufeilbar. Hvorledes det siden er
gaaet Theologien med Bibelforgudelsen og den absolute
Ufeilbarhedstro, skulle vi ikke opholde os ved: vi forbigaae
Fataliteterne med det Udvortes og fæste Opmærksomheden
paa det Sjælelige, det Indvortes.

Uden nærmere Kjendskab til Kritiken af de bibelske
Skrifter og det kirkelige Dogme om Inspirationen
vil Enhver, som med Nutidens almindelige Dannelse
forener en religiøs Trang og Ærbødighed for Skriftens
Ord og gjennemlæser det Gamle og Ny Testamente, overbevise
sig om, at Aabenbaringens Aand ikke lader sig
fange i Bogstaven. Der staaer i anden Samuels Bog
(24,1): „Jehovas Vrede traf igjen Israel; han fristede
nemlig David“; der staaer i første Krønikernes Bog:
(21,1): „Satan var Israels Modstander og han fristede
David til at tælle Israel.“ Begivenheden er den samme.
Holder man sig nu strængt til Bogstaven, saa faaer
man jo ud, at Jehova og Satan er eet og samme Væsen.
Skiller man Jehova fra Satan, da har man med det
Samme skilt Aanden fra Bogstaven, Meningen fra Ordet
og anvendt en Fortolkning. Der staaer hos Propheten
Hoseas (11,1): „Israel var endnu ung, da jeg fik ham
kjær; ud af Ægypten kaldte jeg min Søn.“ Der staaer
hos Matthæus (2,15), at Josephs Flugt med Moder og
Barn til Ægypten skete, „at det maatte gaae i Opfyldelse,
som Herren talede ved Propheten, der siger:
Fra Ægypten kaldte jeg min Søn. Skal dette forstaaes % PRINTED AS IS: the quotation opened at „at det maatte“ is not closed (T018)
bogstaveligt, da er Evangelisten i Modsigelse med Propheten,
og Aanden i Modsigelse med sig selv. Det er
hverken den samme Søn ei heller den samme Begivenhed.
Den Søn, Propheten taler om, er Jødefolket, der i sin
Ungdom var i ægyptisk Trældom. Propheten  seer til%
% --- p. 36 ---
\setcounter{footnote}{0}%
\opage{36}bage paa det Forbigangne,  men Evangelisten siger, at
han seer fremad mod et Tilkommende. Efter Bogstaven
skulde altsaa Aanden, naar den bogstavelig taler om
det Forbigangne, bogstavelig tale om det Tilkommende.
Det er for Meningens Skyld nødvendigt at anbringe en
Fortolkning.

Fortolkningen kan ikke være ufeilbar, med mindre
Læseren kan sikre sig en ufeilbar Fortolker. At Aanden
i Skriftens Ord er ufeilbar, kommer os kun tilgode
under den Betingelse, at den kan gjøre vor Forstaaelse
deraf ufeilbar. Vil nu Skriftens Læser mene, at Aanden
staaer i selv samme Forhold til ham som til Propheterne
og Apostlene, siger ham ordret, hvad han skal troe og
gjøre, da er hans Læsning af Skriften jo i religiøs Forstand
overflødig. Maa han tilstaae for sig selv, at
Aanden har virket i de hellige Forfattere paa en ganske
anden Maade end den virker i ham, da er der herved
sat en psychologisk Skranke for Forstaaelsen. Opgaven
er, at bestemme det rette Forhold mellem det Oprindelige
og det Afledede.

Man klager over det Tvivlsomme i de bibelske
Beretninger, Modsigelser i det Historiske o. s. v., altsammen
noget Underordnet. Selv om alle Vanskeligheder vare hævede, og Bibelen fra Først til Sidst i
ordret Overeensstemmelse med sig selv, det vilde dog
ikke hjælpe os det allermindste, saalænge Vexelvirkningen
mellem den guddommelige og den menneskelige
Aand vedbliver at være uforklarlig. Pintsefestens Under
er et talende Exempel. Opfattes Aandens Udgydelse
umiddelbart efter Bogstaven, da kan det i Karakteer
af Mirakel vel blive Gjenstand for Autoritetstro; men
hvorledes vil det gaae med den personlige Tilegnelse?
% --- p. 37 ---
\setcounter{footnote}{0}%
\opage{37}„Ere ikke alle disse, som tale, Galilæere? Og hvorledes
høre vi dem da hver især i vort Maal, hvori vi
ere fødte? Parther og Meder og Elamiter og vi, som
boe i Mesopotamien og Judæa og Kappadocien, Pontus
og Asien. I Phrygien og Pamphylien, Ægypten og de
Egne af Lybien ved Cyrene og vi her boende Romere,
Jøder og Jødevenner, Kreter og Araber --- vi høre dem
tale i vore Tungemaal?“ (Ap. G. 2, 7 flgd.)
For Autoritetstroen er her nu slet ingen Vanskelighed; Gud er jo almægtig og kan gjøre Alt, hvad
 han vil. Men for Troen som personlig Tilegnelse er
 her en stor psychologisk Vanskelighed. Griber Guds
Aand, naar den vil, paa saa overnaturlig Maade ind i
det sjælelige Liv, at den pludselig, uden al Forberedelse,
meddeler Bevidstheden Færdighed i fremmede Sprog,
da staaer jo Inspirationsmiraklet paa samme Trin som
Miraklet med Bileams Aseninde. Den samme Aand,
der ved overnaturlig Indvirkning kan faae en Menneskesjæl
til at tale et fremmed, ham forhen ubekjendt Sprog,
kan ogsaa faae en Dyresjæl til at tale Menneskesprog.
Men hvad er saa Sjælens Natur, og hvor ere vi med
vor rationelle Bevidsthed? Et virkeligt Mirakel kan
ikke begribes, det er sandt; men skal det ikke staae
for os som noget vildt Fremmed, noget os aldeles Uvedkommende,
da maa i det Mindste dets Mulighed erkjendes, og denne Muligheds Erkjendelse være begrundet
i Selverkjendelse.

Begynde vi ud fra det Rationelle, da er en af to
Antagelser mulig: enten maa Naturens Orden have sin
oprindelige Grund i Naturen selv eller i en guddommelig
Villen. Antager man det Første, saa er der ingen
Aabenbaring, Aabenbaringslæren er Mythologie. Antager
% --- p. 38 ---
\setcounter{footnote}{0}%
\opage{38}man det Sidste, da maa det nærmere afgjøres, hvad
der skal forstaaes ved en guddommelig Villen. Man
kan da sætte den guddommelige Villen i en saa nødvendig
Eenhed med Lovene for de naturlige Phænomeners
Rækkefølge, at den Intet vil og Intet kan ville
uden alene det, som stemmer med vor fornuftige Indsigt
i det Heles Sammenhæng og kan indlemmes i vore
Erfaringers stedse mere udvidede Kredse: der kan da
fra dette rationalistiske Synspunkt heller ikke være
Tale om nogen særegen Gudsaabenbaring. 

Skal Muligheden af en Aabenbaring erkjendes, da
maa den guddommelige Villen udtrykkelig sættes i en
guddommelig Villie, en almægtig selvbestemmende Frihed,
og denne Friheds bevægende Grund i en uendelig Kjærlighed.
Denne Opfattelse maa give Afkald paa videnskabelig
Beviisførelse, men er derfor ikke fornuftstridig.
Ingen menneskelig Fornuft kan udmaale Dybderne i
Naturens Anlæg, thi ingen menneskelig Fornuft kan
udmaale Dybderne i den guddommelige Kjærlighed.
Naturaarsagerne beroe paa Naturanlæg, Naturanlægene
paa guddommelig Selvbestemmen; den virkelige Natur
er en Villiesaabenbarelse af Aanden. Løsreven fra
Aanden har Naturen ingen Selvstændighed. De naturlige
Tings naturlige Aarsager ere alle begrundede i det
Overnaturlige: Naturens materielle Virkelighed bestaaer
i at være Mulighed for den guddommelige Villies overnaturlige
Virken.

Seet fra dette Synspunkt kan Universet med sine
Love siges at hvile paa et Under. Den skabende Aand
har ikke sat sin Frihed til paa den naturnødvendige
Orden ei heller udtømt sit Væsens Indre i Universets
Ydre. Er den skabende, opholdende Villie en Kjær%
% --- p. 39 ---
\setcounter{footnote}{0}%
\opage{39}lighedsvillie, der vil aabenbare sig, da maa Muligheden
for en personlig Kjærlighedsaabenbaring fra Begyndelsen
af høre med i Naturens Anlæg. Aabenbaringsunderet
med sine Mirakler bliver fra denne Side at opfatte
som et Tegn paa Guds personlige Frihed: i Udslag
af Kjærlighedens Almagt træder det i Naturen skjulte
Overnaturlige frem af sin Skjulthed og lader sig tilsyne
i det Naturlige; der træder noget Nyt op i Naturen,
men Naturens Orden bliver ikke afbrudt.

Hvad er der nu til Hinder for at forlige Aabenbaringstroen
med en rationel Verdenserkjendelse? Blandt
Andet to Ting! For det Første den moderne Fixidee,
at Naturlove, Naturanlæg og naturlige Aarsager umuligt
kunne have Noget med det Overnaturlige at gjøre. At
denne Fixidee er ufornuftig, er viist i vort Foregaaende.
For det Andet: det Pludselige, Overraskende og just
derfor Usandsynlige i Aabenbaringsphænomenet, som det
er fremstillet for Autoritetstroen. Men at sætte det
Usandsynlige paa lige Trin med det Umulige, er ufornuftigt.
Der kan, hvad f. Ex. Geniet angaaer, ikke
drages skarpe Grændselinier mellem det, som ifølge
organisk Disposition skyldes Hjernevirksomheden, det,
som skyldes Selvvirksomheden, og det, som skyldes Beaandelsen.
Gaaer det saaledes med den geniale Inspiration,
hvor kan man da forlange at der ved Inspirationsmiraklet
skulde lade sig trække en kritisk Grændse
mellem det, som personlig meddeles af Guds Aand, og
det, som forud ligger i Menneskets Sjæl. Derimod er
det saavel hvad den mirakuløse som hvad den geniale
Inspiration angaaer, en uafviselig Fornuftfordring, at
Naturens og Aandens Love, om end paa høist forskjellig
Maade, maa finde deres Opfyldelse i begge.

% --- p. 40 ---
\setcounter{footnote}{0}%
\opage{40}Derved kastes et nyt Lys paa Aabenbaringsmiraklet
og dets Tilegnelse. I Oldtiden var Indtrykket af det
Overnaturlige stærkt, men Fornuftudviklingen forholdsviis
ringe; i vore Dage er den fornuftige Indsigt i Naturtingenes
Sammenhæng rig og klar, men Indtrykket
af det Overnaturlige fattigt, svagt og forsvindende. Skal
det for Nutidsbevistheden være muligt at tilveiebringe % PRINTED AS IS: Nutidsbevistheden (for Nutidsbevidstheden; read at the image)
psychologisk Ligevægt mellem det Religiøse og det
Rationelle, da kan det kun skee ved en Vexelvirkning
mellem Autoritetstroen og den personligt frie Tilegnelse.

Aabenbaringsordet med sine Vidundere og hellige
Kjendsgjerninger maa, trods den videnskabelige Kritiks
Opløsninger, blive staaende for den religiøse Bevidsthed,
uforandret, uantastet. Fornuften maa betragte de hellige
Livsbilleder som originale Malerier af gamle Mestere
og ikke tillade at Farverne smudskes og Trækkene
udviskes af rationalistiske Fingre. Videnskabelig Autoritet
have disse Billeder ikke; men Troesautoritet,
Aabenbaringsautoritet have de og ville de have indtil
Verdens Ende.

Ogsaa Troesautoriteten maa opløses; men, vel at
mærke, ved Troeskritik til Selvforstaaelse og personlig
Tilegnelse. Den Maade, hvorpaa Opløsningen skeer,
den Dybde, hvori Tilegnelsen foregaaer, er forskjellig i
Forhold til Individernes Sindsbeskaffenhed og det Dannelsestrin,
hvorpaa de befinde sig.

Hvad Sindsbeskaffenheden angaaer, findes der Sjæle
af ædel Natur og rig Begavelse, hvis religiøse Trang er
saa ringe, at de ikke have Sands for en personlig Aand,
men derimod en oprindelig levende Sands for det upersonlige
Almene, en saa genial Sympathie med Allivet
i Naturen, at vi i særlig Forstand maae kalde dem Na%
% --- p. 41 ---
\setcounter{footnote}{0}%
\opage{41}tursjæle. Den engelske Digter Shelley, Byrons Ven,
var en saadan Natursjæl. I sin berømte Ode til Vestenvinden
udaander han sit Væsens Trang og Længsel.
„O, løft mig som en Bølge, et Blad, en Sky! jeg falder
paa Livets Torne, jeg bløder. En Vægt af tunge Timer
har lænket og nedbøiet En, der var Dig altfor lig ---
utæmmelig og hurtigfodet og stolt. Gjør mig til din
Lyre, som Skoven er det; hvad skader det, at mit Løv
falder som dens! Dine larmende Harmoniers Magt vil
aftvinge os begge en dyb Efteraarstone, sød, skjøndt
sorgfuld. Vær Du, o frie Aand, min Aand; vær mig,
Du voldsomt stormende! Driv mine Tanker, naar jeg
er død, over Universet som visne Blade, at de kunne
kalde ny Væxt tillive, og spred, besværget ved dette
Digt, mine Ord blandt Menneskeheden som Aske og
Gnister fra en uudslukket Arne!“ --- Der er igjen Sjæle,
hos hvilke Sandsen for Inderlighed er saa stærk, Selvbevarelsesdriften
saa ideal, Erkjendelsestrangen saa dyb,
at de først finde Ro, naar de ere naaede til Alkjærlighedens
Vold og føle sig rodfæstede i Selvtankens evige
Grund. Disse aandige Lyssjæle ville ikke fare hen
med Stormen og løses op i Elementernes Nat; de
ville være Borgere i Aandens og Frihedens Rige. Det
er ikke Æolsharper, der tone for Vindenes Pust, det er
Selvhedsinstrumenter, stemte i Fryd og Sorg, for hellige
Harmonier, Villiesorganer for den selvbevidste Magt, der
gjør Stormene til sine Engle, Ildsluerne til sine Tjenere.

Hvilke af de to Slags Sjæle ere nu de ægte, oprindelige
Menneskesjæle? Det kan ikke afgjøres i rationel
Psychologie, det er en Tro, der maa prøves i
Livet.

Tage vi de forskjellige Udviklings- og Dannelses%
% --- p. 42 ---
\setcounter{footnote}{0}%
\opage{42}trin, hvorpaa Individerne indenfor Aabenbaringstroens
Kreds kunne befinde sig, i Betragtning, vil Opfattelsen
af det Overleverede ændre sig med hvert Trin, og den
religiøse Bevidsthed skifte Farve. Den barnlige Tro,
Almuetroen og den i Fornuft reflecterede Tro vise hen
til forskjellige Udviklingstrin, og dog maa Troen i Aand
og Væsen være den samme. For Opdragelsesvæsenet og
særlig for dem, der ere satte til at være Ungdommens
Lærere, er her alvorlige Vanskeligheder at overvinde,
psychologiske Knuder saa indviklede, at der maa Aandsdannelse
til, for at kunne løse dem.

Med den barnlige Tro synes det at gaae som af
sig selv. Barnets Sjæl er lige saa modtagelig for Indtryk
fra den oversandselige som fra den sandselige
Verden. Bibelordets vidunderlige Anskuelighed tiltaler
dets Følelse og Phantasie, og det Overnaturlige vækker
ingen Forargelse: Himlen er aaben, og Guds Engle
fare op og ned over Menneskenes Søn. Men alt som
Barnealderen svinder, og Forstanden udvikles, indtræder
Krisen. --- Lad os, for at gjøre det anskueligt, forestille
os et halvt voxent Barn og dets Religionslærer
læsende sammen i Bibelen: om Skabelse, Paradis og
Syndefald i det Gl. Testamente, om Kristi Fødsel hans
Lidelse, Død og Opstandelse i det N. T. Her maa naturligviis
forudsættes, at Læreren troer oprigtig paa det,
hvorom der læses, og at Barnet er opvakt og vel begavet.
Indtil den unge Lærling har naaet sit fjortende
Aar eller saa omtrent, gaaer Alting vel. Drengen har
Glæde deraf og forstaaer det saa godt Altsammen,
baade det om Skabelsen og det om vore første Forældre,
hvor skjønt der var i Edens Have, og hvor taknemmelige
de burde have været, men at Slangen for%
% --- p. 43 ---
\setcounter{footnote}{0}%
\opage{43}førte dem til Ulydighed; det er næsten som om han
selv havde oplevet det. Og nu Evangeliet: Barnet i
Krybben, Englenes Lovsang, Hyrdernes Tilbedelse! Den
barnlige Elev forstaaer det, elsker det og troer det saa
vist, som havde han seet det med egne Øine. Med
Sorgens Alvor i Hjerteangst og Bedrøvelse følger han
den lidende Frelser over Kedrons Bæk til Gethsemane,
til Kajaphas, til Pilati Domhuus; han oplever Korsfæstelsen,
han staaer blandt de fromme Kvinder under
Korset. Frelseren døer; det er i Barnets Sjæl saa høitideligt
stille som ved Alteret i en Kirke. Saa kommer
Opstandelsens Morgen. Mørket svinder, Vagten flygter,
Engelen, klædt i Lys og Lyn, vælter Stenen fra Graven.

Med alt dette foregaaer der lidt efter lidt en Forandring.
Den unge Lærling nyder ogsaa anden Underviisning.
Som Discipel i en lærd Skole eller en Realskole
tilegner han sig Videnskabens første Elementer og
gjør gode Fremskridt, thi han er begavet, har Fatteevne
og en naturlig skarp Forstand. Han gaaer nu i
sit femtende Aar. Indtrykkene fra den virkelige Omgivelse
begynde at samle sig for ham til et Verdenshele,
der er høist uligt det, han har modtaget fra
Aabenbaringslæren; Kulturtankerne trænge sig ind paa
ham; han har hemmelige Tvivl og veed ikke ret, hvorledes.
Imidlertid er han indadtil saa optagen af at
lære og huske og udadtil saa umiddelbart nærværende
blandt Kammeraterne, at han endnu ikke har faaet
Tid til at gruble over Sammenhængen; men i Religionstimen
har han hvert Øieblik en Indvending, der dukker
op med en uvilkaarlig Trang til at spørge. Drives Underviisningen
ganske officielt, saa skal han vel vogte
sig for at komme frem med sit Eget; de tillærte For%
% --- p. 44 ---
\setcounter{footnote}{0}%
\opage{44}klaringer ere mere end nok til at nedslaae alle frygtsomme
Indvendinger. Men har han Tillid til sin Religionslærer
og er saa heldig stillet, at han kan tale med
ham i al Fortrolighed, da ville Betænkeligheder, om
end i al Uklarhed, liste sig frem. Her er et Vendepunkt,
hvor det skal vise sig, om Læreren er sin Opgave
voxen.

Er Læreren en Barnetroende,  som, naar det kommer
til Stykket, ikke veed anden Udvei end at beraabe sig
paa Menigheden og Daabspagten og Hjertelaget og Aanden
i Barnetroen, da kan han ved sit Ords Varme og sine
hjertelige Formaninger vel gjøre Indtryk paa det unge
Menneske, maaskee endog for en Tid tilsyneladende dysse
 hans Sind til Ro; men Fornuftens Væsen gjennemskuer
han ikke, Videnskabens Grundsætninger fatter han ikke:
hvorledes skal han saa kunne hjælpe? Der er faldet
et Tvivlens Frø i den unge Sjæl, og Tvivlen skyder op
og grener sig, thi Tvivlen har Luftrødder ligesom
Mangletræet, der i tropiske Egne danner de ufremkommelige
Skove.

Er Læreren en høikirkelig orthodox Dogmatiker,
vil han tage Sagen paa en anden Maade. „Fornuften
er blind for det Hellige; Nutidens hedensk opblæste
Videnskaber ere alle paa Vildfarelsens Vei. Guds aabenbarede
Ord skal troes bogstavelig. Naar der udtrykkelig
staaer i den hellige Skrift, at Gud skabte Verden i
sex Dage, saa skabte han den i sex Dage, og hvad
betyde saa den ræsonerende Fornufts Indvendinger og
Naturforskernes Drillerier? Gud maa dog selv bedst
vide, hvorlænge han har været om at skabe Verden.
Guds Ord er ufeilbart som Guds Aand; var det gaaet
anderledes til med Skabelsen end skrevet staaer, saa
% --- p. 45 ---
\setcounter{footnote}{0}%
\opage{45}havde Guds Aand nok sagt os det reent ud; thi Gud
er sanddru, Guds hellige Aand kan ikke lyve. Dersom
Himlen ikke var en fast Hvælving, hvor kunde der da
paa saa mange Steder staae, at Himlen aabnede sig?
Hvor skulde Kristus vel være faret hen, da han foer
op til Himlen? Hvorledes kunde Paulus blive henrykket
i den tredie Himmel, dersom der ikke var tre Himle?“
--- Men hvad udrettes der ved slige Beviisgrunde vel
Andet end at det unge Menneske nu ret begynder at
føle sig som Kritiker. Læreren kan imponere ham ved
et kraftigt: „Ti stille og lær dine Lectier, Du er endnu
for grøn til at disputere om Sligt!“ Læreren kan lade
ham forstaae, at han har en grundig Theolog for sig,
og give ham Haab om, at han selv en Gang med Tiden
kan blive en grundig Theolog og slaae alle Physikens
„Børnelærdomme“ til Jorden; dog hvad hjælper det den
stakkels Lærling? Tvivlene krybe i Skjul, men de sidde
derinde.

Er Læreren en gemytlig Blanding af Kritiker og
Theolog, halvt en Rationalist og halvt en Orthodox,
vil han maaskee indvie Drengen i store Hemmeligheder.
Han vil med Beklagelse indrømme ham, at det ikke
staaer saa rigtig til med den hellige Skrift og vore
symbolske Bøger, som man kunde ønske det. Han vil
med Haanden paa Hjertet oprigtig tilstaae, at der saa
vel i det Nye som i det Gamle Testamente findes Mirakler,
hvorpaa man slet ikke kan troe, Miraklet med
Gadarenernes Sviin, Miraklet med Bileams Æsel o. m.
fl. Men derhos vil han da heller ikke forsømme at
lægge sin Lærling alvorlig paa Sinde, at der virkelig
ogsaa er Mirakler, hvorpaa man ubetinget bør troe, at
sige med Modifikation. Pintsemiraklet, f. Ex., er et
% --- p. 46 ---
\setcounter{footnote}{0}%
\opage{46}Hovedmirakel. Der er noget storartet Symbolsk deri;
det maa man absolut troe, for ikke at blive Fritænker.
For at styrke det unge Menneske paa een Gang baade
i Kritiken og Troen, vil den gemytlige Lærer vel ogsaa
være saa god at indvie ham saa smaat i den „kritiske
Indledning“ og i „Apologetiken“. Men hvad vil
Følgen hlive? At den vordende Mand, hvis han med % PRINTED AS IS: hlive (for blive; read at the image)
opvakt Sands bevarer sin skarpe Forstand ubeskaaren,
maa ende som „Fritænker“. Vil han lade sin Forstand
omskjære, det er en anden Sag, saa kan han maaskee
opnaae at blive Theolog.

Er Læreren en Mand af Aandsdannelse, en Personlighed,
som i Kulturbevidsthedens stride Strømme
har arbeidet sig frem, ikke til en forloren Barnetro,
hvormed man kryber tilbage i Ammestuen, men til en i
Sandhed barnlig Tro, en Tro, der er luttret i Selvforstaaelse,
en levende Tro paa det aabenbarede Ord, da vil
han vide at tage Sagen som den er. Han vil hverken
afvise det unge Menneskes Tvivl ei heller indlade sig
i Disput om Enkeltheder; han vil forstaae, at det igjen
er de to Verdensbetragtninger, Oldtidens og Nutidens,
Aabenbaringens og Fornuftens, som nu i denne Sjæl
begynde at brydes med hinanden. Han vil som en
kyndig Lods hjælpe den endnu Uerfarne over de første
farlige Grunde, vise ham Sømærkerne og orientere
ham paa det frie Hav. Det er en psychologisk Krise,
han har for sig, og han vil behandle den med den
modne Mands Overlegenhed. --- „Hvad er der i Veien,
min unge Ven, at Du ikke længere, som før, kan troe
paa Skabelsens Under, paa Paradiset og Syndefaldet,
men begynder endog at tvivle om Kristi Opstandelse
og Himmelfart? For et Par Aar siden kunde vi begge
% --- p. 47 ---
\setcounter{footnote}{0}%
\opage{47}troe derpaa. Jeg har i denne Henseende ikke forandret
mig. Jeg finder endnu Bibelens Ord saa talende, dens
Udtryk saa slaaende, dens Billeder saa friske, saa fulde
af Liv og Kraft, at jeg forynges derved og hver Gang
tilegner mig noget Nyt. Med Dig gik det ligesaa. Jeg
seer endnu din frembrydende Taare, da Du ved det
Sted i Lidelseshistorien, hvor Pilatus siger til Jøderne:
see \emph{hvilket Menneske}! gav dit Hjerte Luft og udbrød:
Var Han med Tornekronen ikke kommet med i
Kongerækkerne, hvor var den lange Historie om dem med
Guldkronerne da afskyelig! Det var et Øieblik, da Du
pludselig syntes at blive ældre end Dig selv i Foryngelse.
Men nu begynder Du at blive gammel; hvad er
der hændet Dig, lad mig høre!“

Det unge Menneske fatter Mod og rykker ud med
sine Bekjendelser. En Dag --- fortæller han --- havde
vi Physik i Klassen; Læreren var blevet varm af Beviserne
for Tyngdens Lov. Det er, sagde han, en Lov
saa nødvendig begrundet i Legemernes Natur, saa fornuftig
og evig, at ingen Magt hverken i Himlen eller
paa Jorden er i Stand til at ophæve den. Da faaer en
af mine Kammerater det Indfald at spørge, om da
Kristi Himmelfart ikke gjorde en Undtagelse? Læreren
taug, men med en \emph{Mine} og et himmelvendt Blik, saa
hele Klassen loe. Jeg loe med; skjøndt det var mig
som et Knivstik i Hjertet. Dette Knivstik gav min
Forstand en underlig Skarphed, det var, som om den
med Eet blev vred og hvas og vilde skjære min Tro
i Stykker. Jeg begyndte nu at tænke over, hvad
jeg havde læst og hørt om vort Planetsystem og vor
Planets Ubetydelighed i Universet, om Urverdenens
Dyr og de Millioner, Billioner Aar, der maa være for%
% --- p. 48 ---
\setcounter{footnote}{0}%
\opage{48}løbne, inden Mennesket fremkom paa Jorden. Jeg er
kun en Begynder, og meget deraf er mig dunkelt; men
jeg har en uhyggelig Anelse, en hemmelig Mistanke om,
at Darwins Lære om Menneskets Afstamning fra en
eller anden Abeform dog nok kunde være den rette.
Aberne have fire Hænder, Mennesket kun to, hvor sandsynligt
er det nu ikke, at de bageste Hænder i Tidernes
Løb ere blevne omdannede til Fødder. Den hellige
Skrift veed ingen Besked om Sligt. Forleden Dag
gjorde Læreren i Naturhistorie os opmærksomme paa, at
ingen Slange æder Støv, endskjøndt det Modsatte ---
tilføiede han med et sarkatisk Smiil --- rigtignok staaer % PRINTED AS IS: sarkatisk (for sarkastisk; read at the image)
i Bibelen. Jeg er undertiden tilmode, som maatte jeg
enten opgive Videnskaben eller blive Fritænker; det
Første kan jeg jo ikke og det Sidste vilde jeg saare
nødig. Hver Gang jeg kommer til at tænke derpaa,
føler jeg mig meget ulykkelig.

Finder Du Dig ulykkelig, svarer Læreren, saa er
Faren endda ikke saa stor. Jeg var kun bange for,
at Du var blevet Friskfyr og i Tro paa Opinionens
Autoritet allerede havde valgt den moderne Bevidstheds
moderne Viden, som man i en Modeboutik vælger
en moderne Hat eller en moderne Kjole. Du vil
dyrke Videnskaben og deri gjør Du Ret; thi Gud har
givet Dig Evner og Gaver. Jo dybere Du trænger ind
i Videns og Videnskabens Aand, desto større Agtelse
vil Du faae for saadanne Mænd som Kepler og Newton,
Laplace og Darwin; men Du vil ogsaa beundre den
Redelighed, hvormed slige Mænd altid skjelne mellem
det de vide, og det de ikke vide. Du vil, naar Du
kommer saa vidt, besinde Dig paa, at der for den
menneskelige Viden er sat en dobbelt Grændse, en re%
% --- p. 49 ---
\setcounter{footnote}{0}%
\opage{49}lativ, der bestandig kan overskrides, og en absolut, der
ikke kan overskrides. Hvis din Lærer i Physik en
Gang ved en offenlig Prøve forelagde Dig en Opgave
fra den høiere Mathematik, som var løst af Laplace,
og forlangte af Dig, at Du, uden særlige Forstudier,
alene ved Hjælp af dine elementære Skolekundskaber
skulde løse den, vilde Du saa ansee dette for et Beviis
paa, at Læreren selv er en stor Mathematiker, eller
vilde Du ikke meget mere ansee det for et vilkaarligt
og ubilligt Forlangende? Du vil dog ikke antage, at
den Lærer, der holder sig til de simplere Opgaver og
ved deres Løsning anvender en Fremstillingsmaade, som
de endnu ikke Viderekomne kunne forstaae, altid gjør
det af den Grund, at han ikke selv veed Mere; Du
vil være beskeden nok til at antage, at han vel kan
have en høiere Viden, uagtet han nedlader sig til din
Fatteevne. Anvend nu dette paa den hellige Skrift og
spørg Dig selv, om det kan ansees for et Beviis paa
den guddommelige Aands Uvidenhed i Psysik og Naturhistorie, % PRINTED AS IS: Psysik (for Physik; read at the image)
at den i sine Aabenbarelser om Verdens Oprindelse
og Menneskeslægtens Afstamning har nedladt
sig til Oldtidens Fatteevne og talt i et Sprog, som
Menneskene kunde forstaae. Din sunde Forstand vil
sige Dig, at hvad der i en Naturlære vilde være en utilgivelig
Feil, kan i en Religionsmeddelelse, hvor det
Naturlige dog altid er og bliver Symbol paa det Aandelige,
være et stort Fortrin. Du taler om at blive Fritænker.
Hvis det, at frigjøre Tanken for indgroede
Fordomme og forud fattede Meninger, er Kjendemærket,
saa er Han, der siger: „I skulle erkjende Sandheden,
og Sandheden skal gjøre Eder frie“, den største
Fritænker. Men skal Legen med at lade sin
% --- p. 50 ---
\setcounter{footnote}{0}%
\opage{50}Tanke gaae i Veiret, som en Dreng lader sin Drage gaae
i Veiret, og flagre for alle Vinde, betyde Tankens Frigjørelse,
da har Fritænkeren ligesaa lidt hjemme i Videnskaben
som i Religionen. Hvor løst hans Tanke
end flagrer, er den dog ikke fri, han har en Traad i
den, som Drengen i Dragen: den har hemmelig Ordre
til at hakke i det, han ikke synes om, og rive det ned,
der er hans Instinkter imod. Du vil i denne Mening,
haaber jeg, aldrig blive Fritænker. Elsker Du Friheden,
jeg mener Aandsfriheden, da kan Du umulig antage,
at den sande Frihed skulde bestaae i at gjøre
Aanden til Naturens Træl og lægge den i Nødvendighedens
Lænker. Aabenbaringens Gud er Frihedens
Gud; hvad troer Du, det bliver til med det Menneskes
Frihed, der fornegter Frihedens Gud og forguder Nødvendigheden?
Det seer storladent ud, at en Mand af
den frie Tanke tør gaae i Kreds med Elementerne og
frigjøre sig selv, næsten lige saa storartet, som at en
Dværg tør krybe ind i en Kæmpes Harnisk, og en
Tommeliden fraternisere med Titaner. Elsker Du Friheden,
da er jo Frihedens Gud personlig din Gud, en
Gud, der ikke er udenfor Dig, men i Dig som dit Væsens
Ophav, og denne Gud er Aand. Tro paa Frihedens,
Personlighedens og Kjærlighedens Aand, og Aanden skal
frigjøre Dig saa vel fra Vantroens Forhaanelse af det
Overnaturlige som fra Overtroens Ringeagt for det Naturlige.
Din Tvivl og Uro betyder, at Du alt har
betraadt Frigjørelsens Vei; lad mig paa Veien medgive
Dig Pauli Ord: „Da jeg var et Barn, talede jeg som
et Barn, tænkte jeg som et Barn, dømte jeg som et
Barn, men da jeg blev Mand, aflagde jeg det Barnagtige.“

Vi have dvælet saa længe ved denne Overgang, thi
% --- p. 51 ---
\setcounter{footnote}{0}%
\opage{51}for vort Opdragelsesvæsen er den i nærværende Tid
den allervigtigste. Kunde vi see de Millioner Sjæle,
som, uagtet de i Barndommen have modtaget religiøse
Indtryk, dog siden hen under Erfaringernes Mangfoldighed i den
naturlige Verden, have tabt det Oprindelige,
det første Enfoldige --- hvad enten de forøvrigt indbilde
sig at have bevaret „Barnetroen“ i Autoritetens Umiddelbarhed
eller de have valgt Emancipationen med „Kjødets
Evangelium“ eller i indifferentisk Sløvhed ingen af Delene
--- kunde vi see dem i et psychologisk Speil, mon det
da ikke skulde vise sig, at den begyndende Forkvakling
stammer fra den Overgangsperiode, hvori Barnet
modnes til Yngling og Ynglingen til Mand, netop i den
Periode, da Religionsunderviisningen afsluttes? At begynde
en Religionsunderviisnng er forholdsviis en let % PRINTED AS IS: Religionsunderviisnng (for Religionsunderviisning; read at the image)
Sag, --- det kan enhver Katechismuspræst og barnetroende
Asyl-Peter gjøre; men at bringe Opdragelsen
til en saadan Afslutning, at en religiøs-personlig Selverkjendelse
i Forening med en virkelig fornuftig Verdenserkjendelse
kan siges at være grundlagt, er en
vanskelig Opgave, hvis Løsning betinges af Aandsdannelse.


% --- p. 52 ---
\setcounter{footnote}{0}%
\lecture{Fjerde Forelæsning.}\opage{52}

Det var i November 1842, at Pseudonymen Victor
Eremitas \emph{Enten Eller} slog ned i vor Literatur som
et Lyn fra klar Himmel. Dette Lynslag kan med Sandhed
siges at betegne et Vendepunkt for Aandslivets
Udvikling i Danmark. Allerede Forordets første Linier
ere charakteristiske. „Det er maaskee dog stundom
faldet Dig ind, kjære Læser, at tvivle en Smule om
Rigtigheden af den bekjendte philosophiske Sætning, at
det Udvortes er det Indvortes, det Indvortes det Udvortes.
Du har maaskee selv gjemt en Hemmelighed,
om hvilken Du følte, at den, i sin Glæde eller i sin
Smerte, var Dig for kjær til at Du kunde indvie Andre
i den. Dit Liv har maaske bragt Dig i Berøring med
Mennesker, om hvilke Du ahnede, at noget Saadant var
Tilfælde, uden at dog Din Magt eller Din Besnærelse
var istand til at bringe det Forborgne til Aabenbarelse.
Maaskee passer intet af Tilfældene paa Dig og Dit Liv,
og dog er Du ikke ukjendt med hiin Tvivl; den har
nu og da som en flygtig Skikkelse svævet Din Tanke
forbi.“ --- Der er i denne Tiltale noget Aandeagtigt,
der spænder Forventningen. Eremiten med det fine,
kloge Ansigt, og de vidunderlig klare, gjennemtrængende
% --- p. 53 ---
\setcounter{footnote}{0}%
\opage{53}Øine synes livagtig at staae for En med en Gaade paa
Læben, en sælsom Opdagelse, og aabne et hemmeligt
Gjemme. Stemmen er fremmed, og dog fortrolig, ikke
høirøstet, men tiltrækkende, sympathetisk; den kjærtegner
Tanken, den fanger Opmærksomheden, den hvisker om,
at der inde i Sjælen er Noget, som kaldes \emph{Inderlighed}.

Aaret efter viste sig en anden Pseudonym, Johannes
de silentio, med \emph{Frygt og Bæven}. Han forudseer sin
Skjæbne i en Tid, da man „har slaaet en Streg over
Lidenskaben for at tjene Videnskaben“, „Kjærligheden“,
siger han, „har dog i Digterne sine Præster, og stundom
hører man en Røst, der veed at holde den i
Hævd; men om Troen høres der intet Ord, hvo taler
til denne Lidenskabs Ære? Philosophien gaaer videre.
Theologien sidder sminket ved Vinduet og beiler til
dens Gunst, falbyder sin Deilighed til Philosophien.“
Man mærker paa Pseudonymens Stemning, at Troens
Lidenskab har bragt hans Sind i lyrisk Kog. Aaret
efter fremtraadte en tredie Pseudonym, Johannes Climacus,
med det philosophiske Spørgsmaal: Hvorvidt kan
Sandheden læres? Her møder, siger han, den Vanskelighed,
som Sokrates i Meno gjør opmærksom paa som
paa en „stridslysten Sætning“, at et Menneske umuligt
kan søge, hvad han veed, og ligesaa umuligt kan søge,
hvad han ikke veed; thi hvad han veed, kan han ikke søge,
da han veed det, og hvad han ikke veed, kan han ikke
søge, thi han veed jo end ikke, hvad han skal søge.
--- Enhver vil forstaae, at der ved denne Henpegen
paa Kierkegaardske Pseudonymer sigtes til hele S. Kierkegaards
litterære Virksomhed, og at det, hvad Aandsdannelsen
angaaer, maa være hans eiendommelige Opfattelse
af Tro og Viden, Religion og Fornuft, Kri%
% --- p. 54 ---
\setcounter{footnote}{0}%
\opage{54}stendom og Civilisation, hvorpaa vi i denne Sammenhæng
lægge Vægt.

Den sidste Deel af \emph{Enten Eller} slutter, som
bekjendt, med „det Opbyggelige, der ligger i den Tanke,
at mod Gud have vi altid Uret.“ Vort Liv bringer 
os i et mangfoldigt Forhold til andre Mennesker.
Nogle af disse elske Ret og Retfærdighed, Andre ikke,
de gjøre Uret. „Din Sjæl er ikke forhærdet mod den
Lidelse, de saaledes paaføre Dig, men Du randsager
og prøver Dig selv, Du forvisser Dig om at Du har
Ret, og Du hviler rolig stærk i denne Overbeviisning;
hvormeget de end krænke mig, siger Du, denne Fred
skulle de dog ikke formaae at frarøve mig, at jeg veed,
jeg har Ret, og at jeg lider Uret.“ Vort Liv bringer
os i mangfoldigt Forhold til andre Mennesker, til Nogle
drages vi med inderligere Kjærlighed end til Andre.
„Hvis nu et saadant Menneske, der var Gjenstand for
din Kjærlighed, gjorde Uret mod Dig, ikke sandt, da
vilde det smerte Dig, Du vilde nøie prøve Alt, men da
vilde Du sige: jeg veed med mig selv at jeg har Ret,
denne Tanke skal berolige mig? O, hvis Du elskede
ham, da vilde den ikke berolige Dig, Du vilde efterforske
Alt,\dots{} Du vilde ønske, at Du maatte have Uret,
Du vilde søge, om Du ikke kunde finde Noget, der
kunde tale til hans Forsvar, og fandt Du ikke det, Du
vilde først finde Hvile i den Tanke, at Du havde Uret.“

Anderledes med vort Forhold til Gud, thi det er et
Uendelighedsforhold. Kun i et uendeligt Forhold til
Gud kan Tvivlen beroliges; kun i et uendeligt frit Forhold til Gud kan Menneskets Bekymring forvandles til
Glæde. „I et uendeligt Forhold til Gud er han, naar
 han erkjender, at Gud altid har Ret, i et uendeligt
% --- p. 55 ---
\setcounter{footnote}{0}%
\opage{55}frit Forhold, naar han erkjender, at han altid har Uret.
Saa er da Tvivlen standset; thi Tvivlens Bevægelse laae
netop deri, at han det ene Øieblik skulde have Ret, i
det andet Øieblik Uret, til en vis Grad have Ret, til
en vis Grad Uret, og dette skulde betegne hans Forhold
til Gud; men et saadant Forhold til Gud er intet
Forhold, og dette var Tvivlens Næring\dots{}. Og skulde
den Tanke, at mod Gud have vi altid Uret, ikke være
begeistrende, thi hvad udtrykker den Andet end at
Guds Kjærlighed altid er større end vor Kjærlighed?“
Den fylder et Menneskes Sjæl med Kraft til at handle,
den gjør ham brændende i Aanden, kun naar han
endelig beregner, udslukkes Aandens Ild. „Om Du
selv var den, der maatte negte Dig dit høieste Ønske,
Du er dog glad, Du siger ikke Gud har altid Ret, thi
deri er ingen Jubel; Du siger: mod Gud har jeg altid
Uret. Om det, der var Dit Ønske, var, hvad Andre
og Du selv i en vis Forstand maa kalde Din Pligt, om
Du ikke blot maatte fornegte Dit Ønske, men paa en
Maade svige Din Pligt, om Du ikke blot tabte Din
Glæde, men selv Æren, Du er dog glad; mod Gud, siger
Du, har jeg altid Uret. Om Du bankede, men der
blev ikke opladt, om Du ledte, men ikke fandt, om Du
arbeidede, men Intet fik, om Du plantede og vandede,
men ingen Velsignelse saae, om Himlen var lukket, og
Vidnesbyrdet udeblev, Du er dog glad i Din Gjerning,
om Straffen, som Fædrenes Brødre havde nedkaldt,
kom over Dig, Du er dog glad, thi mod Gud have vi
altid Uret.“ --- Det er Troens Under, Inspirationens
Tungemaal, Idealets Forkyndelse i Retning af Inderlighed.

I „Frygt og Bæven“ fremstilles Abraham som Troens
Fader; Pointen er Isaaks Offring. Den Digterbegavelse,
% --- p. 56 ---
\setcounter{footnote}{0}%
\opage{56}hvormed Fremstillingen er udstyret, vil altid være Gjenstand
for æsthetisk Beundring, og den dialektiske Skarphed,
hvormed „det Absurde“ i Troen er drevet til
høiere Vanvid, vil endnu i en lang Fremtid give Stof
til religiøse Overveielser og ethiske Discussioner. „De
rede tause i 3 Dage, paa den 4de Dags Morgen sagde
Abraham end ikke et Ord, men opløftede Øiet og saae
Morijabjerget i det Fjerne. Han lod Drengene blive
tilbage og gik ene med Isaak ved Haanden op til
Bjerget. Men Abraham sagde til sig selv: „jeg vil dog
ikke dølge for Isaak, hvorhen denne Gang fører ham.“
Han stod stille, han lagde sin Haand paa Isaaks Hoved
til en Velsignelse, og Isaak bøiede sig for at modtage
den. Og Abrahams Aasyn var Faderlighed, hans Blik
var mildt, hans Tale formanende. Men Isaak kunde
ikke forstaae ham, hans Sjæl kunde ikke opløftes; han
omfattede Abrahams Knæe, han bønfaldt ved hans Fod,
han bad for sit unge Liv, for sin skjønne Forhaabning,
han mindede om Glæden i Abrahams Huus, han mindede
om Sorgen og Eensomheden. Da reiste Abraham
Drengen op og gik med ham ved sin Haand, og hans
Ord vare fulde af Trøst og Formaning. Men Isaak
kunde ikke forstaae ham. Han besteg Morijabjerget,
men Isaak forstod ham ikke. Da vendte han sig bort
et Øieblik fra ham, men da Isaak anden Gang saae
Abrahams Aasyn, da var det forandret, hans Blik var
vildt, hans Skikkelse var Rædsel. Han greb Isaak i
Brystet, kastede ham til Jorden og sagde: „Dumme
Dreng, troer Du jeg er Din Fader? Jeg er en Afgudsdyrker.
Troer Du det er Guds Befaling? Nei det er
min Lyst.“ Da skjælvede Isaak og raabte i sin Angst:
„Gud i Himlen forbarm Dig over mig, Abrahams Gud
% --- p. 57 ---
\setcounter{footnote}{0}%
\opage{57}forbarm Dig over mig, har jeg ingen Fader paa Jorden,
saa vær Du min Fader.“ Men Abraham sagde sagte
til sig selv: „Herre i Himlene, jeg takker Dig; det er
dog bedre at han troer, at jeg er et Umenneske, end
at han skulde tabe Troen paa Dig!“ --- Det er en
Stemningens Mulighed i „Frygt og Bæven“.

Ja --- siger man maaskee --- dybe og mørke ere
Havets Grunde, men endnu dybere og mørkere ere de
Sindets Afgrunde, der aabne sig, naar Troen kogler i
uforstaaet Lidenskab, og Tungsind og Særhed og sælsom
Phantasie gjøre hinanden Selskab for at ledsage
Abraham til Morija Bjerg og stirre ham ind i Øiet, just
som han drager Kniven. Abraham staaer som en ærværdig
Skikkelse svøbt i Oldtidens Taager; de store
Træk aftegne sig klart i det hellige Sagn, saa længe
han sees paa Afstand; men vover man at nærme sig
Offerstedet og gjøre sig hans Handling nærværende, da
gaaer Forstanden tabt, Aabenbaringens Gaade, det
rædselsfulde Troens Absurdum forvirrer dens Tanker.
Denne Betænkelighed er imidlertid, hvad Troens Væsen
angaaer, idetmindste foreløbig hævet derved, at Forfatteren
har tegnet os et andet Forbillede, en Skikkelse
som greben ud af Nutiden, en Troesridder, paa een
Gang Heros og Spidsborger. „Her er han. Bekjendtskabet
stiftes, jeg bliver forestillet for ham. I det
Moment, jeg først tager ham paa mine Øine, kaster
jeg ham i samme Nu fra mig, gjør selv et Spring tilbage,
slaaer Hænderne sammen og siger halv høit:
„Herre Gud! er det Mennesket, er det virkelig ham,
han seer jo ud som en Rodemester.“ Imidlertid er det
dog ham. Jeg slutter mig lidt nærmere til ham, passer
paa den mindste Bevægelse, om der ikke skulde vise sig
% --- p. 58 ---
\setcounter{footnote}{0}%
\opage{58}en lille ueensartet Brøks-Telegraphering fra det Uendelige,
et Blik, en Mine, en Gestus, et Vemod, et Smiil,
der forraadte det Uendelige i sin Ueensartethed med
det Endelige. Nei! Jeg examinerer hans Skikkelse fra
Top til Taa, om der ikke skulde være en Revne, igjennem
hvilken det Uendelige tittede ud. Nei! Han er
heelt igjennem solid. Hans Fodfæste? er kraftigt, tilhører
ganske Endeligheden, ingen pyntet Borgermand,
der Søndag Eftermiddag gaaer ud paa Fresberg, træder
grundigere paa Jorden, han tilhører ganske Verden,
ingen Spidsborger kan tilhøre den mere. Intet er at
opdage af dette fremmede og fornemme Væsen, hvorpaa
man kjender Uendelighedens Ridder. Han glæder sig
ved Alt, deeltager i Alt, og hver Gang man seer ham
deeltage i det Enkelte, skeer det med en Vedholdenhed,
som betegner det jordiske Menneske, hvis Sjæl hænger
fast ved Sligt. Han passer sin Gjerning. Naar man da
seer ham, skulde man troe, han var en Skriverkarl, der
havde fortabt sin Sjæl i det italienske Bogholderi, saa
punktlig er han. Han ferierer om Søndagen. Han
gaaer i Kirke. Intet hemmeligt Blik eller noget Tegn
paa det Incommensurable forraader ham; hvis man
ikke kjendte ham, var det umuligt at udsondre ham
fra den øvrige Mængde; thi hans sunde kraftige Psalmesang
beviser i det Høieste, at han har et godt Bryst.
Om Eftermiddagen gaaer han i Skoven. Han fryder
sig over Alt, hvad han seer, over Menneskevrimlen, de
nye Omnibusser, Sundet, --- naar man møder ham paa
Strandveien, skulde man troe, han var en Kræmmersjæl,
der slog sig løs, netop saaledes glæder han sig;
thi han er ikke Digter, og jeg har forgjæves søgt at
aflure ham den poetiske Incommensurabilitet. Henad
% --- p. 59 ---
\setcounter{footnote}{0}%
\opage{59}Aften gaaer han hjem, hans Gang er ufortrøden som
et Postbuds. Underveis tænker han paa, at hans Kone
vist har en apparte lille Ret varm Mad til ham, naar
han kommer hjem, f. Ex. et stegt Lammehoved med
Grønt til. Hvis han mødte en Ligesindet, da kunde
han vedblive lige til Østerport at samtale med ham
om denne Ret med en Lidenskab, der vilde passe for
en Restaurateur. Tilfældigvis eier han ikke 4 Sk., og
dog troer han fuldt og fast, at hans Kone har hiin
lækkre Ret til ham. Har hun den, da skal det være
et misundelsesværdigt Syn for fornemme Folk, begeistrende
for Menigmand at see ham spise; thi hans
Appetit er stærkere end Esaus. Hans Kone har den
ikke --- besynderligt nok --- han er aldeles den Samme.
Paa Veien kommer han forbi en Byggeplads, han træffer
en anden Mand. De tale et Øieblik sammen, han opfører
i et Nu en Bygning, han disponerer over alle
Kræfter dertil. Den Fremmede forlader ham med den
Tanke, det var vist en Capitalist, medens min beundrede
Ridder tænker: ja kom det derpaa an, kunde jeg sagtens
faae det. Han ligger i et aabent Vindue og betragter
Pladsen, paa hvilken han boer, Alt hvad der
foregaaer, at en Rotte smutter under et Rendesteensbræt,
at Børnene lege, alt beskæftiger ham med en
Ro i Tilværelsen, som var han en Pige paa 16 Aar.
Og dog er han ikke Genie; thi Geniets Incommensurabilitet
har jeg forgjæves søgt at udspionere hos
ham. Han ryger sin Pibe i Aftenstunden; naar man
seer ham, skulde man sværge paa, det var Spekhøkeren
ligeoverfor, der vegeterer i Skumringen. Han lade 5 være % PRINTED AS IS: lade (for lader; read at the image, T027)
lige med en Sorgløshed, som var han en letsindig Døgenigt,
og dog kjøber han hvert Øieblik, han lever, den
% --- p. 60 ---
\setcounter{footnote}{0}%
\opage{60}beleilige Tid til den dyreste Priis, thi han gjør end
ikke det Mindste uden i Kraft af det Absurde.“

Spørgsmaalet om det Absurde bliver altsaa i allerhøieste
Forstand et Aandsspørgsmaal, ja et Salighedsspørgsmaal.
I „Afsluttende uvidenskabelig Efterskrift
til de philosophiske Smuler“ fremstiller Johannes Climacus
Problemet saaledes: „Individets evige Salighed
afgjøres i Tiden ved Forholdet til noget Historisk, der
ydermere er saaledes historisk, at i dets Sammensætning
det er medoptaget, hvilket ifølge sit Væsen ikke
kan blive historisk, og altsaa maa blive det i Kraft af
det Absurde.“ Hvad dette betyder bliver fra et exclusivt
kristeligt Standpunkt paa evangelisk historisk Grundlag
anskuelig paaviist af Anti-Climacus: „\emph{Indøvelse i
Christendom}.“ Det Absurde, der er Troen en Salighed,
er Vantroen en Forargelse. Kristus, Troens
Gjenstand, er Forargelsens Mulighed: „1) den Forargelses
Mulighed, der ikke forholder sig til Christus
(Gud-Mennesket), men til ham som et slet og ret
enkelt Menneske, der colliderer med det Bestaaende.
2) Den væsentlige Forargelses Mulighed i Retning af
Høiheden, at et enkelt Menneske taler eller handler
som var han Gud, siger sig selv at være Gud, altsaa i
Retning af Bestemmelsen Gud i Sammensætningen Gud-Mennesket.
3) Den væsentlige Forargelses Mulighed i
Retning af Ringheden, at Den, der udgiver sig for at
være Gud, viser sig at være det ringe, fattige, lidende,
tilsidst afmægtige Menneske.“

Hermed var den hidtil gjældende Maalestok anulleret
og Staven brudt over de mange lærde, halvhistoriske
og halvphilosophiske Forsøg paa at gjøre
Kristendommen begribelig for den menneskelige Fornuft.
% --- p. 61 ---
\setcounter{footnote}{0}%
\opage{61}Historiske Undersøgelser over Bibelens Ægthed, dogmatiske
Speculationer over Kirkens Lære o. desl. kunne
vel afgive en passende Underholdning for de Lærde
selv, men for Troen betyde de Intet. Troen er Forstanden
en Forargelse og Fornuften en Daarskab; den
Troende troer det Absurde: Troen er et Paradox.
Hvorledes skulle vi da komme til Troen? Ad Inderlighedens
Vei under Timelighedens Tryk, i Endelighedens
Trængsler. „Tænk Dig skjult i en simpel Indfatning,
et hemmeligt Gjemme, hvor den kostbareste
Skat er nedlagt; --- der er en Fjeder, paa hvilken der
skal trykkes, men Fjederen er skjult og Trykket maa
have en vis Kraft, saa at et tilfældigt Tryk ikke kan
være nok, saaledes er Troens Anlæg, Evighedens Haab,
skjult i Menneskets Indre, og Trængselen er Trykket.
Naar der trykkes paa den skjulte Fjeder, og Trykket
er stærkt nok, saa aabner Indholdet sig i al sin Herlighed.
Eller tænk Dig et krybende Dyr, der dog har
Vinger, som det kan bruge, naar det drives til det
Yderste, men til dagligt Brug finder det det ikke Umagen
værd at bruge dem: saaledes er Troens Anlæg, Evighedens
Haab i Menneskets Inderste: han har Vinger,
men han maa bringes til det Yderste for at opdage
dem eller for at faae dem og bruge dem.“

Ved at fremstille Troen som Paradox gav S. Kierkegaard
den ræsonnerende Forstand med dens selvvigtige
Tvivlen om Alt et Slag paa Munden. Med Paradoxet
til Forudsætning udviklede han i „Kjærlighedens Gjerninger“,
i „Sygdommen til Døden“ og i sine mange
opbyggelige Taler en fuldstændig Cyklus af religiøse
Kategorier, Høidelinier for ideale Kredse, Klangfigurer
af en tonende Evighed. --- Man skulde have meent, at en
% --- p. 62 ---
\setcounter{footnote}{0}%
\opage{62}Kristendomsforkynder af saa overlegen Begavelse maatte
have fundet Tilslutning hos samtlige Lærere i Kristendommen,
de videnskabelige saavel som de folkelige.
Men nei! Theologerne saae med hemmelig Gru
paa det sære Phænomen, og der paakom dem en Fornemmelse,
en uhyggelig Anelse, at her var Noget,
hvormed det ikke hængte rigtig sammen. „Hvad for
Noget? Det skulde da vel aldrig være galt fat med
Theologien? Skulde den maaskee kunne befrygte en
Ildebrand?“ Hvor vil Du hen? Theologien har Intet at
befrygte; dens Leer er ildfast, dens Borg er tryg. Udhusene
med Faar og Kvæg og andre Kreaturer ere desuden
assurerede. Nei, det er den stakkels Forfatter, som
med alt sit Vid, sin Kløgt og sin Dialektik er ved at
gaae fra Forstanden. Det er da Skade; man kunde
ellers have brugt ham --- ikke til Kirkehistorie, thi
dertil mangler han Lærdom og storartet Overblik, ikke
til Dogmatik, thi dertil mangler han speculativ Dybsindighed
og Sands for det Objective, men til Apologetik:
han kunde have vundet sig et Navn som kristelig
Apologet. Nu maa han opgives. --- Og saa lød da Parolen
fra Hovedkvarteret, at den hele vidtløftige Kierkegaardske
Litteratur var indtil videre at behandle som
en Bunke „Strøtanker, Aphorismer, Indfald og Glimt“.

Man skulde have ventet, at de Geistlige, Ordets
Forkyndere i Menigheden, maatte have benyttet en
Haandsrækning, hvorved de med Eet kunde blive det
rationalistiske Uvæsen kvit og faae Bibelens Autoritet
ubeskaaren til vederhæftig Afbenyttelse. Men nei! Der
skal være Maade med Alt. Skal Bibelordet ret være
opbyggeligt, maa det holdes i en vis Forbigangenhedens
Skumring. Det skader ikke, at der er lidt Uvished i
% --- p. 63 ---
\setcounter{footnote}{0}%
\opage{63}Baggrunden, lidt Angst for Fritænkeri i Sindene; Ordets
Forkynder har da den Fordeel at kunne styrke Tilhørerne
i Troen og tale et myndigt advarende Ord mod
„disse Tiders Vantro“. For S. Kierkegaard vare almene
Betragtninger over „hine Tider“ og „disse Tider“
ørkesløse Sinkerier. Enhver, som vil, kan i aandelig
Forstand blive samtidig med Kristus paa samme Maade
som Apostlene. Thi om Kristus faaer man, objectivt
seet, Intet at vide. Kristus er ikke Gjenstand hverken
for from Beundring eller kritiserende Nysgjerrighed.
Kristus er Modsigelsens Tegn, Forargelsens Mulighed,
Troens Eenbaarne. Kristendommen er ikke en Lære,
den er en „Existensmeddelelse“. Vil Du existere som
Kristen eller vil Du ikke? det er Spørgsmaalet. Vil
Du ikke, kan Du jo gaae din Vei, vil Du --- see her
er Forbilledet! Saa holdt han det Nye Testamente saa
nær op for Øinene af Præst og Menighed, at han ved
sin „Samtidighed“ virkelig blev noget nærgaaende. Man
svarede ham i Kirkens Navn, at hvad Existensen angaaer,
var der Intet i Veien. „Vi have den samme
Herre, den samme Tro, den samme Daab, den samme
Helligaand som Apostlene. Vel sandt, vi kunne ikke
gjøre saa kraftige Gjerninger som Apostlene, ei heller
leve saa økonomisk som Apostlene, thi Tiderne have
forandret sig, men vi stole trygt paa Herrens Ord: Du
er Petrus, og paa denne Klippe vil jeg bygge min Menighed,
og Helvedes Porte skulle ikke faae Overhaand
over den.“ --- Nu var Fæstningen cerneret og Bombardementet
begyndte.

„Tænk Dig, at der var en mægtig Aand, som havde
tilsagt nogle Mennesker sin Beskyttelse, men paa den
Betingelse, at de skulde indfinde sig et bestemt Sted,
% --- p. 64 ---
\setcounter{footnote}{0}%
\opage{64}hvor der var Fare med at gaae hen --- sæt nu, at disse
Mennesker undlod at indfinde sig paa det bestemte Sted,
men gik hjemme i deres Dagligstue og talte begeistrede
Ord med hinanden om, hvorledes Aanden havde lovet
dem sin mægtige Beskyttelse, saa Ingen skulde formaae
at skade dem: er dette ikke latterligt?\dots{} At Helvedes
Porte ikke skulle faae Magt over hans Kirke,
disse Ord af Christus ere i den seneste Tid bragte
atter og atter i Erindring mod mig, mod min Paastand,
at Christendommen slet ikke er til. Mit Svar er dette.
Hiin Forjættelse hjælper ikke os det allermindste; thi
det Vrøvl vi leve i som var det at være Christen, er
slet ikke hvad Christus og det Nye Testamente forstaaer
ved at være Christen. Christenhed er slet ikke
Christi Kirke; ei heller siger jeg, at Helvedes Porte
have faaet Magt over Christi Kirke, paa ingen Maade.
Nei, jeg siger, at „Christenhed“ er et Vrøvleri, der har
klamret sig fast ved Christendommen som et Spindelvæv
paa en Frugt, og som nu er saa god at ville forvexle
sig med Christendom.“

Vi stride og stræbe, vi hige og søge. Hvad er
det vi søge? „Først Guds Rige“. (Et Slags Novelle.)
Cand. theol. Ludvig From --- han søger. Og naar man
hører, at „en theologisk“ Candidat søger, behøver man
da ikke nogen levende Indbildningskraft for at forstaae,
hvad det er han søger, naturligviis Guds Rige, som
man jo skal søge først. --- Nei, det er det dog ikke;
det han søger, er: et kongeligt Levebrød som Præst;
og er der, hvad jeg med nogle faa Træk skal angive,
først skeet saare Meget, førend han er naaet saa vidt.
Først har han gaaet i den lærde Skole, fra hvilken han
saa er dimitteret. Derpaa har han først taget 2 Exa%
% --- p. 65 ---
\setcounter{footnote}{0}%
\opage{65}miner og, efter 4 Aars Læsning først taget Embedsexamen.
Han er da altsaa theologisk Candidat; og
man skulde maaskee mene, at han, efter først at have
tilbagelagt alt Dette, endelig kan komme til at virke
for Christendommen. Ja, paa det Lav. Nei, først
maa han gaae \textonehalf{} Aar i Seminariet; og naar det er
gaaet, kan der da ikke være Tale om at kunne søge i
de første 8 Aar, hvilke altsaa først maae være tilbagelagte.
Og nu staae vi ved Novellens Begyndelse: de
8 Aar ere forløbne, han søger. Hans Liv, som hidtil
ikke kan siges at have havt noget Forhold til det Ubetingede,
antager pludseligt et saadant Forhold: han
søger ubetinget Alt; skriver det ene Ark stemplet Papir
fuldt efter det andet, render fra Herodes til Pilatus,
anbefaler sig baade hos Ministeren og hos Portneren,
kort han er ganske i et Ubetingets Tjeneste. Ja, en
af hans Bekjendte, der i de sidste Par Aar ikke har
seet ham, mener til sin Forbauselse at opdage, at han
er blevet mindre, hvilket maaskee lader sig forklare af,
at det er gaaet ham som Münchhausens Hund, der var
en Mønde, men ved den megen Løben blev en Grævling.
Saaledes gaaer der 3 Aar. Vor theologiske Candidat
trænger virkelig til Hvile, til efter en saa uhyre
anstrænget Virksomhed at blive sat ud af Virksomhed
eller at komme til Ro i et Embede og pleies lidt af
hans tilkommende Kone, --- thi han er imidlertid først
blevet forlovet. Endeligen --- som Pernille siger til
Magdelone --- slaaer hans „Forløsnings“ Time, saa han
med Overbeviisningens hele Magt vil, af egen Erfaring,
kunne „vidne“ for Menigheden, at i Christendommen
er der Frelse og Forløsning: han faaer et Embede.
Hvad skeer? Ved at indhente en endnu nøiere Efter%
% --- p. 66 ---
\setcounter{footnote}{0}%
\opage{66}retning end han hidtil har havt om Kaldets Indtægter,
opdager han, at disse ere c. 150 Rd. mindre end han
havde troet. Nu er 101 ude. Det ulykkelige Menneske
næsten fortvivler. Han har allerede kjøbt stemplet
Papir for at indgaae med en Ansøgning til Ministeren,
at det maa tillades ham at betragtes som ikke kaldet
--- og for saa igjen at begynde forfra: da en af hans
Bekjendte faaer ham til at opgive dette. Det bliver
altsaa derved, han beholder Kaldet. Han er ordineret
--- og Søndagen kommet, da han skal fremstilles for
Menigheden. Provsten, ved hvem dette skeer, er en
mere end almindelig Mand, han har ikke blot (hvad
de fleste Præster have, og oftest i samme Grad mere,
jo høiere de komme op i Rangen) et uhildet Blik for
den jordiske Profit, men tillige et speculativt Blik paa
Verdenshistorien, Noget, han ikke beholder for sig selv,
men lader Menigheden komme til Gode. Han har,
genialt, valgt til Text de Ord af Apostelen Peter: „see
vi have forladt alle Ting og fulgt Dig“, og forklarer
nu Menigheden, at just i Tider som vore maa der
saadanne Mænd til som Lærere, og anbefaler i Forbindelse
hermed denne unge Mand, om hvem Provsten
veed, hvor nær han var traadt tilbage for de 150 Rd.’s
Skyld. --- Den unge Mand bestiger nu selv Prædikestolen
og Dagens Evangelium er (besynderligt nok!):
tragter \textbf{først} efter Guds Rige. Han holder sin Prædiken.
„En meget god Prædiken, siger Biskoppen, som selv
var tilstede, „en meget god Prædiken; og det frembragte
ordentlig Virkning det hele Parti om „først“
Guds Rige, den Maade, paa hvilken han fremhævede
dette „Først“. „Men synes da Deres Høiærværdighed,
at her var den ønskelige Overeensstemmelse mellem
% --- p. 67 ---
\setcounter{footnote}{0}%
\opage{67}Talen og Livet; paa mig gjorde det næsten et satirisk
Indtryk dette „Først“. „Hvilken Urimelighed; han er
jo kaldet til at forkynde Læren, den sunde, uforfalskede
Lære om at søge først Guds Rige; og det gjorde han
meget godt.“

Ved at see bort fra Livet, og tale om Læren, den
sunde, uforfalskede Lære, faaer man „nymodens religiøse
Betryggelser“. „Engang i en længst, længst forsvunden
Tid forstod man Sagen saaledes: man fordrede
af Den, der vilde være Lærer i Christendom, at hans
Liv afgav Garantien for, hvad han lærte. Dette er
man nu længst kommet bort fra, Verden er bleven
klogere, alvorligere, har lært at ringeagte alt dette
Smaalige og Sygelige med det Personlige, lært at begjære
ene det Objective --- nu fordrer man, at Lærerens
Liv afgiver Garantien for, at det, han siger,
er Commerce, dramatisk Festlighed, Divertissement,
reent objectivt. Nogle Exempler. Er det dette Du vil
tale om, at Christendommen, det nye Testamentes
Christendom, har en Forkjærlighed for eenlig Stand ---
og Du selv er et eenligt Menneske: Kjære, det er ikke
Noget for Dig at tale om, Menigheden kunde jo troe,
at det var Alvor, blive urolig eller den kunde føle sig
fornærmet, at Du saaledes, saa upassende blander din
Personlighed ind. Nei, det har lange Udsigter, før Du
kan blive duelig til med Alvor at tale derom saaledes,
at Du i Sandhed tilfredsstiller Menigheden. Vent til
Du engang allerede har din første Kone i Jorden og
er et godt Stykke henne med den anden: da er Tiden
for Dig, da træder Du frem for Menigheden, lærer og
„vidner“, at Christendommen har Forkjærlighed for
eenlig Stand --- og Du vil ganske tilfredsstille;  thi dit
% --- p. 68 ---
\setcounter{footnote}{0}%
\opage{68}Liv afgiver Garantien for, at det er Fjas, Commerce,
eller at det, Du siger, er: interessant. --- Er det dette
Christelige, Du vil tale om, at Christendommen lærer
Foragt for Titler og Ordener og alle Ærens Narrestreger
--- og Du selv hverken er Rangsperson, eller
Noget, som ligner Sligt: Kjære, det er ikke Noget for
Dig at tale om, Menigheden kunde jo troe, at det var
Alvor eller føle sig fornærmet over den Mangel paa
Dannelse, saaledes at paanøde sin Personlighed. Nei,
bie til Du først selv har lagt Dig en Slump Ordener
til, jo flere jo bedre, bie til Du selv slæber en Ramse
af Titler med Dig, saa Du for Mængden af Titler
neppe selv veed, hvad Du hedder; da er det Tiden,
da træder Du frem, prædiker og „vidner“ --- og Du
vil utvivlsomt tilfredsstille Menigheden; thi dit Liv afgiver
Garantien for, at det er dramatisk Forlystelse,
en interessant Formiddags-Underholdning. Er det dette
Du vil tale om, i Armod at forkynde Christendom, at
det er den sande christelige Forkyndelse --- og Du selv
bogstavelig er en fattig Diævel: Kjære, det er ikke % PRINTED AS IS: Diævel (for Djævel; read at the image, T038)
Noget for Dig at tale om, Menigheden kunde jo troe,
at det var Alvor, blive angst og bange, føle sig aldeles
ude af Stemning og i allerhøieste Grad uhyggelig berørt
ved at Armoden kom En saa nær paa Livet. Nei,
skaf Dig først et fedt Levebrød, og naar Du saa har
hav det saa længe, at Du snart vil avancere til et % PRINTED AS IS: har hav (for har havt; read at the image, T039)
endnu federe: da er den beleilige Tid kommen, da
træder Du frem for Menigheden, prædiker og „vidner“
--- og Du vil ganske tilfredsstille; thi dit Liv afgiver
Garantien for, at de Hele løber ud paa en Spas, som % PRINTED AS IS: de Hele (for det Hele; read at the image, T040)
alvorlige Mænd kunne ønske den, i Theatret eller i
Kirken, en Forfriskelse for at samle nye Kræfter til ---
% --- p. 69 ---
\setcounter{footnote}{0}%
\opage{69}at tjene Penge. --- Og paa den Maade dyrker man Gud
i Kirkerne! Og saa græder disse Fløiels- og Silke-Talere,
Stemmen svigter qvalt i Taarer!“

Der berettes i Evangeliet, at en Engel til visse
Tider foer ned i Bethesda Dam og rørte Vandene, saa
at Den iblandt de mange Syge, som derpaa først kom
i Badet, blev helbredet. Her er ogsaa En, der rører
Vandene, Troens, Tankens og Sprogets Vande, men er
han en Engel? Han synes snarere at være en Dæmon;
thi Enhver som helst iblandt „de Sunde“, der vover et
Spring i de rørte Vande, maa jo blive lam baade paa
Sjæl og Legeme. Hvad er han da, hvad vil han?
„Hvad jeg vil? Ganske simpelt: jeg vil Redelighed.
Jeg er ikke, som man velmenende --- thi Forbittrelsens
og Raseriets og Afmagtens og Vrøvlets Opfattelse af
mig kan jeg ikke tage Hensyn til --- har villet fremstille
mig, jeg er ikke en christelig Strenghed lige over
for en given christelig Mildhed. Paa ingen Maade; jeg
er hverken Mildhed eller Strenghed --- jeg er en menneskelig
Redelighed.“

Det høieste Aands- og Kulturspørgsmaal er sat i
Bevægelse. Hvem har nu Ret, S. Kierkegaard eller
hans Modstandere? Spørgsmaalet er dybtgaaende og
dertil saa skærpet og tilspidset, at det sætter sin Spydsod
mod Tidens Bryst og tvinger den til at staae foran et
Enten-Eller. Man kan vende Øinene bort fra Spørgsmaalet
og lade som om det var glemt eller slet ikke
til. Men det er lige fuldt til, og lader sig ikke afvise.
Naar man læser de samlede Bladartikler, Kierkegaards
Angreb og Geistlighedens Forsvar, og tager „Øieblikkene“
med, danner der sig ud af Forhandlingerne et
eiendommeligt Billede, en Situation, der minder om
% --- p. 70 ---
\setcounter{footnote}{0}%
\opage{70}Helten og Choret i den antike Tragoedie. Dette Skuespil
i vor danske Literatur vil sikkert gaae over i
fremmede Literaturer og blive opført paa det store
Verdenstheater. Om hundrede Aar skrider Paradoxets
Elsker med sine Kjærlighedsgjerninger og sine „1000
Præster“ over Scenen, dramatisk og med fuld Musik
som Don Juan med sine 3000 Elskerinder.

% --- p. 71 ---
\setcounter{footnote}{0}%
\lecture{Femte Forelæsning.}\opage{71}

\emph{Det Nye Testamentes Kristendom er ikke
mere til}. Det er S. Kierkegaards Paastand, hans „eneste
Thesis.“ Og Beviset? er kort og fyndigt. Kristendommen
er „Existensmeddelelse“. Den religiøse Existens i Nutiden
og den religiøse Existens i Apostlenes Tid ere hinanden
saa diametralt modsatte, at den ene kan betragtes
som en Parodie paa den anden. „Hvis det vi
forstaae ved at være Christen virkelig er det at være
Christen: hvad er saa Gud i Himlene? Han er det
latterligste Væsen, der nogensinde har levet, hans Ord
den latterligste Bog, som nogensinde er kommet for
Lyset: at sætte (som han jo gjør i sit Ord) Himmel
og Jord i Bevægelse, at true med Helvede, med evige
Straffe --- for at opnaae det, vi forstaae ved at være
Christne (og vi ere jo sande Christne): nei noget saa
Latterligt er aldrig forekommet! Tænk Dig, at der til
et Menneske traadte en Mand med skarpladt Pistol og
sagde til ham: jeg skyder Dig ned“, eller tænk Dig det
endnu forfærdeligere, at han sagde: jeg bemægtiger
mig Din Person og martrer Dig ihjel ved den rædsomste
Dødsmaade, hvis Du ikke (pas nu paa, nu kommer det!)
hvis Du ikke gjør Dig Livet her paa Jorden saa pro%
% --- p. 72 ---
\setcounter{footnote}{0}%
\opage{72}fitabelt og nydelsesrigt som Dig er muligt, „saa er dette
dog vel den latterligste Tale\dots{}. Den rædsomste Art
Gudsbespottelse er den „Christenhed“ forskylder: at
forvandle Aandens Gud til --- et latterligt Vrøvl; og
den aandløseste Art Gudsdyrkelse, aandløsere end Alt,
hvad der findes og fandtes i Hedenskabet, aandløsere
end det, som Gud at tilbede en Steen, en Oxe, et Insekt,
aandløsere end Alt er: at tilbede under Navn af
Gud --- et Vrøvlehoved!“

\emph{Naar hørte det N. Testamentes Kristendom
op med at være til}? Efterhaanden som man hørte op
med at forsage Verden, offre Livet for Troen og lide
for Læren. Den har altsaa været til endnu i Joh. Husses
Tid, Luther maa vel ogsaa regnes med til det Nye Testamentes
Kristne? Ja, paa en Maade; men tage vi
det strengt, saa er Christendommen egentligen ikke
kommet ind i Verden. Det blev ved Forbilledet og
høiest Apostlene; men disse forkyndte den dog allerede
saa stærkt i Retning af Udbredelse, at her allerede
Misligheden begynder\dots{}. Kun guddommelig Myndighed
kunde imponere Menneskeslægten saaledes, at det blev
til ubetinget Alvor med ubetinget at ville det Evige.
Kun Gudmennesket vilde kunne udholde, hvis man vil
tænke sig det, 1000 Aar igjennem og atter 1000 Aar
ubetinget at virke for Lærens Udbredelse ved at forkynde
den, selv om han ikke fik en eneste Discipel,
hvis han kun kunde faae dem ved at forandre Vilkaaret.
Apostelen har dog nogen selvisk Trang til den Lindring
at faae Tilhængere, hvad Gudmennesket ikke har, der
ikke selvisk trænger til Tilhængere, derfor kun har
Evighedens Priis, ingen Markeds-Priis.“

\emph{Og hvis er saa Skylden, at det N. Testamentes}
% --- p. 73 ---
\setcounter{footnote}{0}%
\opage{73}\emph{Kristendom ikke mere er til}? Naturligviis Menneskenes
Skyld. „I det nye Testamente er, ifølge Christi egen
Lære, det at være Christen, blot menneskelig talt, idel
Qval, en Qval sammenlignet med hvilken alle menneskelige
Lidelser næsten kun er Barnestreger. Det
Christus taler om --- thi han lægger ikke Dølgsmaal
derpaa --- er om at korsfæste Kjødet, hade sig selv,\dots{}
om de mest hjerteskjærende Lidelser ved at hade Fader,
Moder, Hustru, sit eget Barn“\dots{} Det var slet ikke
efter Menneskenes Sind. Saa greb man „til Hyklerie,
man forfalskede Bestemmelsen af det at være Christen.
Det at være Christen, heed det, er idel Salighed,
„hvad var jeg, o, hvad var jeg, hvis jeg ikke var Christen;
o, uskatteerlige Gode at være Christen; jo, det at
være Christen, det giver først ret Livet Betydning,
Glæderne Smag og Lidelserne Lindring.“

Her er S. Kierkegaards \emph{aabenbare Selvmodsigelse}.

Først bestemmer han Kristendommen som en Existensmeddelelse,
en overnaturlig Aandsmeddelelse, en
Gave af Guds frie Naade, og præcicerer Udtrykket % PRINTED AS IS: præcicerer (for præciserer; read at the image)
derhen, at Kristus, Gud-Mennesket selv, det eneste
sande Aandsmenneske, er den eneste sande Kristen.
Apostlene, hans nærmeste Efterfølgere, forraade allerede,
trods Inspirationen, paa selve Pintsefesten et betænkeligt
Hang til Forfalskning, en „selvisk Trang til
den Lindring, at faae mange Tilhængere.“ Og hvorfor?
Fordi Aanden i Apostlene ikke var saa stærk som
Aanden i Kristus. Saa ubetinget er Existensmeddelelsen
afhængig af den overnaturlige, paradox-mirakuløse Aandsmeddelelse,
at den yderste Anstrengelse af menneskelig
Kraft og Villie i dette Forhold bliver en forsvindende
% --- p. 74 ---
\setcounter{footnote}{0}%
\opage{74}Størrelse. Et Menneske kan og bør tilegne sig det
Guddommelige, for saa vidt det bliver ham, den Enkelte,
personlig meddeelt, men intet Menneske kan med
al sin Iver lægge en Alen til sin Væxt eller tilegne sig
end det Mindste af hvad Guds er, naar Gud ikke selv
vil meddele ham det. Under denne Forudsætning maatte
da S. Kierkegaard med sin „menneskelige Redelighed“
være kommet til følgende Resultat: Det N. Testamentes
Kristendom er ikke mere til. Ak hvor beklageligt.
Men lad os ikke jamre derover eller gjøre os selv
ufrugtbare Bebreidelser; thi det er ikke Menneskenes
Skyld, men Guds urandsagelige Villie. Gud gjorde i
Tidens Fylde et Experiment med en overnaturlig Existensmeddelelse,
men lod det kort efter falde igjen.
Hvorledes? Af hvilke Grunde? Blev hans Kjærlighed
kold, blev Aanden træt, gjorde Menneskene ham det
saa broget, at han blev kjed af det Hele? Spørg ikke,
spot ikke! Det er jo, hvad det skal være, --- et Paradox,
et Absurdum. Med denne Forklaring kunde Paradoxets
Ordfører kjækt have hævdet sin Thesis og
conseqvent blevet færdig med hvad han kalder „Præstevrøvlet.“


Men den Vei gik han ikke. Han tog et andet
Tilløb. Han stillede Forbilledet frem, for at indskærpe
Fordringen paa ubetinget Efterfølgelse. Enhver, som
alvorlig vil, maa kunne troe, at Gud-Mennesket har
existeret. I Troen paa denne Existens er Existensmeddelelsen
med dens overnaturlige Betingelser een
Gang for alle given. Hver Troende, som vil, kan blive
samtidig med Kristus, kan træde i hans Spor og følge
ham efter. „Forholdet er: jo mere Du indlader Dig med
Gud og jo mere han elsker Dig, jo mere vil Du blive,
% --- p. 75 ---
\setcounter{footnote}{0}%
\opage{75}menneskeligt talt, ulykkelig for dette Liv, jo mere vil
Du komme til at lide i dette Liv.“ I den Grad er nu
den overnaturlige, paradox-mirakuløse Existensmeddelelse
gjort afhængig af Menneskets Sind og Villie, at
her end ikke er Skygge af Forandring fra Guds Side,
thi Gud er Kjærlighed og indlader sig i Øieblikket
med hver den, der vil indlade sig med ham. Den opstillede
Thesis bliver da at forsvare saaledes: Det N.
Testamentes Kristendom er ikke mere til; men det er
 ikke Guds Skyld, det er Menneskenes Skyld. „Dette
veed jeg paa to Maader af Erfaring. Deels deraf, at
jeg ikke kan holde den Tanke“ (Kjærlighedstanken)
„ud og derfor blot opdagende øiner denne sande christelige
Bestemmelse af det at være Christen, medens
jeg for mit Vedkommende hjælper mig til at udholde
Lidelser med en langt lettere Tanke, en jødisk, ikke
en i høieste Forstand christelig: at jeg lider for mine
Synders Skyld; deels deraf, at jeg ved mit eget Livs
Forhold maatte føres til at blive opmærksom herpaa“\dots{}.
Den „menneskelige Redelighed“ er nu denne:
Det N. Testamentes Kristendom er ikke mere til; thi
jeg, Søren Aaby Kierkegaard, har ikke villet, at den
skulde være til. Havde jeg villet indlade mig dybere
med den guddommelige Kjærlighed og ladet den gjøre
mig absolut ulykkelig, ganske og aldeles ulykkelig i
dette Liv, saa havde der dog i det mindste været een
Kristen til: Kristendommen havde da været til i mig.
Men nu har jeg ikke villet det. Hvorfor ikke? Spot
ikke, spørg ikke derom; det er jo, hvad det skal være,
det Paradoxe, det Absurde. Min Opgave har jeg løst:
at rive Masken af Hykleriet. Præsterne sige, at de
ville, hvad Kristus vil, og for at holde Menneskene hen
% --- p. 76 ---
\setcounter{footnote}{0}%
\opage{76}i denne Indbildning gjøre de andægtige Miner og forfalske
Kristendommen; jeg angiver sandt hvad Kristendom
er, og saa salver jeg mit Hoved og toer mit Ansigt
og siger frit ud: Kristendommen er slet ikke til; jeg
er endnu ikke Kristen, fordi jeg endnu ikke vil være
det. Det er min Redelighed overfor Præsterne.

Man seer, at den sidste Forklaring i Frækhed og
Vanvid ikke staaer tilbage for den første. I en saadan
Frækhed har S. Kierkegaard visselig ikke gjort sig
skyldig. Men hvorved har han undgaaet Vanvidets
Yderste? Ved at bøje af fra Conseqvensen, sætte en
Dup paa Spydsodden og tilstaae sin Skrøbelighed: jeg
kan ikke bære den store Kristustanke, Forbilledet er
for stort og jeg er for lille; med andre Ord: jeg venter
endnu paa en Existensmeddelelse.

Det Kierkegaardske Gudsforhold er himmelvidt forskjelligt
fra det nytestamentlige Gudsforhold. Det
nytestamentlige Gudsforhold er et umiddelbart Aabenbaringsforhold;
det Kierkegaardske er et uendelig dialektisk
Reflexionsforhold. Modsigelsen, den radikale
Selvmodsigelse, som S. Kierkegaard med al sin Reflecteren
aldrig fik op i Bevidstheden, er: at et afledet
Reflexionsforhold til Gud skulde kunne give samme
Existens som det primitive Aabenbaringsforhold. Visselig
har S. Kierkegaard tænkt og handlet, levet og lidt
trods Nogen i Samtiden. Hans nærmeste Eftertid har
endnu ikke ret besindet sig paa, hvad det er for et
Aandsarbeide han har faaet udført og hvor uendelig
Meget, den skylder ham. Men hans Reflexion var negativ,
Paradoxet, det Absurde var hans eneste Positive. Paradoxets
Positive --- det er dets dybe Hemmelighed ---
er selv i sit Inderste negativt. Den salige Ruus i
% --- p. 77 ---
\setcounter{footnote}{0}%
\opage{77}Lidelse, Forsagelse, Opoffrelse er en negativ Existens;
Guds Kjærlighed løser sig op i lutter negative Bestemmelser
og forvandler sig tilsidst til et sjælefortærende
Afgrundsdyb, der er forfærdeligere end Døden. „Gud
er Din Dødsfiende; han, Kjærligheden, han vil af Kjærlighed
til Dig være elsket af Dig: dette betyder, at
Du maa døe, afdøe, ellers kan Du ikke elske ham. Saa
sidder han da, allestedsnærværende og alvidende som
han er, og passer paa Dig, veed end det Mindste der
foregaaer i Dig --- det veed han, Din Dødsfjende!
Vogt Dig for at ønske Noget, vogt Dig for at frygte
Noget; thi hvad Du ønsker, opfyldes ikke, men snarere
det Modsatte, og hvad Du frygter, og jo mere, jo snarere
kommer det over Dig. Thi han elsker Dig, og han vil
være elsket af Dig; begge Dele af Kjærlighed; men
saasnart der er Noget, Du ønsker, saa tænker Du ikke
paa ham, ligesom naar der er Noget Du frygter, eller
hvis Du bringer ham i Forbindelse med Din Ønsken
og Frygten, saa tænker Du jo altsaa ikke paa ham i
og for sig selv, det er, Du elsker ham ikke --- og han
vil være elsket, vil det af Kjærlighed.“ --- Med en
saadan Kjærlighed havner man ikke i Himlen hos Kristus,
men i Nirvana hos Buddha.

Som den Sandhedsaabenbaring, der aabner Livets
hellige Kilder, kan Kristendommen med Grund kaldes
en Existensmeddelelse, for saa vidt den er en uendelig
frugtbar Livsmeddelelse. Men Liv og Existens ere dog
ikke enstydige Ord. Det samme Liv kan give sig Udtryk
i forskjellige Existenser: det evige Liv er de
timelige Existensers uudtømmelige Grund. Det evige
Liv i Kristendommens Væsen bliver ikke staaende ved
en enkelt Tidsalders stereotype Existensschema, det
% --- p. 78 ---
\setcounter{footnote}{0}%
\opage{78}gaaer, uforanderligt, som det er, ind i Forandringer,
evigt, som det er, ind i Tid og Omskiftelse, overnaturligt,
som det er, ind i Naturbetingelser og naturlige
Forhold, saligt, som det er, ind i jordisk Forkrænkelighed,
i verdslige Brydningers Uro, Lidelser og Strid.
Herved bliver Forholdet mellem Meddelelsen og Tilegnelsen
nærmere at bestemme. Her er en fortsat
Vexelvirkning mellem Meddelelse og Tilegnelse, men
Vexelvirkningens Factorer forholde sig ikke altid paa
samme Maade.

Først er den guddommelige Factor almægtig overgribende.
Ordet bliver Kjød, Gudsaabenbareren træder
frem som Menneske, lærer og lider, døer og opstaaer
til Vitterliggjørelse af Livets personlige Seier over
Døden. Han indgaaer til Herlighed; men just som
han har sat sig paa Kongestolen ved Faderens Høire,
kommer Aanden som i et fremfarende vældigt Veir med
Tunger af Ild ned over Apostlene, og Ilden tænder i
disse Vidner, og Tungemaalene gløde i Vidnesbyrdets
Flammer. Det var den store Gjennembrudstid; Himlen
aabnedes og Gravene aabnedes: det skjulte Under kom
i denne Aandsvirksomhed til Udslag i aabenbare Mirakler.


Saaledes er den primitive Existensmeddelelse, Kristendomsexistensen
i det N. Testamente. At maale
Nutidsexistensen med denne Maalestok og gjøre „Christenheden“
ansvarlig for, at den ved egen Brøde har
ladet den nytestamentlige Existens gaae tabt, forraader
en uhyre Misforstaaelse. Apostelexistensen hørte op,
da Gjennembrudsperioden var afsluttet.

Der følger nu Perioder, i hvilke den guddommelige
Meddelelse drager sig tilbage i Form af det skjulte
% --- p. 79 ---
\setcounter{footnote}{0}%
\opage{79}Under, for at Tilegnelsens menneskelige Factor særlig
kan træde frem og komme til bevidst Virksomhed, ---
den katholske og den gammelprotestantiske Periode.

See vi bort fra Forskjellighederne i det Udvortes
og fæste Blikket alene paa Tilegnelsens Væsen, da har
den katholske Kirke ved sin Gjerningshellighed unegtelig
gjort Alt for at anstrenge den menneskelige Villie,
men den har ogsaa gjort Alt for at vedligeholde
Illusionen om en i mirakuløse Gjennembrud sig fortsættende
Aandsmeddelelse. --- Vexelvirkningen mellem
den guddommelige og den menneskelige Factor har
derved antaget et Præg af mechanisk Endelighed.
Den protestantiske Kirke har trængt Forestillingen om
mirakuløse Gjennembrud tilbage og reduceret Aandsmeddelelsen
til et stille perennerende Under. Den
har vel ikke villet slappe Tilegnelsens Kraft, men dog
givet Villien en eensidig indadgaaende og selvopløsende
Retning i Dogmet om Retfærdiggjørelsen ved Troen.

Begge Kirker ere, hver paa sin Maade, i lige stor
Forlegenhed med den virkelige Menneskenatur og den
sig fremarbeidende Existens. Og hvorfor? Fordi de,
berøvede den sprudlende Aandens Kilde, som overnaturligt
nærede den apostoliske Existens, ikke desto mindre
i Efterligningens Form have fastholdt det apostoliske Existensschema.
Vi skulle efter begges Anvisning, forsage Verden,
døde Kjødet, kue Naturdrifterne, kort sagt, betragte
det hele nærværende Liv som en Jammerlighed, der kun
duer til at plage os op til Salighed hisset. Hvad den humane
Existens angaaer, er Katholicismen med sin gjennemførte
Asketik, sit ufordulgte Fiendskab mod fri
fornuftig Oplysning og sit Munkevæsens mørke Foragt
for Livets naturlige Formaal imidlertid mere conseqvent
% --- p. 80 ---
\setcounter{footnote}{0}%
\opage{80}end Protestantismen. Den protestantiske Kirke har
med Bibeholdelse af sit apostoliske Schema villet knæsætte
Humaniteten og tage Kulturskjønheden i Favn;
det er en Tvetydighed, den dyrt har maattet betale.
Hvor er nu den nytestamentlige Existens, hvor er Kjødets
Dødelse, hvor er Guds Rige? „Cand. theol. Ludvig
From han søger“. --- S. Kierkegaards Satire gjør
en følelig Virkning, hvert Slag af hans Svøbe efterlader
blodige Strimer, netop fordi han gjør det nytestamentlige,
af hans Modstandere selv høitidelig anerkjendte,
Existensschema gjældende.

Nutidsexistensen har eiendommelige Krav, som den
hverken kan eller vil opgive; den gjør Naturlivets Betydning
gjældende fra Grunden af og fordrer Fornuften
anerkjendt. Dersom Kristendommen virkelig var et
Paradox, og Troen et Absurdum, saa var det Nye Testamentes
Kristendom rigtignok ikke mere til; men, vel at
mærke, af den Grund, at den aldrig har været til uden
i Menneskenes Indbildning, en Indbildning, der kun
kan holde sig saalænge, indtil man har opdaget det
Paradoxe.

Hvad er Paradoxet? Paradoxet med det Absurde
kommer tilsyne, naar to Begreber, der i
fuld Gjennemsigtighed vise sig indbyrdes modsigende,
sættes sammen i Eenhed. En fiirkantet Cirkel er et
Absurdum. Siger man nu: i Aandens og Sandhedens
høieste Regioner ere alle Cirkler fiirkantede, og dette
skal troes imod al Forstand, saa protesterer Fornuften;
vi have da Paradoxet. Siger man: i Gud-Menneskets
Existens er der Noget blevet historisk, som ifølge sit
Væsen umulig kan være historisk, og altsaa maa være
blevet det i Kraft af det Absurde, og dette skal
% --- p. 81 ---
\setcounter{footnote}{0}%
\opage{81}troes mod al Forstand, saa protesterer Fornuften:
vi have Paradoxet. At Modsigelsen i den fiirkantede
Cirkel og den foregivne Modsigelse i Gud-Menneskets
Væsen imidlertid ere absolut ueensartede Modsigelser
kan Enhver, som vil bruge sin Fornuft, indsee. Cirkelens
Begreb og Fiirkantens Begreb kunne tydeliggjøres
i exacte Bestemmelser og tilsvarende Anskuelsesformer.
Men Gudsbegrebets Indhold kunne vi ikke udtømme
i exacte Bestemmelser og Menneskebegrebets
heller ikke. Naar da Troen fastholder begge Begrebers
Eenhed i Gud-Menneskets Væsen, fatter Fornuften, at
det, som herved forudsættes, ikke er et Absurdum, men
et Mysterium.

Maaskee vil „Cand. theol. From“ anmærke, at det
netop er, hvad han fra Theologien længe har vidst, og
just derfor søger\dots{}. Lad ham søge! Mysteriet er
ikke en gammel theologisk Nat, hvori alle Katte ere
graae, er ikke det ubestemmelige Tusmørke, hvori Religionen
røres sammen med Videnskaben, Troen med Fornuften
til opbyggelige Forklaringer for uklare Hoveder.
Mysteriet er et Grændsebegreb, hvori Tro og Viden
mødes og skilles. Det maa ved omhyggelig Videnskritik
nøiagtig bestemmes, hvor dybt Fornuften kan
naae med Erkjendelsen af det Naturlige; den absolute
Grændse, hvori Erkjendelsen løser sig op, er Mysteriet.
Det maa ved omhyggelig Troeskritik nøiagtig bestemmes,
hvor dybt den religiøse Bevidsthed kan gaae med
Erkjendelsen af det Overnaturlige, uden at komme i
Modsigelse med Fornuften og derved i Strid med sig
selv; den absolute Grændse, hvori den opløses, er Mysteriet.


% --- p. 82 ---
\setcounter{footnote}{0}%
\opage{82}Aabenbaringsmiraklet er et Mysterium, Gud-Menneskets
Væsen et Mysterium. Fæster man nu eensidig
Opmærksomheden paa de mirakuløse Udslag og tilsvarende
Existenskriser, idet man vilkaarlig undgaaer
at tænke paa de i Mysteriet skjulte Mellembestemmelser,
saa faaer man ganske rigtig ud, at det Hele er et
Absurdum. Lægger man med S. Kierkegaard et saadant
Existensschema som Maalestok an paa Nutidsexistensen
og fordrer Commensurabilitet, saa faaer man
ganske rigtig ud, at Tro og Kristendom, Anerkjendelse
af det Absurde, det Paradoxe, ikke mere er til. Men
Aabenbaringsmysteriet er endnu til, og vil fremdeles
være til.

„Idealerne skulle forkyndes!“ Rigtigt. Men Spørgsmaalet
er, hvorledes de skulle realiseres og forme sig
til Existens. Abraham staaer endnu som et Troens
Forbillede; men ikke ligefrem ved sit Existensschema.
„Gud fristede Abraham og sagde“\dots{}. Gud talede altsaa
til Abraham, som en Herre kan tale til sin Tjener.
Dette Schema kunne vi ikke tilegne os. Paa vort Existenstrin
vil Enhver, som beraaber sig paa guddommelige
Bud af umiddelbar Indskydelse, blive anseet for
sindssvag. Derimod kan det, at et Menneske efter
samvittighedsfuld Selvprøvelse og moden Overveielse af
religiøse Grunde finder sig for Gud forpligtet, absolut
forpligtet til at offre endog det, der er ham det Kjæreste
og Dyrebareste i Verden, ikke ansees for Sindssvaghed.
„Gud sagde til Abraham: tag Isaak, din eneste Søn,
som du elsker, og gaae hen\dots{} og offer ham til et Brændoffer“\dots{}
Dersom dette Schema bogstavelig skal gjælde
for Nutiden, saa bliver de Troendes Gud til en Molok, og
det høieste existentielle Udtryk for Tro og Lydighed
% --- p. 83 ---
\setcounter{footnote}{0}%
\opage{83}vilde da være: at offre sine Børn til Molok. Derimod
kan en personlig Aandsfordring, om den end i sin indre
Ubetingethed maa komme i stærk Collision med Samtidens
Forestillinger, vel tænkes at være begrundet og
fuldstændig berettiget. Den, der i Troskab mod Aandsvillien
med Gud alene som Medvider skal gjøre Fordringen
fyldest, vil erfare, hvad Lydighed er, hvad
det er at troe og bære et Ansvar. Naar han da, under
det tunge Tryk af det tunge Ansvar, opløfter sit Hoved
og seer Abraham staaende paa Morija Bjerg med Offerkniven
i Haanden, da tegner Idealet sig paa Mysteriets
dunkle Grund; han seer et lysende, opløftende Exempel
paa Tro og Lydighed.

„Idealerne skulle forkyndes!“ Ja. Men dog vel
ikke derved, at de redigeres in absurdum. Den humoristiske
Figur, S. Kierkegaards Pseudonym har dannet
ved Sammensmeltningen af Troesridder og Spidsborger,
er efter Tegningen kjærnesund; men ifølge det Absurdes
Mechanik bevæger den sig paa Skruer. „Han (Troesridderen)
tømmer i den uendelige Resignation Tilværelsens
dybe Vemod, han kjender Uendelighedens Salighed,
han har fornummet Smerten af at forsage Alt,
det Kjæreste, man har i Verden, og dog smager Endeligheden
ham fuldt saa godt som den, der aldrig kjendte
noget Høiere, thi hans Forbliven i Endeligheden har
intet Spor af en forknyt beængstet Dressur, og dog har
han denne Tryghed til at fryde sig ved den, som var
den det Visseste af Alt. Og dog, dog er hele den
jordiske Skikkelse, han frembringer, en ny Skabning i
Kraft af det Absurde. Han resignerede uendeligt paa
Alt, og da greb han Alt igjen i Kraft af det Absurde.
Han gjør bestandig Uendelighedens Bevægelse,  men han
% --- p. 84 ---
\setcounter{footnote}{0}%
\opage{84}gjør det med en saadan Correcthed og Sikkerhed, at
han bestandig faaer Endeligheden ud, og der er intet
Secund, hvor man aner noget Andet.“

Dette Vidunder belyses nu særligt ved Hjælp af
en Ungersvend, der lægger hele sit Livs Indhold ind i
Kjærligheden til en Prindsesse, uagtet Forholdet er et
saadant, at denne Kjærlighed „umuligt lader sig realisere,
umuligt lader sig oversætte fra Idealiteten til
Realiteten.“ For nøiagtig at angive, hvori det Absurde
ligger, stilles Ridderen af den uendelige Resignation
imod Troens Ridder. Naar Uendelighedens Ridder har
indseet Umuligheden af at faae Prindsessen, resignerer
han paa Ønskets Opfyldelse. Hans Kjærlighed antager
en religiøs Karakteer og forklarer sig i en Kjærlighed
til det evige Væsen, der vel negter ham Opfyldelsen,
men dog forsoner ham i den evige Bevidsthed om, at
han besidder den Elskede i en Erindring, som ingen
Virkelighed kan fratage ham. Han gjør det Umulige
muligt derved, at han udtrykker det aandeligt, og
aandeligt udtrykker han det ved at give Afkald paa
det. Ønsket, der vilde føre ham ud i Virkeligheden,
men strandede paa Umuligheden, bøies nu indefter,
men er derfor ikke tabt, heller ikke glemt. Fra det
Punkt, kvor Uendelighedens Ridder standser, gjør Troens % PRINTED AS IS: kvor (for hvor; read at the image, T049)
Ridder endnu en Bevægelse, forunderligere end Alt, thi
han siger: jeg troer dog, at jeg faaer hende, i Kraft
nemlig af det Absurde, i Kraft af, at for Gud er Alting
muligt. Saaledes er Troens Ridder „den eneste Lykkelige,
Endelighedens Stamherre“, medens Resignationens
Ridder er en „Fremmed og Udlænding“.

Det Grundpunkt, hvorom Bevægelserne dreie sig,
er Grændsepunktet mellem det Mulige og det Umulige.
% --- p. 85 ---
\setcounter{footnote}{0}%
\opage{85}Enhver af de to Elskere er enig med sig selv om, at
det, menneskelig talt, er ham umuligt, ikke blot usandsynligt,
i allerhøieste Grad usandsynligt, men absolut
umuligt at faae Prindsessen. Lad os nu see nærmere
til, hvad heri ligger. For at den menneskelige Forstand
kan være fuldkomment sikker paa Umuligheden,
maa den have et ubedrageligt Tegn. Hvad skulde det
vel være for et Tegn? --- At den kongelige Fader,
om det saa var i fuld Samstemning med sit Statsraad
og hele Diplomatiet, det europæiske, asiatiske,
amerikanske, australiske tilsammen, satte sig imod
Partiet, vilde ikke være det ubedragelige Tegn.
Kongen kunde jo døe, og de diplomatiske Positioner
forandre sig. Selv om Prindsessens Hjerte var
koldt mod den Elskendes Bønner, det vilde dog ikke
være det ubedragelige Tegn. Det var jo ikke umuligt,
at hendes Følelser kunde forandre sig; man har Exempler
paa, at unge Piger have opgivet Knibskheden og omsider
ladet sig røre. Men --- vi sætte dette Tilfælde
--- dersom Prindsessen døde: ja saa have vi Tegnet,
thi det er tilvisse, menneskelig talt, umuligt at ægte hende,
naar hun er død og begraven. Hvad gjør nu Resignationens
Ridder? Han resignerer uendelig paa sit Ønskes Opfyldelse,
men bevarer sin Kjærlighed i Erindringens
Smerte. Men hvad gjør Troens Ridder? Han faaer
den --- døde og begravne --- Prindsesse, og lever med
hende i fuld Endelighed, i udvortes Virkelighed, i Fryd
og Herlighed, i Kraft af det Absurde. Skal dette forstaaes
bogstaveligt, saa maa man unegtelig give Pseudonymen
Ret, naar han udbryder: „det Sidste forbauser
mig, min Hjerne vender sig i mit Hoved ved at tænke
derpaa. Resignationens Bevægelse gjør jeg ved mig
selv, og det jeg derfor vinder, er min evige Bevidsthed
% --- p. 86 ---
\setcounter{footnote}{0}%
\opage{86}i salig Forstaaelse med min Kjærlighed til Gud. Men ved
Troen giver jeg ikke Afkald paa Noget, tvertimod, ved
Troen faaer jeg Alt netop i den Forstand, i hvilken
det hedder, at den, som har Tro som et Sennopskorn, % PRINTED AS IS: Sennopskorn (for Sennepskorn; read at the image, T050)
kan flytte Bjerge.“

Enten Fornuft eller Tro! Hvor Troen begynder,
hører Fornuften op, d. e. saalænge et Menneske endnu
har beholdt en lille Rest af sund Fornuft tilbage, er
den uendelige Resignation hans Høieste, han naaer
kun til det Punkt, hvor Troen skulde begynde. Den
Troende har lagt hele Fornuftstrækningen bagved sig og
gaaer videre. For ham er Grændsen mellem Sandsynligt
og Usandsynligt, Muligt og Umuligt ikke mere til,
thi for Gud ere jo alle Ting mulige. Altsaa enten
Fornuft, og da uden al Tro, eller Tro, og da uden al
Fornuft; dette er det Kierkegaardske Dilemma. Men
see vi nærmere til, mon det da ikke skulde vise sig,
at Tro og Fornuft, skjøndt i menneskelig Forstand to
absolut ueensartede Erkjendelsesfaktorer, dog have
deres Eenhedsgrund i det guddommelige Væsens Dyb,
i Mysteriet?

Pseudonymen i „Frygt og Bæven“ forsikkrer, at
han kan gjøre den uendelige Resignations Bevægelse
ved sig selv --- uden Tro --- i „salig Forstaaelse“ med
sin Kjærlighed til Gud. Men en salig Forstaaelse med
Gud er umulig, naar den Tanke, at for Gud ere alle
Ting mulige, er udelukket. Det hedder i „Sygdommen
til Døden“: Gud er det, at Alt er muligt. Ved at
fornegte denne Tanke ophæver man Gudsbegrebet og
fornegter Gud; men en salig Forstaaelse i Kjærlighed
til en fornegtet Gud er et Absurdum. Skal dette Absurde
absolut være Troens Mærke, saa befinder Ridde%
% --- p. 87 ---
\setcounter{footnote}{0}%
\opage{87}ren af den uendelige Resignation sig jo allerede paa sin
Viis indenfor Troens Omraade. Det Absurde ved Troesridderen
skal være dette, at han midt i Endeligheden veed
at finde sig til Rette med Forskjellen mellem det, i
menneskelig Forstand, Mulige og det Umulige, mellem
det Sandsynlige og det Usandsynlige, uagtet han
dog holder fast ved den Tanke, at for Gud ere alle
Ting mulige. Troesridderen er altsaa paa sin Viis en
ganske fornuftig Mand. Men dersom Resignationens
Ridder netop ved at være fornuftig er troende og Troens
Ridder netop ved at være troende er fornuftig, saa er
jo Grændsen hævet.

Er da den principielle Modsætning mellem Tro og
Fornuft forsvundet? Ingenlunde! Her begyndes ud
fra to absolut ueensartede Udgangspunkter: den naturnødvendige
Tingenes Sammenhæng og den guddommelige
Villies Almagt; men alt som man stiger ned i Dybet
fra den ene Side og op af Dybet fra den anden Side,
viser det sig at den sidste og yderste Grændse mellem
Muligt og Umuligt forsvinder i Mysteriet. Fornuftidealet
er: at bringe alle Virkelighedens Muligheder
ind under Altilværelsens betingede Nødvendighed, Troesidealet
er at bringe Nødvendigheden med den uendelige
Betingelsernes Kjæde ind under den guddommelige Villies
ubetingede Frihed; men ingen Dødelig skal kunne
paavise den sidste og yderste Grændse, hvor den betingede
Nødvendighed hører op, og den absolut ubetingede
Frihed begynder. Frihedens og Nødvendighedens
oprindelige, actuelle Eenhed er skjult i Mysteriet.

I Fremstillingen af Troesidealerne tale de Kierkegaardske
Pseudonymer med forskjellige Tunger og
modsige hverandre paa det Kraftigste, netop i Kraft af
% --- p. 88 ---
\setcounter{footnote}{0}%
\opage{88}det Absurde. Hos Johannes de silentio gjør Troens
Absurdum den Troende til Forjættelsernes Patriarch,
Velsignelsernes Arving, Endelighedens Stamherre, det
lyksaligste Menneske paa Jorden: Abraham faaer
Isaak i Kraft af det Absurde; Troesridderen faaer
Prindsessen --- den for Umulighedens Skyld døde og
begravne Prindsesse ---, han faaer hende virkelig, de
leve tilsammen i Fryd og Gammen, i Kraft af det Absurde.
Men hos Joh. Climacus og hos Anti-Climacus gaaer
det anderledes til. Hos Climacus gjør Troen den reflecterende
Humorist, der saa gjerne vilde forstaae den
og sætte sig ind i den, de forskrækkeligste Existensbryderier,
i Kraft af det Absurde, og hos Anti-Climacus,
den religiøse Karakteer, der er og lever midt i Troen,
gjør Troen Mennesket, menneskelig talt, til den ulykkeligste
Skabning, netop i Kraft af det Absurde. Men
naar det Absurde er lige saa kraftigt til at gjøre et
Menneske lykkeligt som til at gjøre det ulykkeligt, hvorledes
kan da Forskjellen mellem det Lykkelige og det
Ulykkelige blive Troens Kraftmaaler?

Lidelsernes, Smerternes og Forsagelsernes Existensschema
kan i Principet være lige saa vildledende som
Glædernes og Nydelsernes. Kristus er ikke kommet til
Verden for at forstærke Lidelserne eller for at forhøie
Glæderne; at han ikke desmindre gjør begge Dele, betyder,
at han i absolut Forstand er kommet, for Livets
og Lysets og Sandhedens Skyld, til Aabenbaring af
Guds hemmelige Raad, fra Frihedens hellige Dyb, Retfærdighedens
og Kjærlighedens Mysterium.

Noget ganske Andet er det, naar der tales om det
opdragende Element. Fra Opdragelsens og Aandsudviklingens
Standpunkt har S. Kierkegaard umiskjende%
% --- p. 89 ---
\setcounter{footnote}{0}%
\opage{89}ligt grebet det Rette, for saa vidt han, i sin Forkyndelse
af Idealerne, idelig og idelig med en psychologisk
Iagttagers indtrængende Skarphed og en smerterig Selvoplevelses
overbevisende Kraft har indskjærpet den
frivillige Lydighed under Aandens Lov og prædiket
„Lidelsernes Evangelium“. Men Nutidsexistensen og
dens aandelige Vilkaar har han ikke forstaaet. Han
meente at imponere det Bestaaende ved at drage de
nytestamentelige Existensfordringer frem og med Farver
af glødende Aand indbrænde hvert Træk af det
store Gjennembruds Schema i sin Samtids Bevidsthed,
og mærkede ikke, at Samtiden, hvor forhutlet den end
var ved sine „nymodens religiøse Betryggelser“, dog i
et meget væsentligt Punkt havde Ret imod ham. Han
meente at hidføre en aandelig Revolution, en afgjørende
Katastrophe. Endnu paa sit Dødsleie skal han have
sagt: „jeg falder som en Bombe, Bomben springer, og
saa tænder det.“ Men Røret gik forbi, og Katastrophen
udeblev, Bomben sprang, men det tændte ikke. Derimod
opnaaede han unegtelig paa en ret slaaende og
iøinefaldende Maade at overraske og fange det hele
officielle Kirkevæsen i Konsekvenserne af dets egne
Forudsætninger.

Den Enkelte skal ikke have Ret til at haane Menigheden
og sætte sin subjective Reflexion i Aandsmeddelelsens
Sted; men lad det kun være sagt reent
ud, der er i den officielle Kristendomsforkyndelse et
Noget, som fra Grunden af maa forandres. Det er ikke
det aabenbarede Ord, der skal forandres, det er Tilegnelsen,
den frie, personlige Tilegnelse af den indholdsrige
Sandhed, der maa fremstilles i et nyt Lys. Skal
Kristustanken fæste Rod i Sjælene og blive Kulturtidens
% --- p. 90 ---
\setcounter{footnote}{0}%
\opage{90}dybeste, kraftigst bevægende Livstanke, da ville Barnetroens
Billeder og Almuetroens Forestillinger ikke
strække til. Det er en virkelig Forstaaelse af det
Religiøses Forenelighed med det Rationelle, hvorpaa
det fremfor Alt kommer an. Existensens Omdannelse
er ikke katastrophisk; den foregaaer --- under forudsat
Villiesanstrengelse --- indenfra, langsomt og stille, men
sikkert og sundt, ved Opdragelse og fremadskridende
Aandsdannelse.

% --- p. 91 ---
\setcounter{footnote}{0}%
\lecture{Sjette Forelæsning.}\opage{91}

Nutidsexistensen har sit eget Væsen; i Holdning,
Gang og Physionomie adskiller den sig klart og bestemt
fra alle sine Forgængere. Dens Synsmaade er realistisk;
den holder sig til det Virkelige som det staaer og
gaaer. Dog er det ingenlunde dens Mening dermed at
fornegte Aanden eller ringeagte de ideale Fordringer;
men den har en Mistanke til alle abstrakte, høit svævende
Idealer, hvad enten de iklæde sig religiøse, ethiske,
politiske eller æsthetiske Former. Idealer i det Blaa
maae være saa deilige som de ville, de ere dog ufrugtbare
Skyer, luftige Vegetationer, der ikke kunne trives
i Virkelighedens Jordbund; de virke skadeligt, naar de
i Evighedens Navn blodsuge Timeligheden og udmarve
det nærværende Liv, istedenfor at give det Kræfter.
Nutidsexistensen er ingenlunde samvittighedsløs, lider
just heller ikke af Skrupler: dens Samvittighed er som
et Foster, der har vendt sig i Leiet.

At Menneskenaturen paa Grund af sin Dobbelthed
(Dualisme), som sandselig og aandelig, bestandig gjør
Indtryk af, at Mennesket lever som Grændsesoldat mellem
to Riger, et naturligt og et overnaturligt, er en Kjendsgjerning.
Modsætningen mellem det Aandelige og det
% --- p. 92 ---
\setcounter{footnote}{0}%
\opage{92}Legemlige, det Himmelske og det Jordiske, vil, hvor
meget man end for Harmoniens og Eenhedens Skyld
anstrenger sig med at svække og fortrænge det ene
Led, ikke lade sig udslette. Den naturalistiske Stræben:
at forherlige Naturen paa Aandens Bekostning,
og den asketiske: at forhjælpe Aanden til Seier ved
at dræbe Naturen, ere modsvarende Yderligheder. Ingen
af disse Yderligheder er i Stand til at bedaare
Nutiden. Nutidsbevidstheden har Kritik nok til at
indsee, at en naturalistisk Moral ikke blot er lige saa
huul, men i al sin Prunken med Sandselivets Herlighed
endnu mere foragtelig end Asketiken.

Asketen midt i sin Vildfarelse har dog en ubetinget
Anerkjendelse af Aanden som det Høieste. Hans
Opfattelse af Aandens Væsen kan være fattig og uklar;
men har han erkjendt, at Aanden er Villie, da skaaner
han ikke sig selv, men opbyder Alt, for at lyde Aandens
Bud og gjøre dens Fordringer fyldest; han spæger
sit Kjød, han kvæler de naturlige Lyster, han
hærder sin Villie ved Luttring i Ild af Forsagelsens
Smerter. Man kan beklage hans Eensidighed; men man
kan ikke negte ham en vis Høiagtelse.

Naturalisten med „Drifternes Autonomie“ og „Kjødets
Evangelium“ søger sit Støttepunkt i „den moderne
Bevidsthed“. Han kan med en vis Ret beraabe sig
paa, at Naturlivet er Grundlag for Aandslivet, men
han forvansker det Sande ved at nedsætte det foregivne
Aandsliv til en formel Afglands af Bevidstheden
om det naturlige, endelige, forkrænkelige Livs Væsen
og Virken. Han kan gjøre endeel Spræl med de Aandsarbeider,
der allerede ere udførte, og endnu i Fremtiden
skulle udføres, inden det Ideale kommer til sin
% --- p. 93 ---
\setcounter{footnote}{0}%
\opage{93}Ret og Frihedens Værk fuldendes. Men, see vi nærmere
til, saa gaaer det Hele ud paa en Fornegtelse af
Aandens absolut selvgyldige Villie, altsaa af Aanden
som Aand. De splittede Villieskræfter forvandles til
elementære Naturkræfter. Det store Skuespil, som
disse Naturkræfters fornuftnødvendige Brydninger frembringe,
faaer sit aandrige Præg ved at beskues, kritiseres,
og nydes af „den moderne Bevidsthed“. Menneskehedens
høieste Opgave, Dybden i alle dens Længsler,
Lønnen for alle dens Anstrengelser, Maalet for dens
Strid og Stræben maa da sættes i at hæve sig op
paa „den moderne Bevidstheds“ Stade. Hvad Kirken
var for Middelalderens Troende, det er „den moderne
Bevidsthed“ for Nutidens Vidende. Kirken gav Aflad
for allehaande Synder og saa igjennem Fingre endog
med grove Udskeielser, naar blot een Betingelse var
tilstede: Tro paa Kirken og dens Autoritets Ufeilbarhed.
Var denne Betingelse ikke tilstede, da vee Forbryderen;
havde hans Liv end Englenes Reenhed og alle en Helgens gode Gjerninger, han var dog en Kjætter, banlyst
og dømt til evig Fortabelse. „Den moderne Bevidstheds“
kritiske Majestæt er traadt i Kirkens Sted.
Tanken er fri, at sige, naar den ubetinget underkaster
sig „den moderne Bevidstheds“ Ufeilbarhed; thi denne
Ufeilbarhed er ikke blot en Vidensartikel, den er en
Troesartikel. Videnskabens Dyrkere kunne have modstridende
Anskuelser, tvivle om Fornuften, forgude Fornuften
og gjøre sig skyldige endog i grove Inconseqvenser:
siger Intet! Kritiken giver Aflad for allehaande
Modsigelser og seer gjennem Fingre med Tankernes Skrøbelighed,
selv med \emph{grove Inconseqvenser}, dog, vel at
mærke, paa een Betingelse: at Ophavsmændene oprigtig
% --- p. 94 ---
\setcounter{footnote}{0}%
\opage{94}bekjende sig til „den moderne Bevidsthed“. Hvis denne
Betingelse mangler, saa kan den Skyldige pakke ind
og hænge sig selv i Kundskabens Træ. Kjætteren er
dømt --- ikke til Baal og Brand, thi „den moderne
Bevidsthed“ holder paa det Humane --- men til at udstødes
af Videnskaben, berøves Fremtiden og hensættes
i Literaturens Gabestok, Andre til Skræk og Advarsel.

Nutidsbevidstheden, væsenlig seet, falder ingenlunde
sammen med „den moderne Bevidsthed“. Nutidens
Vidende vide meget vel at gjøre Forskjel mellem det
Forbigaaende og det Blivende. Sandheden er hverken
gammeldags eller moderne, den er evig. Den evige
Sandhed kan skifte Form og Udtryk, efter som Tiderne
skifte, men i Væsen og Grund forbliver den altid den
samme.

Den personlige Sandheds Aand, som aabenbarer
sig i Religionen, kan paa Nutidsbevidsthedens Trin erkjendes
i to hinanden modsvarende Virksomhedsformer,
Meddelelsens oprindelige og Tilegnelsens afledede Form,
den første betegnende for Gjennembrudsexistensen, den
sidste afgjørende for Nutidsexistensen. Aldrig før i
Historien er Modsætningen mellem de to Existensformer
traadt saa tydelig frem. Forholdet er hidtil blevet
mere eller mindre fordunklet derved, at man deels har
villet indrette det Afledede i Lighed med det Oprindelige,
deels give det Afledede Skin af at være det Oprindelige
selv i sin videre Fortsættelse, deels endelig
omdanne det Oprindelige ved at tillæmpe det efter det
Afledede. Katholicismen afgiver talende Exempler paa
Forsøg i de to førstnævnte Retninger. Apostolisk Hellighed,
apostoliske Mirakler gjentage sig i skuffende
Efterligninger, selv Frelserens Liv og hellige Vunder
% --- p. 95 ---
\setcounter{footnote}{0}%
\opage{95}copieres i middelalderlige Farvetryk. Den gamle Protestantisme
lod det Oprindelige blive staaende som det
Fuldkomne, det Uopnaaelige, og nøiedes med at tilegne
sig Læren, uden at kunne komme til Klarhed over Existensens
Omdannelse. Medens man i Virkeligheden
gik ind paa en Tillæmpning efter de forandrede Tidsvilkaar
og ydede Statsmyndigheden og det borgerlige
Samfunds mangfoldige Krav fuld Anerkjendelse, følte
man sig dog i sin Samvittighed forpligtet til at indskjærpe
de apostoliske Forskrifters bogstavelige Efterlevelse,
uden at turde vedgaae for sig selv, at det
Sidste stod i skjærende Modsigelse til det Første.
Først i Oplysningsperioden begyndte man at tage Mod
til sig. Ved Fornuftens Lys indsaae man, hvor urimeligt
det var at leve ligefrem efter Apostlenes Anviisning.
Men hvad gjorde man saa? Man gjorde Apostlene
til „fornuftige“ Mænd, underlagde deres Ord en „fornuftig“
Mening og uddrog en jævn borgerlig Moral af
deres Lære. Kristendommens foregivne „Perfectibilitet“
havde til Følge, at det Afledede blev sat i det
Oprindeliges Sted: Copien blev Original, og Originalen
Copi. Det rationalistiske Misgreb var imidlertid for
aabenlyst til, at det i Længden kunde holde sig. Apostelexistensen
er igjen traadt frem i al sin Eiendommelighed,
og Spørgsmaalet er, hvorledes Nutidsexistensen
har at indrette sig, for at det Afledede kan i Aand og
Sandhed svare til det Oprindelige.

Det første Nødvendige er, at Modsætningen træder
ud i fuld Belysning. Aabenbaringsexistensen med de
ideale Karakteermærker bliver staaende uantastet.
Overfor denne Autoritet stiller Nutidsexistensen sig paa
sit historisk bestemte Grundlag med klar Bevidsthed
% --- p. 96 ---
\setcounter{footnote}{0}%
\opage{96}om sine eiendommelige Aandsvilkaar. Nutidsexistensen
veed med sig selv, at den hverken kan eller vil indrette
sig paa at være en ligefrem Efterligning af hiin
ideale Gjennembrudsexistens. Idealexistensen anpriser
for Exempel en eenlig Stand; men Nutidsexistensen
holder paa, at Ægteskab og Familieliv kan i Aandens
Navn være lige saa berettiget som Coelibat.
Naar Idealexistensen foretrækker Armod for Rigdom,
saa er Nutidsexistensen af den Mening, at man uden
at krænke det Ideale kan foretrække Rigdom for Armod.
Naar den ideale Gjennembrudsexistens seer ned
paa politiske Samfundsforhold som paa noget i sig selv
Ligegyldigt og anseer det at være Træl for et lige saa
gunstigt Vilkaar som det at være Fri, saa erklærer
Nutidsexistensen reent ud, at den politiske Tilstand
ingenlunde er ligegyldig, at der i Aandens Kraft maa
kæmpes for borgerlig Orden og borgerlig Frihed, at
Anseelse og Ære er bedre end Ringhed og Vanære.

Hvad Nutiden saaledes udtaler, er dens oprigtige
Hjertens Mening. Ved at give sig som den er
og sige sin Mening reent ud stiller Nutidsexistensen
sig i et sandere Forhold til den hellige Autoritet end
ved at smigre og hykle. Intet kan være Sandhedens
Aand modbydeligere end Hykleri; den guddommelige
Majestæt er ikke tilgængelig for andægtige Complimenter.


Vel kunde det synes, som om Nutidsexistensen ved
at udtale sig i en Retning, der er Idealexistensen saa
modsat, vilde trodse Autoriteten og unddrage sig Aandsfordringerne.
Det er en stor Misforstaaelse. Meningen
er, at de to Existensformer, Idealexistensen med sit
Oprindelige og Realexistensen med sit Afledede, maae
% --- p. 97 ---
\setcounter{footnote}{0}%
\opage{97}træde ud i fuld Modsætning, for at Eenheden kan erkjendes
i sin Dybde, og Aandens Baand imellem dem
knyttes desto fastere. Hvad der skiller Nutidsexistensen
fra det første Oprindelige er noget Mere end „Verdens
Vantro“ og „disse Tiders Ugudelighed“, det er Noget,
hvorover den menneskelige Villie ikke vilkaarlig raader.
Hiint Aandens Gjennembrud skete jo ikke paa Grund
af Menneskenes Hellighed, men Aanden brød igjennem
for at helliggjøre Menneskene. Apostlene vare efter
egen Bekjendelse Syndere, og Apostelmenighedens Medlemmer
vare store Syndere.

Det er Aandsmeddelelsen og de tilsvarende historiske
Vilkaar, hvorpaa Modsætningen væsentlig beroer.
I Gjennembrudsexistensen var den guddommelige Aands
Selvmeddelelse saa mægtigt overstrømmende, at den
drev sine Udvalgte frem til frivilligt Vidnesbyrd, Strid
og Kamp for Lysets og Sandhedens Rige. Over for
dette Lysets og Sandhedens Rige stod Verdensriget,
Mørkets og Løgnens Rige, der maatte bekæmpes med
aandelige Vaaben. Det Aandsfjendtlige var udvortes;
Lidelsen og Trængslen, ja Spliden i Mennesket selv var
særlig betinget af denne Udvorteshed: Aanden var
villig, om end Kjødet var skrøbeligt. I Nutidsexistensen
er Aandsmeddelelsen svagere, ikke et i synlige Udslag
vitterligt Mirakel, men et stille perennerende Under.
De bratte, katastrophiske Overgange fra Mørke til Lys,
fra Vantro til Tro, som karakterisere Gjennembrudsperioden,
ere forvandlede til gradvise, trinvise Omdannelser,
betingede af Villiens Kraft i Tilegnelsens Arbeide.
I Aposteltiden var det Aanden, som paa overnaturlig
Maade greb de menneskelige Sjæle, fyldte dem, løftede
dem, hævede dem op over Tid og Forkrænkelighed. I
% --- p. 98 ---
\setcounter{footnote}{0}%
\opage{98}Nutiden er det omvendt Menneskene, der maae søge
den i sine Frelsesmidler nærværende Aand, gribe den,
holde den fast og ved inderlig Tilegnelse arbeide sig
op igjennem Timelighedens Forviklinger til Sandhed og
evigt Liv. Den hastige Bevægelse ovenfra nedad er
omsat i en langsom Bevægelse nedenfra opad.

Aandens Rige, som i Aabenbaringens Sprog kaldes Guds
Rige, kaldes i Nutidens Sprog Humanitetens Rige. Dersom
det nu kunde vises, at Guds Rige var for indskrænket til at
omfatte Alt, hvad der hører til det sande Menneskelige
i fuld Udstrækning, eller dersom det lod sig godtgjøre,
at Humanitetens Rige havde et eget Grundlag for sig,
hvorpaa det kunde bygges, befæstes og udvides uden
Gudsrigets Bistand, da vilde det Humane og det Religiøse
ganske vist falde udenfor hinanden. Dersom det
Humane virkelig faldt udenfor det Religiøse, da vilde
det jo være ganske i sin Orden, at de to Aandsriger
fremdeles bekæmpede hinanden paa Liv og Død, thi
det enes fuldstændige Seier maatte nødvendig blive det
andets Undergang. Men det viser sig tvertimod, at begge
Riger, oprindelig seet, have samme Grundlag, samme
høieste Formaal, saa at Ingen kan nyde fuld Borgerret
i det ene uden at være Borger i det andet; her
er ikke to Riger, men to Existensformer af eet og samme
guddommeligt-menneskelige Aandens Rige.

Naar det ikke desto mindre seer ud, som var
her virkelig to med hinanden kæmpende Riger, da kan
dette ikke tilskrives den constituerende Grundlov;
det maa være Borgernes Skyld, som ved eensidig at
klamre sig enten til den ene eller til den anden af de
to Existensformer geraade i Splid med hverandre, danne
% --- p. 99 ---
\setcounter{footnote}{0}%
\opage{99}Partier og plage Riget med indvortes Uroligheder. Lad
os forsøge at oplyse dette ved nogle Exempler.

\emph{Aandslovene}. Ved at stille sig exclusivt paa
Gudsrigets Side og tale Aabenbaringens Sprog faaer
Bodsprædikanten Myndighed til at hævde de hellige
Loves ubetingede Gyldighed. Det er jo ikke Menneskebud,
men Guds Bud. Gud selv har udtrykkelig givet
disse Love; Beskrivelsen af den guddommelige Lovgivning
virker imponerende. Jehova tordner fra Sinai,
og Bjerget skjælver, og Tilhørernes Samvittighed skjælver:
det er den Almægtiges Bud, vee Den, der vover
at overtræde dem! --- Pietisternes og Methodisternes Historie
afgiver talende Exempler paa, hvilke Aandens
Rørelser Bodsprædikenen endnu i den nyere Tid formaaer
at fremkalde. Der behøves en hellig Autoritet
for at vække det naturlige Menneske til alvorlig Selvprøvelse.
Der er i Strengheden selv noget Vækkende:
Almuetroen forynges og styrkes derved; de bydende
Ord virke paa Sindet opløftende som Røster fra oven,
forfriskende som Aandepust af guddommelig Majestæt.
Men mærkeligt nok, gjør en saadan Forkyndelse saa
godt som intet Indtryk paa Kulturmennesket. Ved at
stille sig exclusivt paa Humanitetens Standpunkt betragter
han det hele Maskineri af overnaturlige Drivhjul
som et Bevægelsesapparat, den menneskelige Phantasie,
uden ret at vide deraf, selv har construeret. Den
Virkning, Maskineriet gjør paa de troende Sjæle, er
naturligviis en Gjenstand for hans iagttagende Opmærksomhed;
men det er kun en Gjenstand. Det lyner fra
Skyernes Mulm, fra de mørke Vinger, paa hvilke Aanden
svæver, men Lynene ramme ikke den humane Betrag%
% --- p. 100 ---
\setcounter{footnote}{0}%
\opage{100}ter; han staaer rolig under Sinais Torden, uanfægtet
som en Correspondent, der refererer til Bladet.

Og hvad kommer der saa til at staae i Bladet?
„Det er et Tidens Tegn, et mærkeligt Phænomen at
see den veltalende, af Begeistring glødende N. N. samle
Masserne skareviis om sin Prædikestol. En storartet
Opvækkelse! Hans Tale er et Skybrud; Ordene
falde som Regndraaber, vederkvægende efter langvarig
Tørke.“ Saa følger Beskrivelsen af en mærkværdig Mand
med høi Pande, funklende Øine, indfaldne Kinder,
smalle Læber og et langt Skjæg. Miner, Gestus, Talemaader
og Slagord pointeres paa det Nøiagtigste, men
om det Aandelige i Sagen selv meldes der ikke et Ord.

At Gud selv skulde have dicteret Moselovens ti
Bud, vil man paa Humanitetens Vegne finde endnu
urimeligere end at han skulde have dicteret, hvor mange
og hvor lange Gardiner, der skulde ophænges i Tabernaklet,
af hvilke Farver og med hvormange Sløifer.
Humanisterne, som saa klart indsee, at der er Forskjel
mellem sædelige Bud og carmoisinrøde Gardiner, indsee
ogsaa, at Forskjellen netop bestaaer i, at de sædelige
Bud maae have, hvad Gardiner ikke vel kunne have, et
umiskjendeligt Præg af det Fornuftnødvendige, det Almeengyldige.
Det kunde vel ved en foreløbig Betragtning
synes, som var det mindre fornødent, at Gud selv
gav Aabenbaringer om Gardinerne, end at han aabenbarede
Morallovene; men fra et humanistisk Synspunkt
viser Sagen sig omvendt. Om Forhæng, Gardiner, Lampeolie
o. dsl. maatte Jehova nødvendig give ordret Besked,
naar han vilde have Tabernaklet nøiagtig indrettet efter
sit eget Hoved, thi slige Specialiteter, der umulig kunne
udledes af den rene Fornuft, maae just af den Grund
% --- p. 101 ---
\setcounter{footnote}{0}%
\opage{101}henregnes til „Guds Riges“ Inventarium. Hvad Morallovene
derimod angaaer, er en mirakuløs Aabenbaring
aldeles overflødig, thi disse Love have deres tilstrækkelige
Grund i Fornuften og den menneskelige Natur, og
høre udtrykkelig til Humanitetens Rige.

Saa skarpt kan det exclusivt Religiøse og det exclusivt
Humane stille sig mod hinanden. Men at de Religiøse
ved eensidig Fastholden af det Overnaturlige lige saa vist
misforstaae Religionen, som de Humane ved eensidig
at holde sig til det Naturlige misforstaae Humaniteten,
lader sig tydelig paavise.

Religiøse Opvækkelser have til enhver Tid den store
Betydning, at de i Sandhed Opvakte ville gribe Sagen
fra Grunden af, sætte det Oprindelige i Kraft og trænge
ind til Aandslivets Kjerne. Men Grændsen og Skranken
er, at de med al deres Ivren for Aandens Ret og de
hellige Love aldrig naae frem til nogen egenlig Redegjørelse
for det virkelige Livs Naturlove, for Menneskenes
Arbeiden, Virken og Stræben i det Nærværende.
De blive bestandig staaende ved den Forklaring, at det
er en ond, ugudelig, syndig, forfængelig og forkrænkelig
Verden, hvori vi leve. De kunne paa deres Fingre opregne
de mange Overtrædelser, Synder og Udsvævelser,
hvortil Verdens Børn ere hengivne. Men hvad sætte
de i Stedet? Et Hellighedsvæsen som bestaaer i at
gjøre sig Livet saa suurt som muligt. Det er ikke
nok med Overholdelsen af de ti Bud; der er mange
andre Bud at overholde; Guds Børn maae f. Ex. ikke lee,
ikke ryge Tobak, ikke dandse, ikke spille Kort, ikke gaae
paa Komedie o. s. v. Ved at banlyse Livsnydelser og
sætte snævre Grændser for det Tilladelige lægger man
an paa at sikkre Aanden et varigt Herredømme over
% --- p. 102 ---
\setcounter{footnote}{0}%
\opage{102}Kjødet, og vist er det, at dersom der blot var Spørgsmaal
om et Liv med færre og færre Fornøjelser, saa
vilde Brødre og Søstre af den strenge Disciplin staae
paa et overordnet Trin. Men den falske Forudsætning,
at Aand og Natur staae fjendtlig mod hinanden, skinner
igjennem. Saa længe det fra religiøs Side ikke er
indseet, at det Aandelige trænger til det Naturlige, at
det er Aandens Maal at luttre det Naturlige, udvikle,
forme og fremhjælpe det menneskelige Liv i alle Retninger,
saa længe vil Forkyndelsen af Guds Rige være
afmægtig over for Humaniteten, uden væsenlig Indflydelse
paa Nutidsexistensen i det Hele.

Paa Humanitetens Trin er Misligheden af modsat
Art. Her begyndes uden videre med det Naturlige, det
Nærværende, det virkelige Livs mangfoldige Opgaver.
Saa gjælder det om at paavise Lovene. Men hvad er
det for Love? Lutter Naturlove --- gjældende saa vel
for Dyret som for Mennesket. I Menneskelivet skulle
nu disse Love, saa vidt Hjernevirksomheden tillader
det, efterhaanden forvandles til almindelige Fornuftlove
og derved til Aandslove. Men med Forvandlingen gaaer
det kun smaat. Den rene Fornuft med sit kategoriske
Imperativ er en Abstraction, og Naturen gjør Oprør
mod alle Abstractioner. Den virkelige Fornuft ligger
skjult i de virkelige Drifter. De virkelige Drifter ere
ikke blot selviske, men ogsaa sympathetiske. I Bavianernes
indbyrdes Hjælpsomhed er Humanitetens sædelige
Væsen præfigureret. Hos Mennesket bliver ---
især ved den arvelige Hjerneforbedring --- de empiriske
Ideer lidt efter lidt aprioriske. Under fortsat Vexelvirken
af Samfundets mangehaande Kræfter gaae de
sædelige Drifter umærkelig over til sædelige Grund%
% --- p. 103 ---
\setcounter{footnote}{0}%
\opage{103}sætninger, til et System af sædelige Love. De sædelige
Love danne Grundlaget for Humaniteten og den humane
Aand. Næringsvirksomheden, Industrien, Oekonomien,
Kapitalen, Politiken, Videnskaben, Poesien og Kunsten
arbeide i Aandens Tjeneste; Resultatet af alle disse
Arbeider er Menneske-Aandens voxende Magt over den
lavere Natur. Denne Humanitetens Aand har en uendelig
lang Fortid og en uendelig lang Fremtid, men, desværre,
ingen virkelig Nutid; thi den har aldrig havt og vil
aldrig faae nogen virkelig Villieseenhed. Af Vexelvirkningen
mellem Samfundets Enkeltvillier kan Villieseenheden
og dermed den virkelige Aand aldrig resultere.
Enkeltvillierne ere i deres inderste Grund Driftvillier;
men naar Driftvillien skal være det Oprindelige,
saa komme vi --- til Bavianerne.

De exclusivt Religiøse forsee sig paa det Absolute;
de blive staaende ved en overnaturlig Gudsvillie, hvis
ubetingede Bud de i al deres Lydighed ikke formaae
at opfatte som Forkyndelser af en det hele Menneskeliv
gjennemtrængende Aandsvillie. De kunne ikke fyldestgjøre
Religionen, fordi de ikke kunne fyldestgjøre
Humaniteten. De exclusivt Humane forsee sig paa det
Relative. Med deres naturlige Driftvillier ere de ikke
i Stand til at hæve Menneske-Aanden til virkeligt Herredømme
over Naturlivets stridige Elementer, og selvfølgelig
ikke i Stand til at fyldestgjøre Humanitetens
Fordringer. Religionen fordrer det Humane, thi den
fordrer en guddommelig Frigjørelse af Menneskelivets
Kræfter. Humaniteten kræver det Religiøse, thi de
relative Villiers Anstrængelser ere forgjæves, naar de
ikke have den hellige absolute Villie til Bærer og første
Bevæger.

% --- p. 104 ---
\setcounter{footnote}{0}%
\opage{104}\emph{Retfærdighed}. „Skjælver, I Syndere, thi Gud
er retfærdig!“ Naar dette Tilraab i Aabenbaringens
Navn lyder ud over Verden, pleier det sjældent at forfeile
sin Virkning. Just som Menneskene ere allermeest
uenige, ere de netop allermeest enige om at klage over
Urettens Vederstyggelighed, den voxende Uretfærdigheds
Anmasselse, Vold og Vilkaarlighed. Paa det menige
Folks religiøse Følelse vil da hiint formanende Tilraab
virke som et Lynslag og gjøre et mægtigt Ryk i Samvittighederne.
De Troende slaae sig for Brystet, Enhver
bekjender sin Synd for at undflye den kommende
Vrede. Bag denne øieblikkelige Bodfærdighed lurer
ikke sjældent en selvisk Forestilling om, at man ved
Syndernes Bekjendelse strax kan blive Synderne kvit,
at det selvfølgelig ikke saa meget er En selv som det
er den ugudelige, af Uretfærdighed opfyldte Verden,
der maa grue for Straffen. Verdensdommeren, den
Retfærdige, antager sig de Undertrykte, de Forurettede,
og Enhver føler sig jo selv som en Forurettet. „Naar
Menneskesønnen kommer i Skyen med Kraft og megen
Herlighed, skal han skille de Retfærdige fra de Uretfærdige,
som en Hyrde skiller Faarene fra Bukkene,
og da skulle alle Jordens Slægter hyle; men de Udvalgte,
som ere sonede ved Lammets Blod, som have
afført sig deres egen Retfærdighed og iført sig Forsonerens,
skulle tages til Naade.“

Paa Nutidens Kulturmenneske vil en saadan høitidelig
Anmeldelse af Guds Retfærdighed ikke gjøre
noget Indtryk. Kritiken maa finde, at den skarpe Adskillelse
mellem de Retfærdige og de Uretfærdige er
inhuman: der er Ret og Uret paa begge Sider, paa
„Faarenes“ saa vel som paa „Bukkenes.“ Det er heller
% --- p. 105 ---
\setcounter{footnote}{0}%
\opage{105}ikke korrekt at sige: Gud er retfærdig, thi Retfærdighed
er et Begreb, der kun finder Anvendelse paa
menneskelige Forhold. Selv om man vil sige, at Retfærdigheden
er guddommelig, saa er det dog kun en
Abstraction. Den rene Retfærdighed gjælder Intet i
sig selv, det er en Skuemønt, der ikke er i Kurs.
Virkelig Retfærdighed er kun til i det virkelige Sindelag,
der lægger sig for Dagen i Striden om Mit og Dit,
altsaa i den Maade, hvorpaa man opfatter og behandler
de positive Retsspørgsmaal. Men de positive Retsspørgsmaal
ere menneskelige Spørgsmaal, og Benaadningsretten
en menneskelig Ret med en Bismag af Vilkaarlighed,
der karakteriserer Endeligheden og i Et og Alt tilhører
Relativiteten.

Ogsaa fra dette Synspunkt viser Modsætningen
mellem det Religiøse og det Humane sig som en uforsonlig
Splid. Men hvad er Grunden? I den umiddelbart
religiøse Forestilling seer det ud, som var den
menneskelige Retfærdighed uden Gyldighed for Gud.
Det kan da siges om alle Tilsyneladelser af denne
Retfærdighed, hvad en Kirkefader har sagt om Hedningenes
Dyder, at de ere glimrende Laster. Forstaaes
dette saaledes, at menneskelig Retfærdighed,
naar den, fastholdt for sig, vil have Fortjeneste
over for Gud, maa blive et Udtryk for menneskelig
Frækhed, kan herimod Intet indvendes. Naar Retfærdigheden
derimod opfattes saaledes, at dette Væsen
nødvendig maa udelukke den menneskelige Natur
og sætte den guddommelige i Stedet, saa er det Religiøse
rigtignok absolut uforeneligt med det Humane,
men saa er ogsaa den religiøse Grundtanke fordunklet.
De Troende skulle tilegne sig Forsonerens Retfærdighed,
men Forsoneren gjør det jo netop aabenbart, at det
% --- p. 106 ---
\setcounter{footnote}{0}%
\opage{106}Guddommelige selv er menneskeligt. Dersom Forsoneren
blot opfattes som en historisk Individualitet, der
engang har lidt og er død paa Grund af Menneskenes
Uretfærdighed, da er det aldeles umuligt at tilegne sig
hans Retfærdighed. Men Forsonerens Væsen er i gudmenneskelig
Forstand Menneskehedens inderste sande
Væsen. Forsoneren kom netop een Gang for alle udvortes
til Menneskene som historisk Skikkelse, for at han nu
og fremdeles kan komme i Aanden indvortes med sin
Retfærdighed, ikke som en fremmed Person, men som
Væsen og Grund i al sand menneskelig Personlighed. At
tilegne sig Forsonerens Retfærdighed er at guddommeliggjøre
den menneskelige og menneskeliggjøre den guddommelige
Retfærdighed, er i dybeste Forstand at fyldestgjøre
Humaniteten.

Stille vi os paa det humane Standpunkt og begynde
med Striden om Mit og Dit, da bliver det jo indrømmet,
at Retsstridigheder, selv de ivrigste, umulig kunne føre
til retfærdige Forhold, med mindre de begrundes paa Retfærdighedstanken,
føres med et retfærdigt Sindelag
og en retfærdig Villie. Agter man ikke Retfærdighedens
ubetingede Gyldighed, men gjør Retfærdigheden
i og for sig til et tomt Begreb, da vil den endelige
Retsstrid ende med den Stærkeres Ret, og den
relative Retfærdighed forvandle sig til et positivt Udtryk
for den relative Uretfærdighed. Skal den relative
Retfærdighed tilnærmelsesviis seire over Uretfærdighed,
da maa det skee i Kraft af Retfærdighedens Idee. Men
Retfærdighedens Idee er en Villiesidee, den absolute
Retfærdighed forudsætter en absolut retfærdig Villie,
en Villie, som i sidste Instans faaer Bugt med al Uretfærdighed.
Tro paa den seirende Retfærdighed er
uadskillelig fra Tro paa Retfærdighedens Gud. Gud
% --- p. 107 ---
\setcounter{footnote}{0}%
\opage{107}maa for den menneskelige Retfærdigheds Skyld absolut
være retfærdig: det Humane kræver det Religiøse.

Hvad hermed er antydet, kan udvikles i det Uendelige;
men jo videre Udviklingen gaaer, og jo nøiagtigere
Forholdsbestemmelserne angives, desto tydeligere viser
det sig, at Religionen maa være human, og Humaniteten
religiøs.

\emph{Aandsdannelse}. Stille vi Nutidsexistensen paa
sit historiske Grundlag med sine Arbeider og Formaal
frit over for Apostelexistensen med de ideale Gjennembrud,
idet dog begge Existenser erkjendes for væsenlig
sammenhørende, saa faae vi et Indtryk af storstilet
Dualisme. Gudsrigets Tjenere med den himmelstræbende
Livsretning og altopoffrende Verdensforsagelse,
„under Ære og Vanære, under ondt Rygte og godt
Rygte, som Forførere og dog sanddrue, som miskjendte
og dog erkjendte, som de, der døe og dog leve, som
de, der ere revsede og dog ikke ihjelslagne, som bedrøvede
og dog glade, som fattige, der dog gjøre Mange
rige, som de, der have Intet og dog besidde Alt“, ---
og saa Humanitetens Tjenere, fordybede i det Jordiske,
stridende for det Timelige, arbeidende for det Nærværende,
for Øieblikkets Krav i Endelighedens Trældom,
rige og dog fattige, overmætte og dog hungrige, freidige
og dog i hemmelig Angst, glade og dog bekymrede,
altid bekymrede for verdslige Goder, Ære, Magt og
borgerlig Indflydelse, --- hvilken uhyre Modsætning!

Og dog høre begge Existenser ind under eet og
samme Rige, Aandens Rige; Grundloven er for begge
den samme, Frihedens, Retfærdighedens og Kjærlighedens
Grundlov. Den ubetingede Fordring er, at
Aanden skal herske overalt i sit Rige, i de lavere som
% --- p. 108 ---
\setcounter{footnote}{0}%
\opage{108}i de høiere Regioner. Hvis nu Apostelexistensen ikke
stod fast paa Autoritetens hellige Grund, saa vilde
Vidnesbyrdet om Aandens Ret og Maal forsvinde. Hvad
Aanden kan fordre af Mennesket, naar sandselig Drift
og selvisk Egenraadighed gjør Oprør mod dens Herredømme,
hvilke Offre den kræver, hvilke Kræfter den
har, hvilke Midler den tilbyder, naar det i Striden
kommer til det Yderste, vilde reent være glemt; det
Hellige, det absolut Ophøiede, der holder Maalestokken,
vilde synke ned i lutter Endelighed, lutter Relativitet,
og falsk Vægt og Maal indsnige sig overalt og forvirre
de humane Forhold. Saa lidt som der kan gjøres
Ret og Skjel i et borgerligt Samfund med falsk Mynt,
falsk Maal og Vægt, lige saa lidt kan der holdes Orden
i Aandens Rige, naar de ideale Maalestokke ere brudte,
naar det Høiere kan blive det Lavere, og det Lavere
det Høiere, alt som Døgnet skifteviis løfter og styrter
sine Afguder.

Idealexistensen gjælder ikke alene for Guds Skyld,
men først og sidst for Menneskets Skyld, og Nutiden
har for saa vidt Ret, naar den fæster sit Blik paa
Realiteternes Verden og holder paa det Humane.
Det er unegtelig en Indbildning, ja en indgroet
Fordom, at Aandens hellige Gud skulde have foreskrevet
Menneskene en eneste Lov, en eneste Betingelse, som
ikke var begrundet i det menneskelige Væsens Anlæg,
ikke var nødvendig for dets frie, personlige Selvudvikling.
Men at det Personlige kan begrundes paa et
Upersonligt, er ogsaa en Indbildning.

Det er Menighedslivet, der vedblivende danner
Grundlag for Aandslivet. Ordets Tjenere i Menigheden
have den store Opgave at haandhæve og indskærpe
Gudsrigets Vilkaar, ikke med paatagen Apostelmyndig%
% --- p. 109 ---
\setcounter{footnote}{0}%
\opage{109}hed, men, for Tilegnelsens Skyld, med Aandsdannelsens
hele Overlegenhed. Ungdommens Lærere have den store
Opgave ved Aandsdannelsens Midler, at give de aabenbarede
Sandheder et praktisk fatteligt Udtryk i Humanitetens
Sprog og ved Idealernes Autoritet i Opdragelsens
Alvor gjøre dem nærværende. --- Saa er der
vel en Forskjel mellem Kirkeskolen og Kulturskolen,
men Spliden er hævet, en gjensidig Forstaaelse er opnaaet,
saa at begge Skoler, om end ved forskjellige
Midler og ad forskjellige Veie, kunne arbeide i samme
Aand til Opnaaelse af det samme Maal.

Der er vel dem, som mene, at Religionens Tid er
forbi, at Humaniteten alene har Fremtiden for sig. Vi
skulle efter deres Anviisning ikke blive staaende paa
kirkelig Grund og med Luther stirre mod Østen; vi
skulle med en ny Columbus heise „den frie Tankes“
svaiende Flag og freidig styre vor Kurs mod Vesten. „Fjernt
i Vest, i Naturopdagelsernes Atlanterhav ligger Aandens
forjættede Land, skjult som en Perle i Dybet.“ --- Det er
en falsk Anviisning. Den, der blindt hen styrter
sig i Naturalismens Svælg, dukker ikke op „med en Perle
i Haanden“, men --- med en Blære i Munden.

Den samme Livets og Sandhedens Aand, der fordum
ved overordenlig Selvmeddelelse frembragte Apostelexistensen,
meddeler sig endnu paa Tilegnelsens humane
Vilkaar med Foryngelsens og Omdannelsens Kraft til
Nutidsexistensen. Schemaet er et andet, men Indholdet
er det samme. Midlerne have vi, Betingelserne ere os
givne; det kommer kun an paa, at vi bruge dem rettelig.
Til at bruge dem rettelig udkræves der Opdragelse og
Aandsdannelse.

\end{document}
